% Lettres de réclamation et curriculum vitae : forme et contenu — Allemand Tle A — Cours signé OnBuch+
% Programme officiel MINESEC. Compiler avec : tectonic AL29-lettres-de-reclamation-et-curriculum-vitae-forme-et-contenu.tex
\newcommand{\DOCMATIERE}{Allemand}
\newcommand{\DOCNIVEAU}{Tle A}
\newcommand{\DOCMODULE}{Chapitre 5 : Citoyenneté}
\newcommand{\DOCLECON}{AL29}
\newcommand{\DOCTITRE}{Lettres de réclamation et curriculum vitae : forme et contenu}
% OnBuch+ — préambule « cours signés » ENRICHI (compilé avec tectonic / XeLaTeX)
% Système visuel « Pop » d'OnBuch+ : fond crème (jamais blanc), bordures encre
% épaisses, ombres « dures » décalées, titres Archivo Black en capitales.
% Version MATHÉMATIQUES : graphes (pgfplots/tikz), boîtes pédagogiques
% (définition, propriété, méthode, attention, le savais-tu, exercice résolu,
% exercice, TP, bilan), page de garde et signature OnBuch+ ancrée en bas.
%
% Macros attendues dans main.tex AVANT \input{preamble} :
%   \DOCMATIERE  \DOCNIVEAU  \DOCMODULE  \DOCLECON (numéro)  \DOCTITRE
\documentclass[11pt]{article}

\usepackage{fontspec}
\usepackage[ngerman,french]{babel} % français = langue principale ; l'allemand via \all{..} / textebox
\usepackage[a4paper,margin=2cm,top=3.4cm,bottom=2.8cm,headheight=1.4cm,footskip=1.3cm]{geometry}
\usepackage{amsmath,amssymb,amsfonts}
\usepackage{mathtools} % \xmapsto, \coloneqq... souvent utilisés par les modèles
\usepackage{cancel} % simplifications barrées (bilans de forces)
% \abs{z} (module, valeur absolue) et \norm{u} (norme), utilisés par les modèles
\providecommand{\abs}{}\let\abs\relax\DeclarePairedDelimiter{\abs}{\lvert}{\rvert}
\providecommand{\norm}{}\let\norm\relax\DeclarePairedDelimiter{\norm}{\lVert}{\rVert}
\usepackage[table]{xcolor} % [table] : fournit aussi \rowcolors (alternance de lignes)
\usepackage{pagecolor}
\usepackage{tikz}
\usetikzlibrary{arrows.meta,positioning,calc,shapes.geometric,shapes.misc,shapes.arrows,decorations.pathmorphing,decorations.pathreplacing,decorations.markings,patterns,shadows,mindmap,trees,fit,backgrounds,matrix,intersections,angles,quotes}
\usepackage{pgfplots}
\usepackage[version=4]{mhchem} % équations nucléaires \ce{^{235}_{92}U}
\usepackage[european,siunitx]{circuitikz} % schémas électriques
\pgfplotsset{compat=1.18}
\usepgfplotslibrary{fillbetween,groupplots} % aires entre courbes (calcul intégral)
\usepackage{enumitem}
\usepackage{fancyhdr}
% [most] charge toutes les bibliothèques tcolorbox (listings, minted, raster,
% documentation...) et gonfle inutilement la mémoire TeX sur un document long
% (~300 boîtes) : on ne charge que ce qui est réellement utilisé.
\usepackage[skins,breakable]{tcolorbox}
\usepackage{graphicx}
\usepackage{colortbl}
\usepackage{booktabs}
\usepackage{tabularx}
\usepackage{diagbox} % case d'en-tête diagonale des tableaux à double entrée
% Colonnes extensibles centrée / gauche / droite (C, L, R), souvent utilisées par les modèles
\newcolumntype{C}{>{\centering\arraybackslash}X}
\newcolumntype{L}{>{\raggedright\arraybackslash}X}
\newcolumntype{R}{>{\raggedleft\arraybackslash}X}
\usepackage{array}
\usepackage{multicol}
\usepackage{multirow}
\usepackage{siunitx}
% parse-numbers=false : accepte aussi \SI{2,5 \times 10^{-3}}{...}
\sisetup{locale=FR,output-decimal-marker={,},group-minimum-digits=4,parse-numbers=false}
\DeclareSIUnit\ohm{\ensuremath{\symup{\Omega}}} % U+2126 absent de la police
\DeclareSIUnit\division{div} % réglages d'oscilloscope (ms/div, V/div)
\DeclareSIUnit\tr{tr} % tours (vitesse de rotation, ex. \SI{1500}{\tr\per\minute})
\DeclareSIUnit\molar{M} % concentration molaire (ex. \SI{3}{\molar}, \SI{1}{\micro\molar})
\DeclareSIUnit\an{an} % année (le modèle confond parfois avec \year, un compteur TeX) — ex. \SI{5}{\kilo\meter\per\an}
\DeclareSIUnit\FCFA{FCFA} % franc CFA (ex. \SI{500}{\FCFA\per\kilo\gram})
\DeclareSIUnit\degreeBrix{\textdegree Brix} % teneur en sucre d'un sirop/jus (réfractométrie)
\DeclareSIUnit\month{mois} % le modèle utilise parfois \month, un compteur TeX, comme unité de durée
\usepackage{qrcode}
% \degree : commande fréquemment utilisée par les modèles pour un angle ou une
% température en dehors de \SI{}{} (non fournie par les paquets chargés).
\providecommand{\degree}{^{\circ}}
% \qed (fin de démonstration), utilisé par les modèles sans amsthm
\providecommand{\qed}{\hfill$\square$}
% \barre : double barre des valeurs interdites dans un tableau de variations (tkz-tab)
\providecommand{\barre}{\ensuremath{\|}}
% environnement proof (amsthm non chargé) utilisé par les modèles
\ifdefined\proof\else\newenvironment{proof}[1][Démonstration]{\par\noindent\textit{#1.}\ }{\qed\par}\fi
\usepackage{lastpage}
\usepackage{hyperref}
\hypersetup{hidelinks}

% ============================================================
% Palette « Pop » OnBuch+ — mêmes hex que la classe P (plus_ui.dart)
% ============================================================
\definecolor{poporange}{HTML}{F59321}
\definecolor{popdark}{HTML}{E07A0C}
\definecolor{popcream}{HTML}{FFF6E8}
\definecolor{popink}{HTML}{1C1714}
\definecolor{poppurple}{HTML}{7A5AE0}
\definecolor{popgreen}{HTML}{2E9E6A}
\definecolor{poppink}{HTML}{E0457B}
\definecolor{popblue}{HTML}{2F7DE1}
\definecolor{popgold}{HTML}{C98A12}
\definecolor{popmuted}{HTML}{8A7F76}
\definecolor{popline}{HTML}{EADFD2}
\definecolor{popgray}{HTML}{8A7F76}
\colorlet{popgrey}{popgray}
% Alias tolérants (le modèle nomme parfois les couleurs différemment)
\colorlet{popwhite}{white}
\colorlet{popblack}{popink}
\colorlet{popred}{poppink}
\colorlet{popyellow}{popgold}
% Teintes claires (fonds de tableaux/encadrés) — le modèle nomme parfois
% les couleurs avec un suffixe L (ex. popblueL) sans les avoir définies.
\colorlet{poporangeL}{poporange!20}
\colorlet{popdarkL}{popdark!20}
\colorlet{poppurpleL}{poppurple!20}
\colorlet{popgreenL}{popgreen!20}
\colorlet{poppinkL}{poppink!20}
\colorlet{popblueL}{popblue!20}
\colorlet{popgoldL}{popgold!20}
\colorlet{popredL}{popred!20}
\colorlet{popgrayL}{popgray!20}
% Teintes claires pour fonds de tableaux / zones
\colorlet{popblueL}{popblue!10!white}
\colorlet{poporangeL}{poporange!14!white}
\colorlet{popgreenL}{popgreen!12!white}
\colorlet{poppurpleL}{poppurple!12!white}
\colorlet{poppinkL}{poppink!12!white}
\colorlet{popgoldL}{popgold!14!white}

\pagecolor{popcream}

% ============================================================
% Typographie
% ============================================================
\defaultfontfeatures{Ligatures=TeX}
\setmainfont{PlusJakartaSans-Medium.ttf}[Path=../../../fonts/,BoldFont=PlusJakartaSans-Bold.ttf]
% Maths en sans-serif (Fira Math) avec chiffres et lettres droites de Plus
% Jakarta, pour que les formules se fondent dans le texte.
\usepackage[math-style=ISO,bold-style=ISO]{unicode-math}
\setmathfont{FiraMath-Regular.otf}[Path=../../../fonts/]
\setmathfont{PlusJakartaSans-Medium.ttf}[Path=../../../fonts/,range={up/num,up/{Latin,latin},\mathup}]
\usepackage{icomma} % virgule décimale sans espace en mode math
% Caractères mathématiques Unicode absents des polices de texte (Plus Jakarta,
% Archivo Black) : rendus par leur équivalent mathématique, en texte comme en
% formule — sinon ils s'affichent en carré vide (ex. « Arithmétique dans ℤ »).
\usepackage{newunicodechar}
\newunicodechar{ℝ}{\ensuremath{\mathbb{R}}}
\newunicodechar{ℤ}{\ensuremath{\mathbb{Z}}}
\newunicodechar{ℂ}{\ensuremath{\mathbb{C}}}
\newunicodechar{ℕ}{\ensuremath{\mathbb{N}}}
\newunicodechar{ℚ}{\ensuremath{\mathbb{Q}}}
\newunicodechar{𝔻}{\ensuremath{\mathbb{D}}}
\newunicodechar{α}{\ensuremath{\alpha}}
\newunicodechar{β}{\ensuremath{\beta}}
\newunicodechar{γ}{\ensuremath{\gamma}}
\newunicodechar{δ}{\ensuremath{\delta}}
\newunicodechar{ε}{\ensuremath{\varepsilon}}
\newunicodechar{θ}{\ensuremath{\theta}}
\newunicodechar{λ}{\ensuremath{\lambda}}
\newunicodechar{μ}{\ensuremath{\mu}}
\newunicodechar{σ}{\ensuremath{\sigma}}
\newunicodechar{τ}{\ensuremath{\tau}}
\newunicodechar{φ}{\ensuremath{\varphi}}
\newunicodechar{ψ}{\ensuremath{\psi}}
\newunicodechar{ω}{\ensuremath{\omega}}
\newunicodechar{Σ}{\ensuremath{\Sigma}}
% majuscules produites par \MakeUppercase dans les titres (α-aminés -> Α)
\newunicodechar{Α}{\ensuremath{\alpha}}
\newunicodechar{Β}{\ensuremath{\beta}}
\newunicodechar{Γ}{\ensuremath{\Gamma}}
\newunicodechar{Δ}{\ensuremath{\Delta}}
\newunicodechar{Ε}{\ensuremath{\varepsilon}}
\newunicodechar{Θ}{\ensuremath{\Theta}}
\newunicodechar{Λ}{\ensuremath{\Lambda}}
\newunicodechar{Μ}{\ensuremath{\mu}}
\newunicodechar{Τ}{\ensuremath{\tau}}
\newunicodechar{Φ}{\ensuremath{\Phi}}
\newunicodechar{Ψ}{\ensuremath{\Psi}}
\newunicodechar{Ω}{\ensuremath{\Omega}}
\newunicodechar{↦}{\ensuremath{\mapsto}}
\newunicodechar{∈}{\ensuremath{\in}}
\newunicodechar{∉}{\ensuremath{\notin}}
\newunicodechar{∧}{\ensuremath{\wedge}}
\newunicodechar{⇔}{\ensuremath{\Leftrightarrow}}
\newunicodechar{⟺}{\ensuremath{\Longleftrightarrow}}
\newunicodechar{⇒}{\ensuremath{\Rightarrow}}
\newunicodechar{⟹}{\ensuremath{\Longrightarrow}}
\newunicodechar{∘}{\ensuremath{\circ}}
\newunicodechar{∪}{\ensuremath{\cup}}
\newunicodechar{∩}{\ensuremath{\cap}}
\newunicodechar{≡}{\ensuremath{\equiv}}
\newunicodechar{⊂}{\ensuremath{\subset}}
\newunicodechar{⊥}{\ensuremath{\perp}}
\newunicodechar{∃}{\ensuremath{\exists}}
\newunicodechar{∀}{\ensuremath{\forall}}
\newunicodechar{∛}{\ensuremath{\sqrt[3]{\,}}}
\newunicodechar{∜}{\ensuremath{\sqrt[4]{\,}}}
\newunicodechar{⃗}{\ensuremath{{}^{\to}}}
\newunicodechar{ⁿ}{\ifmmode^{n}\else\textsuperscript{\ensuremath{n}}\fi}
\newunicodechar{ˣ}{\ifmmode^{x}\else\textsuperscript{\ensuremath{x}}\fi}
\newunicodechar{ᵇ}{\ifmmode^{b}\else\textsuperscript{\ensuremath{b}}\fi}
\newunicodechar{ʳ}{\ifmmode^{r}\else\textsuperscript{\ensuremath{r}}\fi}
\newunicodechar{ᵘ}{\ifmmode^{u}\else\textsuperscript{\ensuremath{u}}\fi}
\newunicodechar{ʸ}{\ifmmode^{y}\else\textsuperscript{\ensuremath{y}}\fi}
\newunicodechar{ᵃ}{\ifmmode^{a}\else\textsuperscript{\ensuremath{a}}\fi}
\newunicodechar{ᵉ}{\ifmmode^{e}\else\textsuperscript{\ensuremath{e}}\fi}
\newunicodechar{ᵏ}{\ifmmode^{k}\else\textsuperscript{\ensuremath{k}}\fi}
\newunicodechar{ᵖ}{\ifmmode^{p}\else\textsuperscript{\ensuremath{p}}\fi}
\newunicodechar{ᴺ}{\ifmmode^{N}\else\textsuperscript{\ensuremath{N}}\fi}
\newunicodechar{ᶜ}{\ifmmode^{c}\else\textsuperscript{\ensuremath{c}}\fi}
\newunicodechar{ʰ}{\ifmmode^{h}\else\textsuperscript{\ensuremath{h}}\fi}
\newunicodechar{ᵠ}{\ifmmode^{q}\else\textsuperscript{\ensuremath{q}}\fi}
\newunicodechar{ₙ}{\ifmmode_{n}\else\textsubscript{\ensuremath{n}}\fi}
\newunicodechar{ₐ}{\ifmmode_{a}\else\textsubscript{\ensuremath{a}}\fi}
\newunicodechar{ᵢ}{\ifmmode_{i}\else\textsubscript{\ensuremath{i}}\fi}
\newunicodechar{ᵣ}{\ifmmode_{r}\else\textsubscript{\ensuremath{r}}\fi}
\newunicodechar{ₖ}{\ifmmode_{k}\else\textsubscript{\ensuremath{k}}\fi}
\newunicodechar{ᵦ}{\ifmmode_{\beta}\else\textsubscript{\ensuremath{\beta}}\fi}
\newunicodechar{ₓ}{\ifmmode_{x}\else\textsubscript{\ensuremath{x}}\fi}
\newfontfamily\popdisplay{ArchivoBlack-Regular.ttf}[Path=../../../fonts/]
\newfontfamily\poplogo{SpaceGrotesk-Bold.ttf}[Path=../../../fonts/]
\newfontfamily\popsemi{PlusJakartaSans-SemiBold.ttf}[Path=../../../fonts/]
\newfontfamily\popxbold{PlusJakartaSans-ExtraBold.ttf}[Path=../../../fonts/]
\newfontfamily\popxboldls{PlusJakartaSans-ExtraBold.ttf}[Path=../../../fonts/,LetterSpace=3.0]

\setlist[enumerate]{leftmargin=1.7em,itemsep=5pt,topsep=4pt}
\setlist[itemize]{leftmargin=1.7em,itemsep=4pt,topsep=4pt,label={\color{poporange}\textbullet}}
\setlength{\parindent}{0pt}
\setlength{\parskip}{5pt}
\renewcommand{\arraystretch}{1.25}

% pgfplots : style Pop par défaut
\pgfplotsset{
  every axis/.append style={axis line style={popink,line width=1pt},tick style={popink},
    grid style={popline,line width=0.5pt},label style={font=\small},tick label style={font=\footnotesize}},
  popcurve/.style={poporange,line width=1.8pt,no markers,smooth},
}
% Même style disponible dans un \draw TikZ simple (hors pgfplots)
\tikzset{popcurve/.style={poporange,line width=1.8pt}}
\tikzset{popgrid/.style={popline,very thin}}
\colorlet{popcurve}{poporange} % le nom du style est parfois utilisé comme couleur

% ============================================================
% Pill / squiggle
% ============================================================
\newcommand{\poppill}[3]{%
  \tikz[baseline=-0.55ex]\node[draw=popink,line width=1.2pt,rounded corners=2.6pt,%
    fill=#2,inner xsep=6.5pt,inner ysep=3pt]{%
    \popxboldls\fontsize{7}{7}\selectfont\color{#3}\MakeUppercase{#1}};%
}
\newcommand{\popsquiggle}[1]{%
  \tikz[baseline=-0.4ex]{
    \draw[#1,line width=2.6pt,line cap=round]
      plot [smooth] coordinates {(0,0.09) (0.34,0.01) (0.68,0.21) (1.02,0.01) (1.36,0.10)};
  }%
}
% Mot-clé surligné (marqueur orange sous le texte)
\newcommand{\cle}[1]{{\popxbold\color{popdark}#1}}

% ============================================================
% En-tête / pied
% ============================================================
\pagestyle{fancy}
\fancyhf{}
\renewcommand{\headrulewidth}{0pt}
\renewcommand{\footrulewidth}{0pt}
\lhead{\raisebox{-0.15ex}{\poplogo\fontsize{13}{13}\selectfont\color{popink}OnBuch\color{poporange}+}}
\rhead{\poppill{\DOCMATIERE\ \textbullet\ \DOCNIVEAU\ \textbullet\ Leçon \DOCLECON}{white}{popink}}
\lfoot{\popxbold\fontsize{7.6}{7.6}\selectfont\color{popmuted}\MakeUppercase{Cours signé OnBuch+}\ \textbullet\ {\popsemi\DOCTITRE}}
\rfoot{\tikz[baseline=-0.6ex]\node[rounded corners=8.5pt,fill=popink,minimum height=17pt,minimum width=17pt,inner xsep=5pt,inner sep=0]{\popxbold\fontsize{8}{8}\selectfont\color{popcream}\thepage\,/\,\pageref*{LastPage}};}

% ============================================================
% Titres
% ============================================================
\newcommand{\chaptitre}[2]{%
  \begin{tcolorbox}[enhanced,colback=poporange,colframe=popink,boxrule=2.6pt,arc=11pt,
      drop shadow=popink,left=15pt,right=15pt,top=12pt,bottom=11pt,before skip=3pt,after skip=18pt]
    \poppill{#2}{popink}{popcream}\par\vspace{9pt}
    {\raggedright\hyphenpenalty=10000\exhyphenpenalty=10000\popdisplay\fontsize{22}{24}\selectfont\color{popcream}\MakeUppercase{#1}\par}
  \end{tcolorbox}
}
% Section numérotée : pastille orange avec le numéro + titre Archivo
\newcounter{popsec}
\newcommand{\coursec}[1]{%
  \stepcounter{popsec}\setcounter{popsub}{0}%
  \par\vspace{16pt}\noindent
  \begin{minipage}{\linewidth}
  \tikz[baseline=-0.7ex]\node[draw=popink,line width=1.6pt,fill=poporange,rounded corners=4pt,minimum size=22pt,inner sep=0]{\popdisplay\fontsize{12}{12}\selectfont\color{popcream}\thepopsec};%
  \hspace{8pt}{\popdisplay\fontsize{15}{17}\selectfont\color{popink}\MakeUppercase{#1}}\par
  \vspace{4pt}\noindent{\color{poporange}\rule{36pt}{3.6pt}}
  \end{minipage}\par\nopagebreak\vspace{8pt}
}
\newcounter{popsub}
\newcommand{\courssub}[1]{%
  \stepcounter{popsub}%
  \par\vspace{9pt}\noindent{\popxbold\fontsize{11.5}{13}\selectfont\color{popink}{\color{poporange}\thepopsec.\thepopsub}\hspace{6pt}#1}\par\nopagebreak\vspace{4pt}
}

% ============================================================
% Boîtes pédagogiques — même recette : fond blanc, bordure épaisse colorée,
% ombre dure encre, badge plein. Titre optionnel [#1].
% ============================================================
\tcbset{popbase/.style={breakable,enhanced,colback=white,boxrule=2.3pt,arc=9pt,
  shadow={3pt}{-3pt}{0pt}{popink},left=14pt,right=14pt,top=11pt,bottom=11pt,before skip=11pt,after skip=15pt,
  fontupper=\popsemi\fontsize{10.6}{14.5}\selectfont\color{popink},
  fontlower=\popsemi\fontsize{10.6}{14.5}\selectfont\color{popink}}}
% Coche dessinée (le glyphe ✓ n'existe pas dans les polices du document)
\newcommand{\popcheck}[1]{\tikz[baseline=-0.5ex]\draw[#1,line width=1.6pt,line cap=round,line join=round](-0.13,0.0)--(-0.04,-0.09)--(0.13,0.1);}
\providecommand{\coche}{\popcheck{popgreen}}
\providecommand{\resultat}[1]{\par\smallskip\noindent\fcolorbox{popgreen}{popgreenL}{\parbox{\dimexpr\linewidth-2\fboxsep-2\fboxrule\relax}{\textbf{Résultat.} #1}}\par\smallskip}
\newcommand{\popboxhead}[3]{\poppill{#1}{#2}{white}\ifx\relax#3\relax\else\hspace{6pt}{\popxbold\fontsize{10.6}{12}\selectfont\color{popink}#3}\fi\par\vspace{7pt}}

\newtcolorbox{aretenir}[1][]{popbase,colframe=popgreen,before upper={\popboxhead{À retenir}{popgreen}{#1}}}
% Texte en allemand (document de lecture, dialogue, extrait) : cadre
% d'étude. Un titre optionnel (source, auteur) s'affiche dans l'en-tête.
\newtcolorbox{textebox}[1][]{popbase,colframe=popblue,before upper={\popboxhead{Text}{popblue}{#1}\begin{otherlanguage}{ngerman}},after upper={\end{otherlanguage}}}
% Marques ✓ / ✗ (forme correcte / fautive) souvent utilisées par les modèles
\providecommand{\cross}{\textcolor{poppink}{\ensuremath{\times}}}
\providecommand{\xmark}{\cross}\providecommand{\crossmark}{\cross}\providecommand{\cmark}{\ensuremath{\checkmark}}
% Ligne de réponse / justification à compléter (inventée par les modèles)
\providecommand{\justification}[1]{\par\noindent\textit{Justification~:} \dotfill\par}
% \textipa (package tipa non chargé) : la police ne couvre pas l'API -> simple passage
\providecommand{\textipa}[1]{#1}
% Soulignement (ulem non chargé) : \uline, \ul
\providecommand{\uline}[1]{\underline{#1}}\providecommand{\ul}[1]{\underline{#1}}
% \glossaire{\item ...} inventée par les modèles : liste de glossaire
\providecommand{\glossaire}[1]{\begin{itemize}#1\end{itemize}}
% \alert (beamer) inventée par les modèles : texte mis en évidence
\providecommand{\alert}[1]{\textbf{\textcolor{poppink}{#1}}}
% variantes ASCII d'ordinaux écrites en macro
\providecommand{\iere}{\textsuperscript{re}}\providecommand{\ieme}{\textsuperscript{e}}\providecommand{\ere}{\textsuperscript{re}}\providecommand{\eme}{\textsuperscript{e}}\providecommand{\er}{\textsuperscript{er}}\providecommand{\ie}{\textsuperscript{e}}
% Ordinaux français écrits en macro par les modèles : 1\ière, 3\ième
\newcommand{\ière}{\textsuperscript{re}}\newcommand{\ième}{\textsuperscript{e}}
% \exemple (inventée par les modèles) : étiquette « Exemple : »
\providecommand{\exemple}{\textbf{Exemple~:}\ }
% \square utilisable aussi hors mode math (cases à cocher)
\AtBeginDocument{\let\squareMath\square\renewcommand{\square}{\ensuremath{\squareMath}}}
% Mot ou phrase en allemand à l'intérieur d'un paragraphe français
\newcommand{\all}[1]{\ifmmode\text{\textit{#1}}\else\begingroup\selectlanguage{ngerman}\itshape #1\endgroup\fi}
\let\esp\all % alias toléré
\newtcolorbox{exemplebox}[1][]{popbase,colframe=poporange,before upper={\popboxhead{Exemple}{poporange}{#1}}}
\newtcolorbox{definition}[1][]{popbase,colframe=popblue,before upper={\popboxhead{Définition}{popblue}{#1}}}
\newtcolorbox{propriete}[1][]{popbase,colframe=poppurple,before upper={\popboxhead{Propriété}{poppurple}{#1}}}
\newtcolorbox{methode}[1][]{popbase,colframe=popgold,before upper={\popboxhead{Méthode}{popgold}{#1}}}
\newtcolorbox{attention}[1][]{popbase,colframe=poppink,before upper={\popboxhead{Attention}{poppink}{#1}}}
\newtcolorbox{experience}[1][]{popbase,colframe=popblue,colback=popblueL,before upper={\popboxhead{Expérience}{popblue}{#1}}}
% « Le savais-tu ? » : carte violette pleine (contexte, culture, Cameroun)
\newtcolorbox{savaistu}[1][]{popbase,colback=poppurple,colframe=popink,
  fontupper=\popsemi\fontsize{10.4}{14.2}\selectfont\color{white},
  before upper={\poppill{Le savais-tu\,?}{white}{poppurple}\ifx\relax#1\relax\else\hspace{6pt}{\popxbold\color{white}#1}\fi\par\vspace{7pt}}}
% Exercice résolu : énoncé en haut, \tcblower puis la solution sur fond crème
\newcounter{popexo}\newcounter{popexr}
\newtcolorbox{exoresolu}[1][]{popbase,colframe=poporange,colbacklower=popcream,
  segmentation style={popink,line width=1.2pt,dashed},
  before upper={\stepcounter{popexr}\popboxhead{Exercice résolu \thepopexr}{poporange}{#1}},
  before lower={\poppill{Solution}{popink}{popcream}\par\vspace{6pt}}}
% Exercice (non résolu en place) — la correction est dans la partie Corrigés
\newtcolorbox{exercice}[1][]{popbase,colframe=popink,
  before upper={\stepcounter{popexo}\popboxhead{Exercice \thepopexo}{popink}{#1}}}
% Corrigé
\newtcolorbox{corrige}[1][]{popbase,colframe=popgreen,colback=popgreenL,
  before upper={\popboxhead{Corrigé}{popgreen}{#1}}}
% Objectifs / prérequis / bilan
\newtcolorbox{objectifs}{popbase,colframe=popink,colback=poporangeL,
  before upper={\popboxhead{Objectifs de la leçon}{popink}{}}}
\newtcolorbox{prerequis}{popbase,colframe=popink,colback=popblueL,
  before upper={\popboxhead{Prérequis}{popblue}{}}}
\newtcolorbox{bilan}{popbase,colframe=popink,colback=white,boxrule=2.8pt,
  before upper={\popboxhead{Fiche bilan}{poporange}{L'essentiel à connaître pour l'examen}}}
% Figure encadrée avec légende
\newtcolorbox{popfigure}[1][]{enhanced,breakable=false,colback=white,colframe=popline,boxrule=1.4pt,arc=8pt,
  left=10pt,right=10pt,top=10pt,bottom=8pt,before skip=10pt,after skip=12pt,halign=center,
  lower separated=false,halign lower=center,
  fontlower=\popsemi\fontsize{9}{11.5}\selectfont\color{popmuted},
  #1}
\newcommand{\legende}[1]{\par\vspace{6pt}{\popsemi\fontsize{9}{11.5}\selectfont\color{popmuted}#1}}

% ============================================================
% Signature OnBuch+
% ============================================================
\newcommand{\poplogoinline}[1]{{\poplogo\fontsize{#1}{#1}\selectfont\color{popink}OnBuch\color{poporange}+}}
\newcommand{\signatureonbuch}{%
  \begin{tcolorbox}[enhanced,colback=popink,colframe=popink,boxrule=2.2pt,arc=10pt,
      drop shadow=poporange,left=14pt,right=14pt,top=10pt,bottom=10pt,before skip=0pt,after skip=0pt]
    \tikz[baseline=-0.7ex]\node[circle,fill=popgreen,draw=popcream,line width=1.3pt,minimum size=22pt,inner sep=0]{\popcheck{white}};%
    \hspace{9pt}\begin{minipage}[c]{0.62\linewidth}
      {\popxbold\fontsize{12}{13}\selectfont\color{popcream}Signé\ \textbullet\ {\poplogo OnBuch\color{poporange}+}}\par\vspace{2pt}
      {\raggedright\popsemi\fontsize{8}{10.5}\selectfont\color{popline}\MakeUppercase{Équipe pédagogique OnBuch+}\\\MakeUppercase{Conforme au programme officiel MINESEC}\par}
    \end{minipage}\hfill
    {\poplogo\fontsize{22}{22}\selectfont\color{popcream}OnBuch\color{poporange}+}
  \end{tcolorbox}%
}

% ============================================================
% Page de garde
% #1 titre · #2 matière · #3 niveau · #4 module · #5 n° leçon · #6 durée ·
% #7 description / objectifs (texte ou itemize)
% ============================================================
\newcommand{\pagegarde}[7]{%
  \thispagestyle{empty}
  \newgeometry{a4paper,margin=1.7cm,top=1.6cm,bottom=1.4cm}%
  \begin{tikzpicture}[remember picture,overlay]
    \fill[poporange!18] ([xshift=-3.2cm,yshift=-3.6cm]current page.north east) circle (4.6cm);
    \fill[poppurple!14] ([xshift=2.4cm,yshift=7.6cm]current page.south west) circle (3.8cm);
  \end{tikzpicture}%
  {\poplogo\fontsize{30}{30}\selectfont\color{popink}OnBuch\color{poporange}+}\hfill
  \raisebox{0.6ex}{\poppill{Cours signé \textbullet\ Premium}{popink}{popcream}}\par
  \vspace{1pt}\popsquiggle{poppurple}\par
  \vspace{8pt}
  \begin{tcolorbox}[enhanced,colback=poporange,colframe=popink,boxrule=2.8pt,arc=13pt,
      drop shadow=popink,left=18pt,right=18pt,top=12pt,bottom=13pt,after skip=10pt]
    \poppill{#2 \textbullet\ #3}{popink}{popcream}\hspace{5pt}\poppill{#4}{white}{popink}\par\vspace{8pt}
    {\popdisplay\fontsize{14}{14}\selectfont\color{popink}LEÇON #5}\par\vspace{4pt}
    {\raggedright\hyphenpenalty=10000\exhyphenpenalty=10000\popdisplay\fontsize{25}{27}\selectfont\color{popcream}\MakeUppercase{#1}\par}
    \vspace{8pt}
    {\popxbold\fontsize{9.5}{9.5}\selectfont\color{popink}\MakeUppercase{Durée conseillée : #6}}
  \end{tcolorbox}
  \begin{tcolorbox}[enhanced,colback=white,colframe=popink,boxrule=2.2pt,arc=10pt,
      drop shadow=popink,left=16pt,right=16pt,top=10pt,bottom=10pt,after skip=8pt]
    \poppill{Dans cette leçon}{poppurple}{white}\par\vspace{5pt}
    \popsemi\fontsize{9.3}{12.6}\selectfont\color{popink}#7
  \end{tcolorbox}
  \noindent\begin{minipage}[t]{0.487\linewidth}
    \begin{tcolorbox}[enhanced,colback=popgreen,colframe=popink,boxrule=2.2pt,arc=10pt,
        drop shadow=popink,left=11pt,right=11pt,top=10pt,bottom=10pt,height=3.1cm,valign=center]
      \tcbox[colback=white,colframe=popink,boxrule=1.3pt,arc=5pt,boxsep=0pt,nobeforeafter,
        left=3pt,right=3pt,top=3pt,bottom=3pt,valign=center]{\qrcode[height=1.9cm]{https://play.google.com/store/apps/details?id=com.luvvix.onbuch&hl=fr}}%
      \hspace{8pt}\begin{minipage}[c]{0.5\linewidth}
        {\popdisplay\fontsize{11}{12.5}\selectfont\color{white}\MakeUppercase{Télécharge OnBuch}}\par\vspace{4pt}
        {\raggedright\popsemi\fontsize{8.4}{11}\selectfont\color{white}Gratuit \textbullet\ Google Play\\Tuteur IA, annales, concours, résultats\par}
      \end{minipage}
    \end{tcolorbox}
  \end{minipage}\hfill
  \begin{minipage}[t]{0.487\linewidth}
    \begin{tcolorbox}[enhanced,colback=poppurple,colframe=popink,boxrule=2.2pt,arc=10pt,
        drop shadow=popink,left=11pt,right=11pt,top=10pt,bottom=10pt,height=3.1cm,valign=center]
      \tikz[baseline=-0.7ex]\node[circle,fill=white,draw=popink,line width=1.3pt,minimum size=1.6cm,inner sep=0]{\popdisplay\fontsize{24}{24}\selectfont\color{poppurple}+};%
      \hspace{8pt}\begin{minipage}[c]{0.6\linewidth}
        {\popdisplay\fontsize{11}{12.5}\selectfont\color{white}\MakeUppercase{Passe à OnBuch+}}\par\vspace{4pt}
        {\raggedright\popsemi\fontsize{8.4}{11}\selectfont\color{white}Tous les cours signés, Tuteur IA illimité, fiches PDF — onglet OnBuch+ de l'app\par}
      \end{minipage}
    \end{tcolorbox}
  \end{minipage}
  \vspace{8pt}
  \popcontenu
  \vfill
  \signatureonbuch
  \restoregeometry
  \newpage
}

% Rangée « Ce cours contient » (page de garde)
\newcommand{\poptile}[3]{% #1 couleur #2 gros texte #3 légende
  \begin{tcolorbox}[enhanced,colback=#1,colframe=popink,boxrule=2pt,arc=9pt,shadow={2.5pt}{-2.5pt}{0pt}{popink},
      left=6pt,right=6pt,top=9pt,bottom=9pt,halign=center,valign=center,height=1.75cm,width=\linewidth,nobeforeafter]
    {\popdisplay\fontsize{12.5}{13}\selectfont\color{white}\MakeUppercase{#2}}\par\vspace{3pt}
    {\centering\popsemi\fontsize{7.8}{9.5}\selectfont\color{white}#3\par}
  \end{tcolorbox}}
\newcommand{\popcontenu}{%
  \par\noindent{\popxboldls\fontsize{7.5}{7.5}\selectfont\color{popmuted}CE COURS SIGNÉ CONTIENT}\par\vspace{7pt}
  \noindent\begin{minipage}[t]{0.235\linewidth}\poptile{popblue}{Cours}{complet, illustré, pas à pas}\end{minipage}\hfill
  \begin{minipage}[t]{0.235\linewidth}\poptile{poporange}{Méthodes}{et exercices résolus}\end{minipage}\hfill
  \begin{minipage}[t]{0.235\linewidth}\poptile{poppink}{Exercices}{type examen + corrigés}\end{minipage}\hfill
  \begin{minipage}[t]{0.235\linewidth}\poptile{popgreen}{Fiche bilan}{l'essentiel à retenir}\end{minipage}\par}

% Sommaire
\newtcolorbox{sommaire}{enhanced,colback=white,colframe=popink,boxrule=2.2pt,arc=10pt,
  drop shadow=popink,left=18pt,right=18pt,top=13pt,bottom=13pt,before skip=6pt,after skip=20pt,
  before upper={\poppill{Sommaire}{poppurple}{white}\par\vspace{9pt}\popsemi\fontsize{10.6}{14.5}\selectfont\color{popink}}}

% Signature de fin de cours — poussée tout en bas de la dernière page
\newcommand{\findecours}{%
  \par\vfill\nopagebreak
  \begin{center}\popsquiggle{poporange}\end{center}\vspace{4pt}
  \signatureonbuch
}
\AtBeginDocument{\def\llbracket{\mathopen{[\mkern-3.5mu[}}\def\rrbracket{\mathclose{]\mkern-3.5mu]}}} % intervalles d'entiers (glyphe ⟦ absent de la police)
% Glyphes absents de Fira Math : redéfinis à partir de glyphes disponibles
\AtBeginDocument{%
  \def\setminus{\mathbin{\backslash}}\let\smallsetminus\setminus
  \def\vdots{\mathord{\vbox{\baselineskip3.2pt\lineskiplimit0pt\kern4pt\hbox{.}\hbox{.}\hbox{.}}}}%
  \def\ddots{\mathinner{\mkern1mu\raise6.5pt\hbox{.}\mkern2mu\raise3.6pt\hbox{.}\mkern2mu\raise0.7pt\hbox{.}\mkern1mu}}%
  \def\triangle{\mathord{\scalebox{0.9}{$\symup{\Delta}$}}}%
  \def\nabla{\mathord{\raisebox{\depth}{\scalebox{1}[-1]{$\symup{\Delta}$}}}}%
  \def\checkmark{\popcheck{popdark}}%
  \def\star{\mathbin{\ast}}\def\bigstar{\mathord{\ast}}%
  \def\diamond{\mathbin{\scalebox{0.6}{$\Diamond$}}}%
  \def\bigcup{\mathop{\mathchoice{\raisebox{-0.3em}{\scalebox{1.7}{$\cup$}}}{\scalebox{1.25}{$\cup$}}{\cup}{\cup}}\displaylimits}%
  \def\bigcap{\mathop{\mathchoice{\raisebox{-0.3em}{\scalebox{1.7}{$\cap$}}}{\scalebox{1.25}{$\cap$}}{\cap}{\cap}}\displaylimits}%
  \def\lesseqqgtr{\mathrel{\lessgtr}}\def\bigoplus{\mathop{\oplus}\displaylimits}%
}
\providecommand{\ssi}{si et seulement si} % abréviation fréquente
\AtBeginDocument{\providecommand{\wideparen}{\overparen}} % arc de cercle

%% FIGFIX-BEGIN v4 — correctifs globaux des figures (docs/onbuch-generation/tools/figfix)
\usepackage{adjustbox}\usepackage{etoolbox}
% toute figure TikZ/pgfplots plus large que la ligne est réduite à la largeur disponible
\BeforeBeginEnvironment{tikzpicture}{\begin{adjustbox}{max width=\linewidth}}
\AfterEndEnvironment{tikzpicture}{\end{adjustbox}}
% légendes pgfplots lisibles par-dessus les courbes
\pgfplotsset{every axis/.append style={legend style={fill=white,fill opacity=0.92,text opacity=1,draw=black!15,inner sep=3pt}}}
% étiquettes d'arêtes (syntaxe edge["..."]) avec fond blanc
\tikzset{every edge quotes/.append style={fill=white,inner sep=1.2pt}}
%% FIGFIX-END
\begin{document}

\pagegarde{Lettres de réclamation et curriculum vitae : forme et contenu}{Allemand}{Tle A}{Chapitre 5 : Citoyenneté}{AL29}{2 h}{Cette leçon mixte initie les élèves de Tle A à deux genres de textes professionnels incontournables : la lettre de réclamation (Beschwerdebrief) et le curriculum vitae (Lebenslauf). Elle combine la compréhension d'un texte support, l'enrichissement du lexique thématique, l'étude du subjonctif II de politesse et de la structure de la lettre formelle, ainsi qu'une production guidée.
\vspace{4pt}
\begin{multicols}{2}\small
\begin{itemize}
\item À la fin de cette leçon, je sais comprendre globalement et en détail une lettre de réclamation et un CV en allemand.
\item À la fin de cette leçon, je sais identifier les rubriques d'un CV allemand (persönliche Daten, Schulbildung, Kenntnisse...).
\item À la fin de cette leçon, je sais employer le vocabulaire thématique : der Lebenslauf, die Bewerbung, die Beschwerde, die Kenntnisse.
\item À la fin de cette leçon, je sais former et utiliser le subjonctif II de politesse (ich hätte gern, könnten Sie, ich würde bitten).
\item À la fin de cette leçon, je sais appliquer la structure et les formules de la lettre formelle allemande (Sehr geehrte Damen und Herren, Mit freundlichen Grüßen).
\item À la fin de cette leçon, je sais distinguer registre formel et informel dans une correspondance.
\end{itemize}\end{multicols}}

\begin{sommaire}
\begin{multicols}{2}
\begin{enumerate}
\item Texte support : Eine Beschwerde und ein Lebenslauf
\item Compréhension du texte
\item Lexique thématique : Bewerbung, Lebenslauf, Beschwerde
\item Grammaire 1 : le subjonctif II de politesse
\item Grammaire 2 : structure et formules de la lettre formelle
\item Production guidée et méthode du commentaire
\item Activité d'intégration
\item Méthodes et exercices résolus
\item Exercices
\item Corrigés des exercices
\item Fiche bilan
\end{enumerate}
\end{multicols}
\end{sommaire}

\begin{objectifs}
\begin{itemize}
\item À la fin de cette leçon, je sais comprendre globalement et en détail une lettre de réclamation et un CV en allemand.
\item À la fin de cette leçon, je sais identifier les rubriques d'un CV allemand (persönliche Daten, Schulbildung, Kenntnisse...).
\item À la fin de cette leçon, je sais employer le vocabulaire thématique : der Lebenslauf, die Bewerbung, die Beschwerde, die Kenntnisse.
\item À la fin de cette leçon, je sais former et utiliser le subjonctif II de politesse (ich hätte gern, könnten Sie, ich würde bitten).
\item À la fin de cette leçon, je sais appliquer la structure et les formules de la lettre formelle allemande (Sehr geehrte Damen und Herren, Mit freundlichen Grüßen).
\item À la fin de cette leçon, je sais distinguer registre formel et informel dans une correspondance.
\item À la fin de cette leçon, je sais répondre à des questions de compréhension et rédiger une courte lettre de réclamation ou un mini-CV.
\item À la fin de cette leçon, je sais traduire des phrases de thème et de version utilisant la politesse et le lexique professionnel.
\end{itemize}
\end{objectifs}

\begin{prerequis}
\begin{itemize}
\item Le subjonctif II présent (formes de base : wäre, hätte, könnte, würde) vu en Première
\item Les verbes modaux au présent et au passé
\item La structure de la phrase allemande (place du verbe, inversion)
\item Le vocabulaire de la vie quotidienne et des achats (kaufen, bestellen, bezahlen)
\item Les pronoms personnels et possessifs
\end{itemize}
\end{prerequis}

\newpage

\coursec{Texte support : Eine Beschwerde und ein Lebenslauf}

Imagine : tu as commandé un téléphone sur un site allemand et le colis arrive abîmé à Douala. Comment réclamer un remboursement en allemand correct et poli ? Et demain, pour un stage à l'ambassade d'Allemagne à Yaoundé ou un emploi en Suisse, on te demandera un \all{Lebenslauf} impeccable. La lettre formelle allemande obéit à des codes précis : formules de politesse, subjonctif II, structure rigoureuse. Celui qui les maîtrise ouvre des portes ; celui qui les ignore reste à la porte.

\textbf{Problématique :} quels sont les éléments de structure et les formules de langue qui font d'une lettre de réclamation et d'un CV allemands des documents efficaces et acceptés ?

\courssub{Compréhension du texte : la lettre de réclamation}

Lisons d'abord la lettre de Jean Mbarga, un client camerounais qui écrit à une entreprise allemande. Lis le texte une première fois sans dictionnaire : cherche seulement \emph{qui} écrit, \emph{à qui}, \emph{pourquoi} et \emph{ce qu'il demande}.

\begin{textebox}[Eine Beschwerde — Brief eines Kunden]
Jean Mbarga\\
Rue 1.234, Quartier Bonanjo\\
Douala, Kamerun\\[0.4em]
TechWorld GmbH\\
Kundenservice\\
Hauptstraße 45\\
10115 Berlin, Deutschland\\[0.4em]
Douala, den 15. Mai 2024\\[0.4em]
\textbf{Betreff: Beschwerde über einen defekten Laptop — Bestellnummer 78432}\\[0.4em]
Sehr geehrte Damen und Herren,\\[0.4em]
vor zwei Wochen habe ich bei Ihnen einen Laptop bestellt. Das Paket ist gestern in Douala angekommen, aber leider funktioniert das Gerät nicht: der Bildschirm bleibt schwarz, und der Akku lädt nicht. Ich bin sehr enttäuscht, denn ich brauche den Laptop dringend für mein Studium.\\[0.4em]
Ich hätte gern einen Ersatz oder eine Rückerstattung des Kaufpreises. Könnten Sie mir bitte mitteilen, wie ich das defekte Gerät zurückschicken kann? Ich würde Sie bitten, mir so bald wie möglich zu antworten.\\[0.4em]
Ich danke Ihnen im Voraus für Ihre Hilfe.\\[0.4em]
Mit freundlichen Grüßen\\
Jean Mbarga
\end{textebox}

\textbf{Glossaire des mots nouveaux} (à noter dans ton cahier de vocabulaire, avec article et pluriel) :

\begin{itemize}
\item \cle{die Beschwerde}, -n : la réclamation, la plainte
\item \cle{der Absender}, - : l'expéditeur ; \cle{der Empfänger}, - : le destinataire
\item \cle{der Betreff}, -e : l'objet (de la lettre)
\item \cle{bestellen} (hat bestellt) : commander ; \cle{die Bestellnummer}, -n : le numéro de commande
\item \cle{das Paket}, -e : le colis ; \cle{das Gerät}, -e : l'appareil
\item \cle{defekt} : défectueux, en panne ; \cle{funktionieren} : fonctionner
\item \cle{der Bildschirm}, -e : l'écran ; \cle{der Akku}, -s : la batterie ; \cle{laden} (lädt, lud, hat geladen) : charger
\item \cle{enttäuscht} : déçu ; \cle{dringend} : urgent(ement)
\item \cle{der Ersatz} : le remplacement ; \cle{die Rückerstattung}, -en : le remboursement ; \cle{der Kaufpreis}, -e : le prix d'achat
\item \cle{zurückschicken} (verbe à particule séparable) : renvoyer, retourner
\item \cle{im Voraus} : d'avance
\end{itemize}

\begin{methode}[Lire un texte formel en trois étapes]
\begin{enumerate}
\item \textbf{Identifier le cadre :} qui écrit (\all{Absender}) ? À qui (\all{Empfänger}) ? Quel est l'objet (\all{Betreff}) ? Ici : un client camerounais écrit au service clients d'une entreprise berlinoise au sujet d'un laptop défectueux.
\item \textbf{Dégager le problème exposé :} repère les faits et les dates. Ici : commande il y a deux semaines, arrivée hier, écran noir, batterie qui ne charge pas.
\item \textbf{Formuler la demande :} que veut exactement l'auteur ? Ici : un remplacement (\all{Ersatz}) ou un remboursement (\all{Rückerstattung}), et une réponse rapide.
\end{enumerate}
Cette démarche en trois temps (cadre, problème, demande) fonctionne pour toute lettre formelle : réclamation, candidature, demande de renseignements.
\end{methode}

\courssub{Compréhension du texte : le curriculum vitae}

Après sa réclamation, Jean prépare un stage : il demande à sa sœur Amina de lui montrer son \all{Lebenslauf}. Voici un extrait.

\begin{textebox}[Lebenslauf — Auszug]
\textbf{Lebenslauf}\\[0.4em]
\textbf{Persönliche Daten}\\
Name: Amina Fotso\\
Geburtsdatum und -ort: geboren am 12.03.2007 in Yaoundé\\
Adresse: Quartier Nkol-Eton, Yaoundé, Kamerun\\
Staatsangehörigkeit: kamerunisch\\[0.4em]
\textbf{Schulbildung}\\
seit 2023: Universität Yaoundé I, Germanistik\\
2023: Baccalauréat série A, Lycée de Nkol-Eton, Yaoundé\\
2017: CEP, École primaire de Nkol-Eton\\[0.4em]
\textbf{Sprachkenntnisse}\\
Französisch: Muttersprache\\
Deutsch: gut (Note 1 am Goethe-Institut Kamerun)\\
Englisch: Grundkenntnisse\\[0.4em]
\textbf{EDV-Kenntnisse}\\
Word, Excel, PowerPoint\\[0.4em]
\textbf{Hobbys}\\
Fußball, Lesen, Kochen
\end{textebox}

\textbf{Glossaire :}

\begin{itemize}
\item \cle{der Lebenslauf}, \all{Lebensläufe} : le curriculum vitae (CV)
\item \cle{die persönlichen Daten} (pluriel) : les données personnelles
\item \cle{das Geburtsdatum}, \all{Geburtsdaten} : la date de naissance ; \cle{der Geburtsort}, -e : le lieu de naissance
\item \cle{die Staatsangehörigkeit}, -en : la nationalité
\item \cle{die Schulbildung} (singulier) : la scolarité, la formation scolaire
\item \cle{die Sprachkenntnisse} (pluriel) : les connaissances linguistiques ; \cle{die Grundkenntnisse} (pluriel) : les notions de base
\item \cle{die Muttersprache}, -n : la langue maternelle
\item \cle{die EDV-Kenntnisse} (pluriel) : les connaissances en informatique (\all{EDV} = \all{elektronische Datenverarbeitung})
\item \cle{das Hobby}, -s : le loisir, le passe-temps
\end{itemize}

\begin{aretenir}[Les cinq rubriques classiques du Lebenslauf]
\begin{enumerate}
\item \all{Persönliche Daten} : nom, date et lieu de naissance, adresse, nationalité.
\item \all{Schulbildung} : la scolarité, \textbf{de l'expérience la plus récente à la plus ancienne} (ordre antichronologique, avec \all{seit} pour ce qui dure encore).
\item \all{Sprachkenntnisse} : les langues avec un niveau : \all{Muttersprache} (langue maternelle), \all{fließend} (courant), \all{gut} (bon), \all{Grundkenntnisse} (notions).
\item \all{EDV-Kenntnisse} : l'informatique.
\item \all{Hobbys} : les loisirs — rubrique courte, deux ou trois éléments sincères.
\end{enumerate}
Le CV allemand tient sur \textbf{une page}, sans phrase complète : des rubriques nominales, des dates en colonne de gauche.
\end{aretenir}

\begin{center}
\begin{tabularx}{\linewidth}{|X|X|X|}
\hline
\rowcolor{popblueL} Français & Deutsch & Remarque \\
\hline
le curriculum vitae & der Lebenslauf, -läufe & mot composé : \all{Leben} + \all{Lauf} \\
\hline
la candidature & die Bewerbung, -en & de \all{sich bewerben} (postuler) \\
\hline
la réclamation & die Beschwerde, -n & de \all{sich beschweren} (se plaindre) \\
\hline
Madame, Monsieur & Sehr geehrte Damen und Herren & salutation impersonnelle standard \\
\hline
cordialement & Mit freundlichen Grüßen & sans virgule finale \\
\hline
les connaissances & die Kenntnisse (Pl.) & presque toujours au pluriel \\
\hline
les notions (de base) & die Grundkenntnisse (Pl.) & niveau le plus bas déclaré \\
\hline
l'expéditeur / le destinataire & der Absender / der Empfänger & paire à mémoriser ensemble \\
\hline
\end{tabularx}
\end{center}

\courssub{Grammaire : Structure et formules de la lettre formelle}

La lettre allemande formelle suit un plan immuable, plus rigide encore que la lettre administrative française. Chaque élément a sa place et son nom :

\begin{popfigure}
\begin{tikzpicture}[
  zone/.style={draw=popblue, thick, fill=popblueL, rounded corners=2pt, minimum width=6.2cm, minimum height=0.7cm, align=center, font=\small},
  etiq/.style={font=\small\bfseries, text=popdark, right=0.25cm, align=left}
]
\node[zone] (abs) at (0,0) {Jean Mbarga, Rue 1.234, Douala};
\node[zone] (emp) at (0,-0.95) {TechWorld GmbH, Berlin};
\node[zone] (dat) at (0,-1.9) {Douala, den 15. Mai 2024};
\node[zone] (bet) at (0,-2.85) {Betreff: Beschwerde über einen defekten Laptop};
\node[zone] (anr) at (0,-3.8) {Sehr geehrte Damen und Herren,};
\node[zone, minimum height=1.1cm, align=center] (txt) at (0,-5.0){Text: Problem + Wunsch\\ (ich hätte gern \ldots)};
\node[zone] (gru) at (0,-6.15) {Mit freundlichen Grüßen};
\node[zone] (unt) at (0,-7.1) {Jean Mbarga};
\node[etiq] at (3.4,0) {\all{der Absender} (expéditeur)};
\node[etiq] at (3.4,-0.95) {\all{der Empfänger} (destinataire)};
\node[etiq] at (3.4,-1.9) {\all{das Datum} (date)};
\node[etiq] at (3.4,-2.85) {\all{der Betreff} (objet)};
\node[etiq] at (3.4,-3.8) {\all{die Anrede} (salutation initiale)};
\node[etiq] at (3.4,-5.0) {\all{der Text} (corps de la lettre)};
\node[etiq] at (3.4,-6.15) {\all{die Grußformel} (formule finale)};
\node[etiq] at (3.4,-7.1) {\all{die Unterschrift} (signature)};
\foreach \n in {abs,emp,dat,bet,anr,txt,gru,unt}
  \draw[-{Stealth}, thick, poporange] (\n.east) -- ++(0.25,0);
\end{tikzpicture}
\legende{Schéma de la lettre formelle allemande : huit éléments, toujours dans le même ordre.}
\end{popfigure}

\begin{propriete}[Règles de forme de la lettre formelle allemande]
\begin{itemize}
\item \textbf{L'expéditeur} (\all{der Absender}) figure en haut à gauche, \textbf{le destinataire} (\all{der Empfänger}) en dessous.
\item \textbf{La date} est précédée du lieu et suit le modèle : \all{Douala, den 15. Mai 2024} — article \all{den} + date ordinale avec un point (\all{15.} = quinzième).
\item \textbf{L'objet} (\all{der Betreff}) annonce en une ligne le sujet ; on écrit souvent simplement \all{Betreff:} suivi du thème.
\item \textbf{La salutation} (\all{die Anrede}) : \all{Sehr geehrte Damen und Herren,} si l'on ne connaît pas le destinataire ; \all{Sehr geehrter Herr Müller,} / \all{Sehr geehrte Frau Ngono,} si on le connaît. Après la virgule, la ligne suivante commence par une \textbf{minuscule} (sauf nom propre ou \all{Sie}).
\item \textbf{La formule finale} (\all{die Grußformel}) : \all{Mit freundlichen Grüßen} (souvent abrégé \all{MfG} dans les courriels), \textbf{sans virgule} à la fin — contrairement au français.
\item \textbf{La signature} (\all{die Unterschrift}) : prénom + nom, sans formule comme « votre serviteur ».
\end{itemize}
\end{propriete}

\begin{attention}[La formule finale : pas de virgule, pas de longueur française]
Le français exige de longues formules (« Je vous prie d'agréer, Madame, Monsieur, l'expression de mes salutations distinguées »). L'allemand est sobre : \all{Mit freundlichen Grüßen} suffit partout. Deux pièges : (1) \textbf{jamais de virgule} après \all{Grüßen} ; (2) ne traduis jamais mot à mot une formule française : \all{Hochachtungsvoll} existe mais sonne très solennel et démodé dans une simple réclamation.
\end{attention}

\courssub{Grammaire : Le subjonctif II de politesse (Konjunktiv II)}

Relisons la lettre de Jean Mbarga en soulignant mentalement toutes les formules de politesse. Trois phrases méritent toute ton attention :

\all{Ich hätte gern einen Ersatz oder eine Rückerstattung.} « J'aimerais un remplacement ou un remboursement. »

\all{Könnten Sie mir bitte mitteilen, wie ich das Gerät zurückschicken kann?} « Pourriez-vous m'indiquer comment je peux renvoyer l'appareil ? »

\all{Ich würde Sie bitten, mir so bald wie möglich zu antworten.} « Je vous prierais de me répondre dès que possible. »

\begin{definition}[Le subjonctif II (Konjunktiv II) : formation et emploi]
Le \cle{Konjunktiv II} exprime l'irréel, le souhait, la politesse. En correspondance formelle, il remplace l'indicatif pour adoucir la demande (équivalent du conditionnel français). Deux formations coexistent :
\begin{enumerate}
\item \textbf{Forme simple (Präteritum + Umlaut + endings) :} pour \all{haben, sein, werden} et les verbes modaux.
\item \textbf{Forme composée (würde + infinitif) :} pour tous les autres verbes (surtout à l'oral et à l'écrit courant).
\end{enumerate}
\end{definition}

\begin{propriete}[Conjugaison du Konjunktiv II — Verbes auxiliaires et modaux (Présent)]
\begin{center}
\begin{tabular}{|l|c|c|c|c|c|c|}
\hline
\rowcolor{popblueL} Personne & \all{haben} & \all{sein} & \all{werden} & \all{können} & \all{müssen} & \all{wollen} \\
\hline
ich & hätte & wäre & würde & könnte & müsste & wollte \\
du & hättest & wärest & würdest & könntest & müsstest & wolltest \\
er/sie/es & hätte & wäre & würde & könnte & müsste & wollte \\
wir & hätten & wären & würden & könnten & müssten & wollten \\
ihr & hättet & wäret & würdet & könntet & müsstet & wolltet \\
sie/Sie & hätten & wären & würden & könnten & müssten & wollten \\
\hline
\end{tabular}
\end{center}
\textbf{Autres modaux :} \all{dürfen} $\to$ \all{dürfte}, \all{sollen} $\to$ \all{sollte}, \all{mögen} $\to$ \all{möchte} (forme figée très fréquente : \all{ich möchte} = « je voudrais »).
\end{propriete}

\begin{propriete}[Forme composée : \all{würde} + infinitif]
Pour les verbes pleins (sauf \all{haben, sein, werden, modaux}), on utilise \all{würde} + infinitif à la fin de la phrase.
\begin{itemize}
\item \all{Ich würde gern kommen.} (J'aimerais venir.)
\item \all{Würden Sie mir helfen?} (M'aideriez-vous ?)
\item \all{Er würde das nicht machen.} (Il ne ferait pas ça.)
\end{itemize}
\emph{Remarque :} À l'écrit formel, on préfère souvent la forme simple pour les verbes courants (\all{ginge, käme, fände}), mais \all{würde} + infinitif est toujours correct et plus sûr pour l'élève.
\end{propriete}

\begin{exemplebox}[Le subjonctif II dans la lettre de Jean Mbarga]
\begin{itemize}
\item \all{Ich \textbf{hätte} gern \ldots} (verbe \all{haben}, forme simple) — « J'aimerais… »
\item \all{\textbf{Könnten} Sie mir \ldots mitteilen?} (verbe \all{können}, forme simple, inversion sujet-verbe pour la question polie) — « Pourriez-vous… ? »
\item \all{Ich \textbf{würde} Sie bitten, \ldots zu antworten.} (verbe \all{bitten} via \all{würde} + infinitif) — « Je vous prierais… »
\end{itemize}
\end{exemplebox}

\courssub{Le ton de la lettre de réclamation : courtoisie et fermeté}

\begin{propriete}[Le ton de la lettre de réclamation : courtois mais ferme]
Une bonne réclamation allemande combine deux registres :
\begin{itemize}
\item \textbf{La courtoisie} : subjonctif II (\all{ich hätte gern}, \all{könnten Sie}), remerciement anticipé (\all{Ich danke Ihnen im Voraus}), salutations complètes (\all{Mit freundlichen Grüßen}) ;
\item \textbf{La fermeté} : faits précis et datés (\all{vor zwei Wochen}, \all{gestern}), demande explicite (\all{Ersatz oder Rückerstattung}), délai demandé (\all{so bald wie möglich}).
\end{itemize}
On n'insulte jamais, on ne menace pas : on constate, on demande, on fixe un délai. C'est ce ton que le correcteur du Bac attend dans ta production écrite.
\end{propriete}

\begin{attention}[Faux amis capitaux : \all{die Beschwerde} et \all{die Bewerbung}]
Ces deux noms se ressemblent et piègent chaque année des candidats au Bac :
\begin{itemize}
\item \cle{die Beschwerde}, -n = la \textbf{réclamation}, la plainte (verbe : \all{sich beschweren über} + accusatif, « se plaindre de ») ;
\item \cle{die Bewerbung}, -en = la \textbf{candidature} (verbe : \all{sich bewerben um} + accusatif, « postuler à »).
\end{itemize}
\all{Ich schreibe eine Beschwerde.} « J'écris une réclamation. » — \all{Ich schreibe eine Bewerbung.} « J'écris une candidature. » Mémorise aussi la différence de préposition : on se plaint \emph{de} (\all{über}), on postule \emph{à} (\all{um}).
\end{attention}

\courssub{Méthode : Rédiger une lettre de réclamation formelle}

\begin{methode}[Écrire une Beschwerde en 5 étapes]
\begin{enumerate}
\item \textbf{En-tête :} \all{Absender} (toi), \all{Empfänger} (entreprise), \all{Datum} (lieu, \all{den} + date).
\item \textbf{Objet :} \all{Betreff:} + nom du problème + numéro de commande/référence.
\item \textbf{Salutation :} \all{Sehr geehrte Damen und Herren,} (ou nom si connu).
\item \textbf{Corps du texte (3 paragraphes types) :}
      \begin{enumerate}
      \item \textbf{Contexte :} \all{Am [Datum] habe ich \ldots bestellt / gekauft.} (Indicatif Perfekt/Präteritum).
      \item \textbf{Problème :} \all{Leider ist \ldots defekt / fehlt / stimmt nicht.} (Indicatif Présent).
      \item \textbf{Demande (Politesse) :} \all{Ich hätte gern \ldots / Könnten Sie \ldots / Ich würde Sie bitten, \ldots} (Konjunktiv II).
      \end{enumerate}
\item \textbf{Clôture :} \all{Ich danke Ihnen im Voraus.} + \all{Mit freundlichen Grüßen} + \textbf{Signature manuscrite} + Nom tapé.
\end{enumerate}
\end{methode}

\begin{experience}[Jeu de rôle : « Au service clients »]
\textbf{Consigne :} En binôme. L'élève A est le client (Jean Mbarga ou toi-même), l'élève B est le service clients de TechWorld GmbH.
\begin{enumerate}
\item L'élève A téléphone (oral) ou écrit un courriel (écrit) pour réclamer : article commandé, défaut constaté, demande précise (échange/remboursement), délai.
\item L'élève B répond : accuse réception, propose solution (étiquette retour, remboursement, délai), formule de politesse.
\item \textbf{Contraintes linguistiques :} Utiliser \all{Sie}-Form, \all{Konjunktiv II} (\all{hätte gern, könnte, würde}), vocabulaire du glossaire.
\end{enumerate}
\textbf{Évaluation par les pairs :} Structure respectée ? Subjonctif II présent ? Ton courtois/ferme ? Vocabulaire précis ?
\end{experience}

\courssub{Production guidée et exercices}

\begin{exoresolu}[Analyser et corriger une lettre]
Énoncé : Voici une lettre d'élève. Identifie 3 erreurs de forme (structure, ponctuation, majuscules) et 2 erreurs de langue (grammaire, vocabulaire), puis réécris-la correctement.\\[0.3em]
\all{paul essono\\
yaoundé\\
an sportschuh gmbh\\
münchen\\
3. april 2024\\
betreff: falsche lieferung\\
sehr geehrte damen und herren,\\
ich habe schuhe bestellt. die schuhe sind zu klein. ich will umtausch. antworten sie mir schnell.\\
mit freundlichen grüßen,\\
paul essono}
\tcblower
\textbf{Solution détaillée :}
\textbf{Erreurs de forme :} (1) \all{Absender} et \all{Empfänger} sans majuscules aux noms propres (\all{Paul Essono}, \all{Yaoundé}, \all{SportSchuh GmbH}, \all{München}). (2) Date incorrecte : il faut \all{Yaoundé, den 3. April 2024}. (3) Virgule interdite après \all{Grüßen}.
\textbf{Erreurs de langue :} (1) \all{ich will umtausch} : trop brutal, pas de politesse $\to$ \all{Ich hätte gern einen Umtausch}. (2) \all{antworten sie mir schnell} : impératif familier/brutal $\to$ \all{Könnten Sie mir bitte bald antworten?}.
\textbf{Version corrigée :}
\all{Paul Essono\\
Quartier Melen, Yaoundé\\
SportSchuh GmbH\\
Hauptstraße 10\\
80331 München\\
Yaoundé, den 3. April 2024\\
Betreff: Beschwerde über falsche Lieferung — Bestellnr. 4711\\
Sehr geehrte Damen und Herren,\\
am 1. April habe ich bei Ihnen Schuhe in Größe 42 bestellt (Bestellnummer 4711). Leider sind die gelieferten Schuhe in Größe 40 zu klein. Ich hätte gern einen Umtausch in die richtige Größe. Könnten Sie mir bitte mitteilen, wie ich die falschen Schuhe zurückschicken kann? Ich würde Sie bitten, mir bald zu antworten.\\
Ich danke Ihnen im Voraus.\\
Mit freundlichen Grüßen\\
Paul Essono}
\end{exoresolu}

\begin{exercice}[À toi de jouer]
\begin{enumerate}
\item Dans la lettre de Jean Mbarga, relève trois marques de courtoisie et deux marques de fermeté.
\item Remets dans l'ordre les éléments d'une lettre : \all{Grußformel — Betreff — Absender — Anrede — Unterschrift — Empfänger — Datum — Text}.
\item Conjugue au Konjunktif II (forme simple ou composée selon le verbe) : \all{ich (haben) gern}, \all{Sie (können) mir helfen?}, \all{ich (bitten) Sie, warten}, \all{er (gehen) nach Hause}, \all{wir (sein) zufrieden}.
\item Traduis en allemand (lettre formelle) : « Monsieur, Madame, j'ai commandé un livre le 10 mai. Il est arrivé endommagé. Je voudrais un remboursement. Pourriez-vous m'envoyer une étiquette de retour ? Je vous prie de me répondre vite. Cordialement, Marie Njoya. »
\item Complète le CV d'Amina avec les rubriques manquantes : elle parle aussi espagnol (notions) et aime la lecture et le chant.
\end{enumerate}
\end{exercice}

\begin{corrige}[Exercice « À toi de jouer »]
\begin{enumerate}
\item Courtoisie : \all{Sehr geehrte Damen und Herren}, \all{Ich hätte gern}, \all{Könnten Sie bitte}, \all{Ich danke Ihnen im Voraus}, \all{Mit freundlichen Grüßen}. Fermeté : faits datés (\all{vor zwei Wochen}, \all{gestern}), demande claire (\all{Ersatz oder Rückerstattung}), délai (\all{so bald wie möglich}).
\item Ordre correct : \all{Absender — Empfänger — Datum — Betreff — Anrede — Text — Grußformel — Unterschrift}.
\item \all{ich hätte gern}, \all{Könnten Sie mir helfen?}, \all{ich würde Sie bitten, zu warten}, \all{er ginge nach Hause} (ou \all{er würde nach Hause gehen}), \all{wir wären zufrieden}.
\item \all{Marie Njoya\\
Quartier Bastos, Yaoundé\\
Buchhandlung Lesen \& Wissen\\
Hauptstraße 5\\
10117 Berlin\\
Yaoundé, den 15. Mai 2024\\
Betreff: Beschwerde über beschädigte Lieferung — Bestellnr. 8821\\
Sehr geehrte Damen und Herren,\\
am 10. Mai habe ich ein Buch bei Ihnen bestellt (Bestellnummer 8821). Leider ist das Buch beschädigt angekommen: der Einband ist zerrissen. Ich hätte gern eine Rückerstattung des Kaufpreises. Könnten Sie mir bitte ein Rücksendeetikett zusenden? Ich würde Sie bitten, mir so bald wie möglich zu antworten.\\
Ich danke Ihnen im Voraus.\\
Mit freundlichen Grüßen\\
Marie Njoya}
\item Rubrique \all{Sprachkenntnisse} : ajouter \all{Spanisch: Grundkenntnisse}. Rubrique \all{Hobbys} : \all{Lesen, Singen}.
\end{enumerate}
\end{corrige}

\begin{savaistu}[Le CV allemand : une tradition de précision]
En Allemagne, en Autriche et en Suisse, le \all{Lebenslauf} était longtemps signé et daté à la main en bas de page — la signature attestait l'exactitude des informations. Aujourd'hui encore, beaucoup d'employeurs germanophones apprécient un CV d'une page, sans faute, avec les dates alignées en colonne. Au Cameroun, le Goethe-Institut de Yaoundé et les entreprises germanophones (ambassades, ONG, firmes logistiques de Douala) attendent ce format : maîtriser le \all{Lebenslauf} est un vrai passeport professionnel.
\end{savaistu}

\begin{aretenir}[Ce qu'il faut retenir de cette leçon (AL29)]
\begin{itemize}
\item La lettre formelle allemande compte huit éléments dans un ordre fixe : \all{Absender, Empfänger, Datum, Betreff, Anrede, Text, Grußformel, Unterschrift}.
\item Le \all{Lebenslauf} se construit en rubriques nominales : \all{Persönliche Daten, Schulbildung, Sprachkenntnisse, EDV-Kenntnisse, Hobbys}, en ordre antichronologique.
\item La politesse passe par le \cle{Konjunktiv II} : \all{ich hätte gern}, \all{könnten Sie}, \all{ich würde Sie bitten}, \all{ich möchte}.
\item Formation Konjunktif II : formes simples (\all{hätte, wäre, würde, könnte, müsste, wollte, sollte, dürfte, möchte}) ou \all{würde} + infinitif.
\item Le ton d'une réclamation est courtois mais ferme : faits datés, demande explicite, délai.
\item Ne jamais confondre \all{die Beschwerde} (réclamation, \all{über} + Akk.) et \all{die Bewerbung} (candidature, \all{um} + Akk.).
\item Pas de virgule après \all{Mit freundlichen Grüßen}.
\end{itemize}
\end{aretenir}

\coursec{Compréhension du texte}

Dans la section précédente, nous avons lu et observé le texte support \all{Eine Beschwerde und ein Lebenslauf} : une lettre de réclamation adressée par un client à une entreprise, suivie du curriculum vitae d'une jeune diplômée. Nous allons maintenant vérifier et approfondir notre compréhension, étape par étape : d'abord une lecture globale (qui écrit, à qui, pourquoi), puis une lecture détaillée (quel produit, quel problème, quelle demande, quelles informations dans le CV). C'est exactement la démarche exigée à l'examen : comprendre globalement, puis repérer les informations précises, enfin rédiger des réponses en phrases complètes.

\begin{textebox}[Rappel du texte support : Eine Beschwerde und ein Lebenslauf]
Sehr geehrte Damen und Herren,\\
am 12. März habe ich in Ihrem Online-Shop einen Laptop der Marke „TechnoStar“ bestellt. Die Lieferung kam zwar schnell, aber leider funktioniert der Bildschirm nicht: Er bleibt schwarz, obwohl das Gerät eingeschaltet ist.\\
Ich bitte Sie, mir so schnell wie möglich ein neues Gerät zu schicken oder mir das Geld zurückzuerstatten. Eine Kopie der Rechnung liegt bei.\\
Ich hoffe auf eine schnelle Lösung und danke Ihnen im Voraus.\\
Mit freundlichen Grüßen\\
Jean Mbarga, Yaoundé

\textbf{Lebenslauf — Amina Ousmanou}\\
Ausbildung: Abitur am Gymnasium in Maroua; danach Studium der Wirtschaftswissenschaften an der Universität Jaunde.\\
Sprachkenntnisse: Deutsch (gut), Englisch (Grundkenntnisse), Französisch (Muttersprache).\\
Computerkenntnisse: Word, Excel, PowerPoint.\\
Praktische Erfahrung: ein dreimonatiges Praktikum in einer Bank in Duala.
\end{textebox}

\textbf{Glossaire :} der Online-Shop, -s (boutique en ligne) ; die Lieferung, -en (livraison) ; der Bildschirm, -e (écran) ; das Gerät, -e (appareil) ; zurückzuerstatten (rembourser) ; die Rechnung, -en (facture) ; die Ausbildung, -en (formation) ; die Sprachkenntnisse (pl., connaissances linguistiques) ; die Grundkenntnisse (pl., notions de base) ; das Praktikum, die Praktika (stage) ; die Lösung, -en (solution) ; die Erfahrung, -en (expérience).

\courssub{Compréhension globale : Wer? Wem? Warum?}

La première étape consiste à répondre à trois questions fondamentales : \textbf{qui écrit ?} (\all{Wer schreibt?}), \textbf{à qui ?} (\all{An wen?}), \textbf{pourquoi ?} (\all{Warum?}). Ces questions permettent de situer la situation de communication avant d'entrer dans les détails. Remarquez la construction : \all{an wen} (à qui) exige l'accusatif après la préposition \all{an} quand elle exprime le destinataire.

\begin{exemplebox}[Réponses globales]
\all{Wer schreibt den Brief? — Jean Mbarga aus Yaoundé schreibt den Brief.} « Qui écrit la lettre ? — Jean Mbarga, de Yaoundé, écrit la lettre. »\\
\all{An wen schreibt er? — Er schreibt an den Online-Shop (an die Firma).} « À qui écrit-il ? — Il écrit à la boutique en ligne (à l'entreprise). »\\
\all{Warum schreibt er? — Er schreibt, weil der Laptop nicht funktioniert. Er möchte sich beschweren.} « Pourquoi écrit-il ? — Parce que l'ordinateur ne fonctionne pas. Il veut se plaindre. »
\end{exemplebox}

\begin{attention}[Le verbe réfléchi \all{sich beschweren}]
\all{Sich beschweren} (se plaindre) est un verbe réfléchi : \all{ich beschwere mich}, \all{er beschwert sich}. Il se construit avec \all{bei + Dativ} (auprès de qui) et \all{über + Akkusativ} (à propos de quoi) : \all{Er beschwert sich bei der Firma über den defekten Laptop.} « Il se plaint auprès de l'entreprise au sujet de l'ordinateur défectueux. » N'oubliez jamais le pronom réfléchi : \all{Er beschwert sich}, et non \all{Er beschwert}.
\end{attention}

\begin{popfigure}
\begin{center}
\begin{tikzpicture}[mindmap, grow cyclic, every node/.style={concept}, concept color=popblueL, root concept/.append style={concept color=popblue, text=white, font=\bfseries}, level 1/.append style={level distance=3.4cm, sibling angle=90, concept color=poporangeL}, text=popdark]
\node[root concept]{Beschwerdebrief}
  child { node[concept, align=center]{Wer?\\ Jean Mbarga} }
  child { node[concept, align=center]{Wem?\\ der Online-Shop} }
  child { node[concept, align=center]{Warum?\\ Laptop defekt} }
  child { node[concept, align=center]{Was will der\\ Schreiber?\\ neues Gerät oder\\ Geld zurück} };
\end{tikzpicture}
\end{center}
\legende{Carte mentale de compréhension globale : les quatre questions clés autour du Beschwerdebrief.}
\end{popfigure}

\courssub{Compréhension détaillée : la lettre de réclamation}

Après la lecture globale, on cherche les informations précises. Pour une lettre de réclamation, il faut toujours repérer trois éléments : \cle{das Produkt} (le produit), \cle{das Problem} (le problème), \cle{die Forderung} (la demande). On ajoute souvent la date de l'achat et les pièces jointes (\all{die Beilagen}).

\begin{exemplebox}[Questions détaillées sur la lettre]
\all{Was hat der Kunde bestellt? — Er hat einen Laptop bestellt.} « Qu'a commandé le client ? — Il a commandé un ordinateur portable. »\\
\all{Wann hat er den Laptop bestellt? — Er hat ihn am 12. März bestellt.} « Quand l'a-t-il commandé ? — Il l'a commandé le 12 mars. »\\
\all{Was ist das Problem? — Der Bildschirm funktioniert nicht; er bleibt schwarz.} « Quel est le problème ? — L'écran ne fonctionne pas ; il reste noir. »\\
\all{Was will der Kunde? — Er will ein neues Gerät oder sein Geld zurück.} « Que veut le client ? — Il veut un nouvel appareil ou le remboursement de son argent. »\\
\all{Was liegt dem Brief bei? — Eine Kopie der Rechnung liegt bei.} « Qu'est-ce qui est joint à la lettre ? — Une copie de la facture est jointe. »
\end{exemplebox}

\begin{center}
\begin{tabularx}{\linewidth}{|X|X|X|}
\hline
\rowcolor{popblueL} Question & Information du texte & Réponse complète \\
\hline
Welches Produkt? & einen Laptop der Marke „TechnoStar“ & Er hat einen Laptop bestellt. \\
\hline
Welches Problem? & der Bildschirm bleibt schwarz & Der Bildschirm funktioniert nicht. \\
\hline
Welche Forderung? & ein neues Gerät / Geld zurück & Er will Ersatz oder eine Rückerstattung. \\
\hline
Welche Beilage? & eine Kopie der Rechnung & Eine Rechnungskopie liegt bei. \\
\hline
\end{tabularx}
\end{center}

\courssub{Compréhension détaillée : le curriculum vitae}

Pour le CV, les questions portent sur la formation (\all{die Ausbildung}), les langues (\all{die Sprachkenntnisse}), les compétences informatiques (\all{die Computerkenntnisse}) et l'expérience (\all{die praktische Erfahrung}). Ce sont les quatre rubriques classiques du \all{Lebenslauf} allemand, avec les données personnelles (\all{die persönlichen Daten}).

\begin{exemplebox}[Questions détaillées sur le CV]
\all{Welche Ausbildung hat Amina? — Sie hat das Abitur in Maroua gemacht und danach in Jaunde Wirtschaftswissenschaften studiert.} « Quelle formation Amina a-t-elle suivie ? — Elle a passé le baccalauréat à Maroua, puis a étudié les sciences économiques à Yaoundé. »\\
\all{Welche Sprachkenntnisse hat Amina? — Sie spricht gut Deutsch und hat Grundkenntnisse in Englisch.} « Quelles connaissances linguistiques a Amina ? — Elle parle bien allemand et a des notions de base en anglais. »\\
\all{Welche Computerkenntnisse hat sie? — Sie kann Word, Excel und PowerPoint benutzen.} « Quelles compétences informatiques a-t-elle ? — Elle sait utiliser Word, Excel et PowerPoint. »\\
\all{Welche praktische Erfahrung hat sie? — Sie hat ein dreimonatiges Praktikum in einer Bank in Duala gemacht.} « Quelle expérience pratique a-t-elle ? — Elle a fait un stage de trois mois dans une banque à Douala. »
\end{exemplebox}

\begin{methode}[Répondre à une question de compréhension]
\begin{enumerate}
\item \textbf{Repérer} dans la question le mot interrogatif (\all{wer}, \all{wem}, \all{warum}, \all{was}, \all{wann}, \all{welche}) : il annonce le type d'information attendu.
\item \textbf{Reprendre les mots de la question} pour construire le début de la réponse : \all{Was hat der Kunde bestellt?} devient \all{Der Kunde hat ... bestellt.}
\item \textbf{Compléter avec l'information du texte}, en la reformulant brièvement, sans copier de longs passages.
\item \textbf{Vérifier} que la réponse est une phrase complète : sujet + verbe conjugué en 2e position.
\end{enumerate}
\end{methode}

\begin{attention}[La place du verbe dans la réponse]
À l'oral comme à l'écrit, une réponse allemande commence par le sujet suivi du verbe conjugué en \textbf{2e position}. Ne traduisez jamais mot à mot le français.\\
Erreur fréquente : \all{Hat bestellt einen Laptop.} (faux : il manque le sujet et le verbe n'est pas en 2e position).\\
Correct : \all{Der Kunde hat einen Laptop bestellt.}\\
Autre piège : avec le Perfekt, le participe passé va à la fin : \all{Er hat einen Laptop bestellt}, jamais \all{Er hat bestellt einen Laptop}.\\
Dans une subordonnée introduite par \all{weil}, le verbe conjugué va à la fin : \all{Er schreibt, weil der Laptop nicht funktioniert.}
\end{attention}

\courssub{Vrai ou faux : justification par le texte}

L'exercice « vrai ou faux » (\all{Richtig oder falsch?}) exige de justifier chaque réponse par une phrase du texte. C'est un entraînement précieux à la citation précise, très utile pour le commentaire de texte au Baccalauréat.

\begin{exoresolu}[Richtig oder falsch?]
Énoncé. Dites si les phrases suivantes sont vraies ou fausses et justifiez par le texte.\\
1. Der Laptop kam nach drei Monaten an.\\
2. Der Bildschirm funktioniert nicht.\\
3. Der Kunde will ein neues Gerät oder sein Geld zurück.\\
4. Amina spricht fließend Englisch.\\
5. Amina hat in einer Bank gearbeitet.
\tcblower
Solution détaillée.\\
1. \textbf{Falsch.} Im Text steht: „Die Lieferung kam zwar schnell.“ (La livraison est arrivée vite, pas après trois mois.)\\
2. \textbf{Richtig.} Im Text steht: „Leider funktioniert der Bildschirm nicht: Er bleibt schwarz.“\\
3. \textbf{Richtig.} Der Kunde schreibt: „Ich bitte Sie, mir so schnell wie möglich ein neues Gerät zu schicken oder mir das Geld zurückzuerstatten.“\\
4. \textbf{Falsch.} Im Lebenslauf steht: „Englisch (Grundkenntnisse)“ — sie hat nur Grundkenntnisse, sie spricht nicht fließend Englisch.\\
5. \textbf{Richtig (teilweise).} Sie hat ein dreimonatiges Praktikum in einer Bank in Duala gemacht — es war ein Praktikum, keine feste Arbeit.
\end{exoresolu}

\courssub{Reformulation : résumer la lettre (Zusammenfassung)}

Résumer (\all{zusammenfassen}) signifie : reprendre uniquement les informations essentielles (qui, quoi, problème, demande) en deux ou trois phrases, avec ses propres mots, sans détails secondaires ni citations longues. Le résumé se rédige au Präsens ou au Perfekt, jamais à la première personne si l'on n'est pas l'auteur de la lettre.

\begin{exemplebox}[Zusammenfassung en deux phrases]
\all{Jean Mbarga aus Yaoundé beschwert sich bei einem Online-Shop, weil der bestellte Laptop einen defekten Bildschirm hat. Er verlangt ein neues Gerät oder die Rückerstattung seines Geldes.}\\
« Jean Mbarga de Yaoundé se plaint auprès d'une boutique en ligne parce que l'ordinateur portable commandé a un écran défectueux. Il exige un nouvel appareil ou le remboursement de son argent. »
\end{exemplebox}

\begin{exercice}[À vous de répondre]
Répondez en phrases complètes en allemand.\\
1. Wann hat Jean den Laptop bestellt?\\
2. Warum ist der Kunde nicht zufrieden?\\
3. Was liegt dem Brief bei?\\
4. Welche Sprachen spricht Amina?\\
5. Wo hat Amina ihr Praktikum gemacht?\\
6. Fassen Sie den Brief in zwei Sätzen zusammen.
\end{exercice}

\begin{corrige}[Exercice « À vous de répondre »]
1. Er hat den Laptop am 12. März bestellt.\\
2. Er ist nicht zufrieden, weil der Bildschirm nicht funktioniert und schwarz bleibt.\\
3. Eine Kopie der Rechnung liegt dem Brief bei.\\
4. Sie spricht gut Deutsch, hat Grundkenntnisse in Englisch, und Französisch ist ihre Muttersprache.\\
5. Sie hat ihr Praktikum in einer Bank in Duala gemacht.\\
6. Modèle : \all{Ein Kunde beschwert sich über einen defekten Laptop. Er möchte ein neues Gerät oder sein Geld zurück.}
\end{corrige}

\begin{aretenir}[Les étapes de la compréhension]
\begin{itemize}
\item \textbf{Lecture globale} : répondre à \all{Wer? Wem? Warum?}
\item \textbf{Lecture détaillée} : repérer produit, problème, demande (lettre) ; formation, langues, compétences, expérience (CV).
\item \textbf{Rédaction} : phrases complètes, verbe en 2e position, participe passé à la fin au Perfekt, verbe à la fin après \all{weil}.
\item \textbf{Résumé} : deux phrases, ses propres mots, uniquement l'essentiel.
\end{itemize}
\end{aretenir}

\coursec{Lexique thématique : Bewerbung, Lebenslauf, Beschwerde}

Après avoir lu et compris la lettre de réclamation et le curriculum vitae de la section précédente, nous allons maintenant \cle{organiser et approfondir le vocabulaire} qui permet de parler du monde du travail, des candidatures et des réclamations. Ce lexique est indispensable : il revient dans presque toutes les épreuves d'expression écrite de Terminale (lettre formelle, résumé, thème). La règle d'or en allemand : \textbf{on n'apprend jamais un nom seul}, mais toujours avec son \cle{article} (der, die, das) et son \cle{pluriel}.

\courssub{Le champ lexical du CV : der Lebenslauf}

\begin{definition}[der Lebenslauf]
\all{Der Lebenslauf} (pluriel : \all{die Lebensläufe}) est le document qui présente de manière structurée la \textbf{formation}, l'\textbf{expérience professionnelle} et les \textbf{compétences} d'une personne. C'est l'équivalent allemand du « CV » français (curriculum vitae). On dit aussi \all{einen Lebenslauf schreiben} (rédiger un CV) ou \all{einen Lebenslauf beifügen} (joindre un CV).
\end{definition}

Observons d'abord les rubriques types d'un CV allemand, que vous avez déjà rencontrées dans le texte support :

\begin{center}
\begin{tabularx}{\linewidth}{|X|X|X|}
\hline
\rowcolor{popblueL} \textbf{Deutsch (avec pluriel)} & \textbf{Français} & \textbf{Exemple d'emploi} \\
\hline
\all{der Lebenslauf, -läufe} & le CV & \all{Ich füge meinen Lebenslauf bei.} \\
\hline
\all{die persönlichen Daten (Pl.)} & les données personnelles & \all{Name, Adresse und Geburtsdatum} \\
\hline
\all{die Schulbildung (Sg.)} & la formation scolaire & \all{meine Schulbildung in Yaoundé} \\
\hline
\all{die Berufserfahrung, -en} & l'expérience professionnelle & \all{zwei Jahre Berufserfahrung} \\
\hline
\all{die Kenntnisse (Pl.)} & les connaissances, compétences & \all{gute EDV-Kenntnisse} \\
\hline
\all{die Sprachkenntnisse (Pl.)} & les connaissances linguistiques & \all{Sprachkenntnisse: Deutsch, Englisch} \\
\hline
\all{das Hobby, -s} & le loisir, le passe-temps & \all{Meine Hobbys sind Fußball und Musik.} \\
\hline
\all{der Abschluss, -schlüsse} & le diplôme & \all{der Abschluss an der Universität} \\
\hline
\all{das Zeugnis, -se} & le bulletin, le certificat & \all{das Abiturzeugnis} \\
\hline
\all{die Unterschrift, -en} & la signature & \all{mit Unterschrift und Datum} \\
\hline
\end{tabularx}
\end{center}

\begin{attention}[« die Kenntnisse » : un pluriel obligatoire]
\all{Die Kenntnis} (la connaissance) s'emploie presque toujours \textbf{au pluriel} quand on parle de compétences : \all{Sprachkenntnisse} (connaissances linguistiques), \all{EDV-Kenntnisse} (connaissances en informatique), \all{gute Kenntnisse in Mathematik}. Ne dites jamais \all{*meine Sprachkenntnis}. En français, le mot « connaissance » existe au singulier, mais pour exprimer des compétences, on utilise presque toujours le pluriel « mes connaissances » ; l'allemand est tout aussi strict : \all{die Kenntnisse} (pluriel) est obligatoire.
\end{attention}

\begin{exemplebox}[Le CV en phrases]
\all{Ich habe meinen Abschluss 2022 an der Universität Yaoundé erworben.} « J'ai obtenu mon diplôme en 2022 à l'université de Yaoundé. »

\all{Meine Sprachkenntnisse: Deutsch (fließend), Englisch (gut), Französisch (Muttersprache).} « Mes connaissances linguistiques : allemand (courant), anglais (bon), français (langue maternelle). »

\all{Ich habe zwei Jahre Berufserfahrung als Verkäuferin.} « J'ai deux ans d'expérience professionnelle comme vendeuse. »
\end{exemplebox}

\courssub{Le champ lexical de la candidature : die Bewerbung}

\begin{definition}[die Bewerbung]
\all{Die Bewerbung} (pluriel : \all{die Bewerbungen}) désigne la \textbf{candidature}, c'est-à-dire l'ensemble des documents (lettre de motivation + CV + diplômes) qu'on envoie pour obtenir un poste. Le verbe correspondant est \all{sich bewerben um + Akkusativ} : « poser sa candidature pour / postuler à ».
\end{definition}

\begin{center}
\begin{tabularx}{\linewidth}{|X|X|X|}
\hline
\rowcolor{popgreenL} \textbf{Deutsch (avec pluriel)} & \textbf{Français} & \textbf{Exemple d'emploi} \\
\hline
\all{die Bewerbung, -en} & la candidature & \all{Ihre Bewerbung ist eingegangen.} \\
\hline
\all{das Bewerbungsschreiben, -} & la lettre de motivation & \all{ein überzeugendes Bewerbungsschreiben} \\
\hline
\all{der Bewerber, - / die Bewerberin, -nen} & le / la candidat(e) & \all{Viele Bewerber für eine Stelle.} \\
\hline
\all{die Stelle, -n} & le poste, l'emploi & \all{eine freie Stelle} \\
\hline
\all{die Stellenanzeige, -n} & l'offre d'emploi & \all{eine Stellenanzeige in der Zeitung} \\
\hline
\all{das Vorstellungsgespräch, -e} & l'entretien d'embauche & \all{zum Vorstellungsgespräch einladen} \\
\hline
\all{die Anlage, -n} & la pièce jointe & \all{Anlagen: Lebenslauf, Zeugnisse} \\
\hline
\all{der Arbeitgeber, -} & l'employeur & \all{ein großer Arbeitgeber in Douala} \\
\hline
\end{tabularx}
\end{center}

\begin{propriete}[« sich bewerben um » + accusatif]
Le verbe réfléchi \all{sich bewerben} se construit avec la préposition \all{um} suivie de l'\textbf{accusatif} : elle introduit \textbf{la chose visée} (le poste, la bourse, la place). Attention : \all{bewerben} est un verbe \textbf{fort} : \all{ich bewerbe mich, ich bewarb mich, ich habe mich beworben}.

\all{Ich bewerbe mich um die Stelle als Übersetzer.} « Je pose ma candidature pour le poste de traducteur. »

\all{Sie hat sich um ein Stipendium beworben.} « Elle a demandé une bourse. »

Piège : ne traduisez jamais « postuler pour » par \all{*sich bewerben für}. La préposition allemande est \all{um}, pas \all{für}.
\end{propriete}

\begin{exemplebox}[La candidature en phrases]
\all{Ich habe Ihre Stellenanzeige in der Zeitung gelesen.} « J'ai lu votre offre d'emploi dans le journal. »

\all{Wir laden Sie zu einem Vorstellungsgespräch ein.} « Nous vous invitons à un entretien d'embauche. »

\all{In der Anlage finden Sie meinen Lebenslauf und meine Zeugnisse.} « Vous trouverez en pièce jointe mon CV et mes diplômes. »
\end{exemplebox}

\courssub{Le champ lexical de la réclamation : die Beschwerde}

\begin{definition}[die Beschwerde]
\all{Die Beschwerde} (pluriel : \all{die Beschwerden}) désigne la \textbf{plainte}, la \textbf{réclamation} : le fait de signaler un problème (produit défectueux, service insuffisant, erreur de facture). Le verbe est \all{sich beschweren über + Akkusativ} : « se plaindre de ». Dans le commerce, on utilise aussi le mot plus technique \all{die Reklamation, -en} (réclamation adressée à un vendeur).
\end{definition}

\begin{center}
\begin{tabularx}{\linewidth}{|X|X|X|}
\hline
\rowcolor{poporangeL} \textbf{Deutsch (avec pluriel)} & \textbf{Français} & \textbf{Exemple d'emploi} \\
\hline
\all{die Beschwerde, -n} & la plainte, réclamation & \all{eine Beschwerde einreichen} \\
\hline
\all{die Reklamation, -en} & la réclamation (commerce) & \all{Ihre Reklamation vom 3. Mai} \\
\hline
\all{der Fehler, -} & l'erreur, le défaut & \all{ein Fehler in der Rechnung} \\
\hline
\all{der Ersatz (Sg.)} & le remplacement & \all{Ich bitte um Ersatz.} \\
\hline
\all{die Rückerstattung, -en} & le remboursement & \all{eine Rückerstattung verlangen} \\
\hline
\all{die Bestellung, -en} & la commande & \all{meine Bestellung vom Montag} \\
\hline
\all{defekt (Adj.)} & défectueux, en panne & \all{Das Gerät ist defekt.} \\
\hline
\all{die Lieferung, -en} & la livraison & \all{Die Lieferung kam zu spät.} \\
\hline
\all{die Qualität, -en} & la qualité & \all{schlechte Qualität} \\
\hline
\end{tabularx}
\end{center}

\begin{propriete}[« sich beschweren über » + accusatif]
Le verbe réfléchi \all{sich beschweren} se construit avec la préposition \all{über} suivie de l'\textbf{accusatif} : elle introduit \textbf{l'objet de la plainte}. Pour indiquer la personne à qui l'on se plaint, on utilise \all{bei + Dativ}.

\all{Er beschwert sich über die schlechte Qualität.} « Il se plaint de la mauvaise qualité. »

\all{Die Kundin beschwert sich beim Chef über den Fehler.} « La cliente se plaint auprès du chef de l'erreur. »

Piège : ne dites pas \all{*sich beschweren gegen} ni \all{*sich beschweren auf}. C'est toujours \all{über + Akkusativ}.
\end{propriete}

\courssub{Les formules de la lettre formelle}

Ces formules figent l'ouverture et la clôture de toute lettre officielle (candidature ou réclamation). Elles doivent être apprises \textbf{par cœur}, avec leur orthographe exacte.

\begin{center}
\begin{tabularx}{\linewidth}{|X|X|X|}
\hline
\rowcolor{poppurpleL} \textbf{Deutsch} & \textbf{Français} & \textbf{Emploi} \\
\hline
\all{Sehr geehrte Damen und Herren,} & Madame, Monsieur, & destinataire inconnu \\
\hline
\all{Sehr geehrter Herr Müller,} & Monsieur, & destinataire masculin connu \\
\hline
\all{Sehr geehrte Frau Ngono,} & Madame, & destinataire féminin connu \\
\hline
\all{Bezug nehmend auf Ihr Schreiben vom 3. Mai ...} & Me référant à votre lettre du 3 mai ... & ouverture de réponse \\
\hline
\all{Ich hätte gern nähere Informationen über ...} & J'aimerais avoir de plus amples informations sur ... & demande polie (Konjunktiv II) \\
\hline
\all{Ich bitte um Ersatz / um Rückerstattung.} & Je vous prie de procéder au remplacement / remboursement. & réclamation \\
\hline
\all{Mit freundlichen Grüßen} & Je vous prie d'agréer mes salutations distinguées & formule de politesse finale \\
\hline
\end{tabularx}
\end{center}

\begin{attention}[Pièges des formules de politesse]
\begin{itemize}
\item \all{Sehr geehrter Herr Müller} : \all{geehrter} prend la terminaison \textbf{-er} au masculin, \textbf{-e} au féminin (\all{geehrte Frau}), \textbf{-e} au pluriel (\all{geehrte Damen und Herren}). C'est un adjectif décliné (déclinaison forte) !
\item \all{Mit freundlichen Grüßen} s'écrit \textbf{sans virgule} à la fin, contrairement au français « Cordialement, ».
\item Après la formule d'appel (\all{Sehr geehrte Damen und Herren,}), la phrase suivante commence par une \textbf{minuscule} en allemand : \all{Sehr geehrte Damen und Herren,}\\\all{ich habe Ihre Anzeige gelesen ...}
\item Faux ami : \all{die Anlage} ne signifie pas « installation » ici, mais « pièce jointe » d'une lettre.
\end{itemize}
\end{attention}

\courssub{Familles de mots : un verbe, plusieurs noms}

L'allemand construit facilement des noms à partir de verbes. Apprendre ces familles permet de multiplier son vocabulaire sans effort.

\begin{center}
\begin{tabularx}{\linewidth}{|X|X|X|}
\hline
\rowcolor{popgoldL} \textbf{Verbe} & \textbf{Noms dérivés} & \textbf{Français} \\
\hline
\all{sich bewerben} & \all{die Bewerbung, der Bewerber, das Bewerbungsschreiben} & postuler → candidature, candidat, lettre de motivation \\
\hline
\all{kennen} & \all{die Kenntnis, die Kenntnisse} & connaître → connaissance, compétences \\
\hline
\all{bestellen} & \all{die Bestellung, der Besteller} & commander → commande, client \\
\hline
\all{sich beschweren} & \all{die Beschwerde} & se plaindre → plainte \\
\hline
\all{liefern} & \all{die Lieferung, der Lieferant} & livrer → livraison, fournisseur \\
\hline
\all{erstatten} & \all{die Rückerstattung} & rembourser → remboursement \\
\hline
\end{tabularx}
\end{center}

\begin{popfigure}
\begin{center}
\begin{tikzpicture}[
  root/.style={rectangle, rounded corners, draw=popblue, fill=popblue!20, font=\bfseries, align=center, minimum height=0.9cm},
  child node/.style={rectangle, rounded corners, draw=poporange, fill=poporange!20, align=center, minimum height=0.8cm},
  link/.style={-{Stealth[length=2.5mm]}, thick, popdark}]
\node[root, align=center] (v) at (0,0){\all{sich bewerben}\\ (postuler)};
\node[child node, align=center] (a) at (-4.6,-2){\all{die Bewerbung, -en}\\ la candidature};
\node[child node, align=center] (b) at (0,-2){\all{der Bewerber, -}\\ le candidat};
\node[child node, align=center] (c) at (4.6,-2){\all{das Bewerbungsschreiben, -}\\ la lettre de motivation};
\draw[link] (v) -- (a);
\draw[link] (v) -- (b);
\draw[link] (v) -- (c);
\end{tikzpicture}
\end{center}
\legende{Arbre de la famille de mots : du verbe \all{sich bewerben} aux noms dérivés.}
\end{popfigure}

\begin{popfigure}
\begin{center}
\begin{tikzpicture}[
  box/.style={rectangle, rounded corners, draw=popdark, thick, align=center, minimum width=3.6cm, minimum height=0.75cm, font=\small}]
\node[box, fill=popblue!20, draw=popblue] (l) at (0,0) {\textbf{\all{Lebenslauf}}};
\node[box, fill=popblue!10, draw=popblue!60] at (0,-1.1) {\all{die Schulbildung}};
\node[box, fill=popblue!10, draw=popblue!60] at (0,-2) {\all{die Berufserfahrung}};
\node[box, fill=popblue!10, draw=popblue!60] at (0,-2.9) {\all{die Sprachkenntnisse}};
\node[box, fill=popblue!10, draw=popblue!60] at (0,-3.8) {\all{das Hobby}};
\node[box, fill=popgreen!20, draw=popgreen] (b) at (5.2,0) {\textbf{\all{Bewerbung}}};
\node[box, fill=popgreen!10, draw=popgreen!60] at (5.2,-1.1) {\all{die Stelle}};
\node[box, fill=popgreen!10, draw=popgreen!60] at (5.2,-2) {\all{der Bewerber}};
\node[box, fill=popgreen!10, draw=popgreen!60] at (5.2,-2.9) {\all{das Vorstellungsgespräch}};
\node[box, fill=popgreen!10, draw=popgreen!60] at (5.2,-3.8) {\all{die Anlage}};
\node[box, fill=poporange!20, draw=poporange] (r) at (10.4,0) {\textbf{\all{Beschwerde}}};
\node[box, fill=poporange!10, draw=poporange!60] at (10.4,-1.1) {\all{der Fehler}};
\node[box, fill=poporange!10, draw=poporange!60] at (10.4,-2) {\all{die Bestellung}};
\node[box, fill=poporange!10, draw=poporange!60] at (10.4,-2.9) {\all{der Ersatz}};
\node[box, fill=poporange!10, draw=poporange!60] at (10.4,-3.8) {\all{die Rückerstattung}};
\end{tikzpicture}
\end{center}
\legende{Carte mentale lexicale : les trois champs de la leçon avec leurs mots-clés.}
\end{popfigure}

\courssub{Texte d'application : un jeune diplômé camerounais}

\begin{textebox}[Zwei Briefe aus Douala]
\all{Arnaud Etoundi aus Douala hat sein Abitur gemacht und danach drei Jahre an der Universität studiert. Jetzt liest er eine Stellenanzeige: Eine deutsche Firma sucht einen Übersetzer. Arnaud möchte sich um diese Stelle bewerben. Er schreibt ein Bewerbungsschreiben und fügt seinen Lebenslauf bei. In seinem Lebenslauf stehen seine persönlichen Daten, seine Schulbildung und seine Sprachkenntnisse: Deutsch, Englisch und Französisch. Seine Hobbys sind Fußball und Lesen.}

\all{„Sehr geehrte Damen und Herren,“, beginnt er seinen Brief, „bezug nehmend auf Ihre Anzeige bewerbe ich mich um die Stelle als Übersetzer. Ich hätte gern mehr Informationen über die Arbeitszeit. Mit freundlichen Grüßen“.}

\all{Eine Woche später bekommt Arnaud ein Paket: Er hat bei einer deutschen Internetfirma einen Laptop bestellt. Aber das Gerät ist defekt! Er schreibt sofort eine Beschwerde: „Sehr geehrte Damen und Herren, am Montag habe ich meine Bestellung erhalten. Leider funktioniert der Laptop nicht. Ich bitte um Ersatz oder um Rückerstattung des Preises. Mit freundlichen Grüßen“.}

\all{Die Firma antwortet schnell: Sie entschuldigt sich für den Fehler und schickt einen neuen Laptop. Arnaud ist zufrieden — und zwei Tage später lädt ihn die deutsche Firma zu einem Vorstellungsgespräch ein. Vielleicht beginnt bald seine Berufserfahrung als Übersetzer!}
\end{textebox}

\textbf{Glossaire :} \all{die Stellenanzeige, -n} (l'offre d'emploi) ; \all{beifügen} (joindre) ; \all{die Arbeitszeit, -en} (l'horaire de travail) ; \all{das Gerät, -e} (l'appareil) ; \all{erhalten, erhielt, hat erhalten} (recevoir) ; \all{sich entschuldigen für + Akk.} (s'excuser de) ; \all{einladen zu + Dat.} (inviter à).

\textbf{Questions de vocabulaire :}
\begin{enumerate}
\item Quels mots du texte appartiennent au champ lexical de la candidature ?
\item Comment dit-on « remboursement » et « remplacement » en allemand ?
\item Quelle formule Arnaud utilise-t-il pour ouvrir et pour clore ses deux lettres ?
\end{enumerate}

\begin{exoresolu}[Complétez avec le mot qui convient]
Complétez les phrases avec : \all{Lebenslauf, Bewerbung, Stelle, Beschwerde, Ersatz, Kenntnisse, Vorstellungsgespräch}.

1. \all{Ich bewerbe mich um die \dots\ als Sekretärin.}
2. \all{In meinem \dots\ stehen meine Schulbildung und meine Hobbys.}
3. \all{Das Handy ist defekt. Ich schreibe eine \dots\ an die Firma.}
4. \all{Ich bitte um \dots\ des defekten Geräts.}
5. \all{Meine Deutsch-\dots\ sind sehr gut.}
6. \all{Morgen habe ich ein \dots\ bei der Firma in Bonabéri.}
\tcblower
\textbf{Solution détaillée :}

1. \all{Stelle} — on postule \emph{pour un poste} : \all{sich bewerben um die Stelle}.

2. \all{Lebenslauf} — la formation et les loisirs figurent dans le CV.

3. \all{Beschwerde} — un produit défectueux provoque une réclamation.

4. \all{Ersatz} — \all{um Ersatz bitten} : demander le remplacement (mot sans article pluriel, masculin singulier).

5. \all{Kenntnisse} — toujours au pluriel pour les compétences : \all{Deutschkenntnisse}.

6. \all{Vorstellungsgespräch} — l'entretien d'embauche (mot composé : \all{vorstellen} + \all{das Gespräch}).
\end{exoresolu}

\begin{exercice}[À vous de jouer]
1. Traduisez en allemand : « Je me plains de la mauvaise qualité de la livraison. » / « Elle pose sa candidature pour le poste de secrétaire. »

2. Donnez l'article et le pluriel de : \all{Stelle, Fehler, Hobby, Bewerbung, Bestellung, Abschluss}.

3. Construisez une phrase avec \all{sich bewerben um} et une avec \all{sich beschweren über} en utilisant le contexte camerounais (marché de Mokolo, lycée de Yaoundé).
\end{exercice}

\begin{corrige}[Exercice : À vous de jouer]
1. \all{Ich beschwere mich über die schlechte Qualität der Lieferung.} / \all{Sie bewirbt sich um die Stelle als Sekretärin.}

2. \all{die Stelle, -n ; der Fehler, - ; das Hobby, -s ; die Bewerbung, -en ; die Bestellung, -en ; der Abschluss, -schlüsse.}

3. Exemples : \all{Ich bewerbe mich um die Stelle als Lehrer am Gymnasium Yaoundé.} / \all{Die Kundin beschwert sich über die schlechte Qualität der Tomaten auf dem Markt in Mokolo.}
\end{corrige}

\begin{savaistu}[Le CV allemand : une tradition différente]
En Allemagne, en Autriche et en Suisse, le CV comporte traditionnellement une \textbf{photo} du candidat, la \textbf{date et le lieu de naissance}, et il est \textbf{daté et signé} en bas de la dernière page : \all{Yaoundé, den 15. Mai 2025 — Unterschrift}. Cette tradition s'éloigne du CV français moderne, qui tend à devenir anonyme (sans photo ni âge) pour lutter contre les discriminations. Autre particularité : les Allemands joignent souvent tous leurs \all{Zeugnisse} (bulletins et certificats) dans un dossier complet appelé \all{Bewerbungsmappe} (pluriel : \all{-n}). Au Cameroun, où les entreprises germanophones et les coopérations allemandes sont présentes, connaître ces codes est un vrai atout pour un jeune diplômé !
\end{savaistu}

\begin{aretenir}[L'essentiel du lexique]
\begin{itemize}
\item Trois champs : \all{der Lebenslauf} (le CV), \all{die Bewerbung} (la candidature), \all{die Beschwerde} (la réclamation).
\item Deux constructions à préposition fixe : \all{sich bewerben \textbf{um} + Akk.} et \all{sich beschweren \textbf{über} + Akk.}
\item \all{die Kenntnisse} : presque toujours au pluriel.
\item Formules à mémoriser : \all{Sehr geehrte Damen und Herren,} / \all{Mit freundlichen Grüßen} (sans virgule finale).
\item Chaque nom s'apprend avec article et pluriel : \all{die Stelle, -n ; der Fehler, - ; das Hobby, -s}.
\end{itemize}
\end{aretenir}

\coursec{Grammaire 1 : le subjonctif II de politesse}

Dans la section précédente, nous avons enrichi notre vocabulaire autour de \all{der Lebenslauf}, \all{die Bewerbung} et \all{die Beschwerde}. En relisant attentivement la lettre de réclamation de notre texte support, une remarque s'impose : l'auteur ne donne jamais d'ordre direct. Il n'écrit pas \all{Schicken Sie mir die Rechnung!} mais des tournures beaucoup plus douces. Cette douceur n'est pas un hasard : elle repose sur un temps grammatical précis, le \cle{subjonctif II} (\all{der Konjunktiv II}), que nous allons maintenant observer, former et employer.

\courssub{Observation : que remarque-t-on dans la lettre ?}

Reprenons trois phrases tirées de la lettre de réclamation étudiée :

\begin{exemplebox}[Les phrases clés du texte support]
\all{Ich hätte gern eine schriftliche Bestätigung.} « Je voudrais une confirmation écrite. »

\all{Ich würde Sie bitten, das Gerät zu ersetzen.} « Je vous prierais de remplacer l'appareil. »

\all{Könnten Sie mir bitte die Rechnung schicken?} « Pourriez-vous m'envoyer la facture, s'il vous plaît? »
\end{exemplebox}

Observons les formes verbales : \all{hätte}, \all{würde}, \all{könnten}. Aucune n'est un présent, aucune n'est un impératif. Toutes trois ressemblent à des imparfaits (\all{hatte}, \all{wurde}, \all{konnten}) mais portent un \cle{tréma} (Umlaut) sur la voyelle radicale. Et toutes trois expriment la même chose qu'en français un \emph{conditionnel} : je voudrais, je prierais, pourriez-vous.

\begin{aretenir}[L'idée centrale]
En allemand, pour adoucir une demande, un souhait ou un conseil, on n'utilise ni l'impératif ni le présent, mais le \cle{subjonctif II}. C'est l'équivalent fonctionnel du \emph{conditionnel de politesse} français : « je voudrais », « pourriez-vous », « vous devriez ».
\end{aretenir}

\courssub{Formation du subjonctif II}

\begin{propriete}[Règle de formation]
Le subjonctif II se forme sur le \cle{radical de l'imparfait} (Präteritum) du verbe :
\begin{itemize}
\item pour les \textbf{verbes forts} et les auxiliaires, la voyelle radicale \textbf{a, o, u} prend un \textbf{Umlaut} : \all{haben} $\rightarrow$ \all{hatte} $\rightarrow$ \all{hätte} ; \all{können} $\rightarrow$ \all{konnte} $\rightarrow$ \all{könnte} ; \all{dürfen} $\rightarrow$ \all{durfte} $\rightarrow$ \all{dürfte} ; \all{sein} $\rightarrow$ \all{war} $\rightarrow$ \all{wäre} ; \all{mögen} $\rightarrow$ \all{mochte} $\rightarrow$ \all{möchte} ;
\item pour les \textbf{verbes faibles}, le subjonctif II serait identique à l'imparfait (\all{machte}, \all{sagte}) : on le remplace donc presque toujours par \all{würde} + infinitif : \all{ich würde machen}, \all{ich würde sagen}.
\end{itemize}
Les terminaisons sont celles de l'imparfait faible : \all{-e, -est, -e, -en, -et, -en}.
\end{propriete}

\begin{center}
\begin{tabular}{|l|l|l|l|}
\hline
\rowcolor{popblueL} Infinitif & Präteritum & Konjunktiv II & Français \\
\hline
sein & war & wäre & serait \\
\hline
haben & hatte & hätte & aurait \\
\hline
werden & wurde & würde & deviendrait / -rait (auxiliaire) \\
\hline
können & konnte & könnte & pourrait \\
\hline
müssen & musste & müsste & devrait (obligation) \\
\hline
dürfen & durfte & dürfte & pourrait (permission) \\
\hline
sollen & sollte & sollte & devrait (conseil) \\
\hline
mögen & mochte & möchte & voudrait \\
\hline
\end{tabular}
\end{center}

Remarquez que \all{sollen} $\rightarrow$ \all{sollte} ne prend pas d'Umlaut : sa voyelle radicale \emph{o} est brève et suivie d'une géminée (\emph{ll}), ce qui bloque la modification vocalique ; le Konjunktiv II est donc identique au Präteritum. Le sens (conseil) suffit à le distinguer. Le verbe \all{mögen} (aimer) est un cas particulier essentiel : son Präteritum \all{mochte} prend l'Umlaut \rightarrow \all{möchte} (« je voudrais »), forme très fréquente à l'oral et à l'écrit.

\begin{popfigure}
\begin{center}
\begin{tikzpicture}[font=\small]
% sein
\draw[thick, popblue] (0,0) rectangle (4.6,5.6);
\node[fill=popblue, text=white, font=\bfseries, minimum width=4.6cm] at (2.3,5.3) {sein : wäre};
\node at (2.3,4.7) {ich wäre};
\node at (2.3,4.2) {du wärest};
\node at (2.3,3.7) {er/sie/es wäre};
\node at (2.3,3.2) {wir wären};
\node at (2.3,2.7) {ihr wäret};
\node at (2.3,2.2) {sie/Sie wären};
\node[text=popmuted] at (2.3,1.5) {serait, serions... (aussi: du wärst, ihr wärt)};

% haben
\draw[thick, poporange] (5.0,0) rectangle (9.6,5.6);
\node[fill=poporange, text=white, font=\bfseries, minimum width=4.6cm] at (7.3,5.3) {haben : hätte};
\node at (7.3,4.7) {ich hätte};
\node at (7.3,4.2) {du hättest};
\node at (7.3,3.7) {er/sie/es hätte};
\node at (7.3,3.2) {wir hätten};
\node at (7.3,2.7) {ihr hättet};
\node at (7.3,2.2) {sie/Sie hätten};
\node[text=popmuted] at (7.3,1.5) {aurait, aurions...};

% werden
\draw[thick, popgreen] (0,-6.2) rectangle (4.6,-0.6);
\node[fill=popgreen, text=white, font=\bfseries, minimum width=4.6cm] at (2.3,-0.9) {werden : würde + Inf.};
\node at (2.3,-1.5) {ich würde};
\node at (2.3,-2.0) {du würdest};
\node at (2.3,-2.5) {er/sie/es würde};
\node at (2.3,-3.0) {wir würden};
\node at (2.3,-3.5) {ihr würdet};
\node at (2.3,-4.0) {sie/Sie würden};
\node[text=popmuted] at (2.3,-4.7) {-rait, -rions... (auxiliaire)};

% können
\draw[thick, poppurple] (5.0,-6.2) rectangle (9.6,-0.6);
\node[fill=poppurple, text=white, font=\bfseries, minimum width=4.6cm] at (7.3,-0.9) {können : könnte};
\node at (7.3,-1.5) {ich könnte};
\node at (7.3,-2.0) {du könntest};
\node at (7.3,-2.5) {er/sie/es könnte};
\node at (7.3,-3.0) {wir könnten};
\node at (7.3,-3.5) {ihr könntet};
\node at (7.3,-4.0) {sie/Sie könnten};
\node[text=popmuted] at (7.3,-4.7) {pourrait, pourrions...};
\end{tikzpicture}
\end{center}
\legende{Tableau de conjugaison au subjonctif II : les formes indispensables de la lettre formelle. Note : les formes \emph{du wärst / ihr wärt} sont des raccourcis courants de \emph{wärest / wäret}.}
\end{popfigure}

\courssub{Les quatre emplois de la politesse}

\textbf{1. La demande polie : \all{ich hätte gern} / \all{ich würde gern} / \all{ich möchte}.}\par
C'est la tournure reine de la correspondance et du commerce. \all{gern} marque le souhait, le subjonctif II l'adoucit. \all{ich möchte} (de \all{mögen}) est l'alternative la plus courante pour « je voudrais + verbe ».

\begin{exemplebox}[Demandes polies]
\all{Ich hätte gern einen Termin.} « Je voudrais un rendez-vous. »

\all{Ich hätte gern mehr Informationen über den Kurs.} « Je voudrais plus d'informations sur le cours. »

\all{Ich würde gern mit dem Direktor sprechen.} « Je voudrais parler au directeur. » (\all{würde gern} + infinitif en fin de phrase)

\all{Ich möchte das Produkt umtauschen.} « Je voudrais échanger le produit. » (\all{möchte} + infinitif en fin de phrase)
\end{exemplebox}

\textbf{2. La question polie : \all{Könnten Sie bitte...?} / \all{Würden Sie bitte...?}}\par
L'auxiliaire \all{können} ou \all{werden} au subjonctif II transforme un ordre en requête courtoise. L'infinitif se place en fin de phrase.

\begin{exemplebox}[Questions polies]
\all{Könnten Sie mir bitte die Rechnung schicken?} « Pourriez-vous m'envoyer la facture, s'il vous plaît? »

\all{Könnten Sie das bitte bis Freitag prüfen?} « Pourriez-vous vérifier cela d'ici vendredi, s'il vous plaît? »

\all{Würden Sie bitte das Formular ausfüllen?} « Voudriez-vous remplir le formulaire, s'il vous plaît? »
\end{exemplebox}

\textbf{3. Le conseil : \all{Sie sollten...} / \all{Sie müssten...}.}\par
Pour conseiller sans imposer : \all{sollte} (devrait, conseil moral) et \all{müsste} (il faudrait que, nécessité plus forte).

\begin{exemplebox}[Conseils]
\all{Sie sollten das Produkt umtauschen.} « Vous devriez échanger le produit. »

\all{Sie müssten den Brief per Einschreiben schicken.} « Vous devriez envoyer la lettre en recommandé. »
\end{exemplebox}

\textbf{4. Le souhait irréel / le regret : \all{Wenn ich doch...!}}\par
Le subjonctif II exprime aussi le regret ou le rêve inaccessible, souvent avec \all{doch} ou \all{nur}. On utilise ici le \emph{subjonctif II du passé} (\all{hätte/wäre} + Partizip II).

\begin{exemplebox}[Souhaits et regrets]
\all{Wenn ich doch mehr Zeit hätte!} « Si seulement j'avais plus de temps! » (présent irréel)

\all{Hätte ich nur früher reklamiert!} « Si seulement j'avais réclamé plus tôt! » (passé irréel : \all{hätte} + \all{reklamiert})
\end{exemplebox}

\courssub{Comparaison français / allemand}

\begin{center}
\begin{tabularx}{\linewidth}{|X|X|X|}
\hline
\rowcolor{popblueL} Français & Deutsch & Remarque \\
\hline
Je voudrais un rendez-vous. & Ich hätte gern einen Termin. & \all{hätte gern} + nom (accusatif) \\
\hline
Je voudrais vous parler. & Ich würde gern mit Ihnen sprechen. & \all{würde gern} + infinitif final \\
\hline
Je voudrais partir. & Ich möchte abreisen. & \all{möchte} + infinitif final \\
\hline
Pourriez-vous m'aider? & Könnten Sie mir helfen? & inversion : verbe en tête, infinitif final \\
\hline
Vous devriez écrire une lettre. & Sie sollten einen Brief schreiben. & conseil = \all{sollte} \\
\hline
Je serais content. & Ich wäre froh. & \all{sein} $\rightarrow$ \all{wäre} \\
\hline
\end{tabularx}
\end{center}

La correspondance est presque mécanique : là où le français met un \emph{conditionnel}, l'allemand met un \emph{subjonctif II}. C'est une excellente nouvelle pour le thème : il suffit de repérer le conditionnel français et de choisir la bonne forme allemande.

\courssub{L'échelle de la politesse}

\begin{popfigure}
\begin{center}
\begin{tikzpicture}[font=\small]
\draw[-{Stealth[length=3mm]}, very thick, popline] (0,0) -- (12.5,0);
\node[below, text=popmuted] at (0.3,0) {moins poli};
\node[below, text=popmuted] at (12.2,0) {plus poli};
\node[draw, fill=poppinkL, rounded corners, align=center, minimum width=3.4cm] (a) at (2,1.1) {\textbf{Impératif}\\ „Schicken Sie die Rechnung!“\\ (Envoyez la facture!)};
\node[draw, fill=popgoldL, rounded corners, align=center, minimum width=3.4cm] (b) at (6.2,1.1) {\textbf{Présent}\\ „Ich will eine Antwort.“\\ (Je veux une réponse.)};
\node[draw, fill=popgreenL, rounded corners, align=center, minimum width=3.6cm] (c) at (10.4,1.1) {\textbf{Konjunktiv II}\\ „Könnten Sie bitte...?“\\ (Pourriez-vous...?)};
\draw[-{Stealth}, thick, poporange] (a.east) -- (b.west);
\draw[-{Stealth}, thick, popgreen] (b.east) -- (c.west);
\end{tikzpicture}
\end{center}
\legende{L'échelle de la politesse : dans une lettre formelle, on se place toujours à droite.}
\end{popfigure}

\begin{attention}[Pièges fréquents des francophones]
\begin{itemize}
\item \textbf{L'Umlaut est obligatoire.} Écrire \all{ich habe gern} au lieu de \all{ich hätte gern} est une double faute : le présent \all{habe} ne forme pas la tournure de politesse, et sans Umlaut la forme n'existe pas. On mémorise : \all{hätte}, \all{wäre}, \all{würde}, \all{könnte}, \all{müsste}, \all{dürfte}, \all{sollte}, \all{möchte}.
\item \textbf{Pas d'impératif dans une lettre formelle.} \all{Geben Sie mir eine Antwort!} (« Donnez-moi une réponse! ») sonne comme un ordre militaire. Dites : \all{Ich würde Sie bitten, mir eine Antwort zu geben.} (infinitif \all{zu geben} en fin de proposition subordonnée).
\item \textbf{\all{Ich will} est impoli.} Le francophone traduit « je veux » par \all{ich will} et paraît brutal. Remplacez toujours par \all{ich hätte gern} (nom), \all{ich würde gern} (verbe) ou \all{ich möchte} (verbe).
\item \textbf{Place de l'infinitif.} Avec \all{würde}, \all{könnte}, \all{sollte}, \all{möchte}, l'infinitif va en \emph{fin} de phrase (ou de proposition) : \all{Könnten Sie mir bitte die Rechnung \textbf{schicken}?}
\item \textbf{Verbes à particule séparable.} À l'infinitif final, la particule se ressoude : \all{anrufen} $\rightarrow$ \all{anrufen}, \all{zurückrufen} $\rightarrow$ \all{zurückrufen}, \all{umtauschen} $\rightarrow$ \all{umtauschen}. Ex. : \all{Könnten Sie mich bitte zurückrufen?}
\end{itemize}
\end{attention}

\begin{methode}[Traduire « je voudrais » et « pourriez-vous » (Thème)]
\begin{enumerate}
\item Repérez dans la phrase française le \textbf{conditionnel de politesse} (je voudrais, pourriez-vous, vous devriez, j'aimerais).
\item Choisissez la forme allemande :
      \begin{itemize}
      \item « je voudrais + nom » $\rightarrow$ \all{ich hätte gern} + nom (accusatif) \emph{ou} \all{ich möchte} + nom.
      \item « je voudrais + verbe » $\rightarrow$ \all{ich würde gern} + infinitif final \emph{ou} \all{ich möchte} + infinitif final.
      \item « pourriez-vous » $\rightarrow$ \all{Könnten Sie} + infinitif final.
      \item « vous devriez » (conseil) $\rightarrow$ \all{Sie sollten} + infinitif final.
      \end{itemize}
\item Vérifiez l'Umlaut sur chaque forme (\all{hätte}, \all{würde}, \all{könnte}, \all{möchte}, \all{sollte}...).
\item Placez l'infinitif en fin de phrase, ajoutez \all{bitte} dans les questions, et respectez la casse (majuscules aux noms).
\end{enumerate}
\end{methode}

\begin{exoresolu}[Thème : une demande polie]
Énoncé : traduisez en allemand — « Pourriez-vous m'envoyer le formulaire, s'il vous plaît? Je voudrais aussi une copie du contrat. »
\tcblower
Solution détaillée : « pourriez-vous » est un conditionnel de politesse $\rightarrow$ \all{Könnten Sie} ; le verbe « envoyer » (\all{schicken}) va à l'infinitif en fin de phrase ; « m' » devient \all{mir} (datif). Première phrase : \all{Könnten Sie mir bitte das Formular schicken?} Pour « je voudrais + nom », on choisit \all{ich hätte gern} + nom à l'accusatif : « une copie du contrat » = \all{eine Kopie des Vertrags} (génitif après \all{Kopie}). Phrase complète : \all{Könnten Sie mir bitte das Formular schicken? Ich hätte auch gern eine Kopie des Vertrags.}
\end{exoresolu}

\begin{exercice}[À vous : cinq demandes polies]
Traduisez en allemand, en employant le subjonctif II :
\begin{enumerate}
\item Je voudrais un entretien avec le directeur.
\item Pourriez-vous me rappeler demain, s'il vous plaît?
\item Vous devriez envoyer la lettre en recommandé.
\item Je voudrais échanger ce produit.
\item Pourriez-vous m'envoyer une confirmation écrite?
\end{enumerate}
\end{exercice}

\begin{corrige}[Exercice « cinq demandes polies »]
\begin{enumerate}
\item \all{Ich hätte gern ein Gespräch mit dem Direktor.} (ou \all{Ich würde gern mit dem Direktor sprechen.} / \all{Ich möchte mit dem Direktor sprechen.})
\item \all{Könnten Sie mich bitte morgen zurückrufen?} (verbe à particule : \all{zurückrufen}, la particule reste soudée à l'infinitif en fin de phrase)
\item \all{Sie sollten den Brief per Einschreiben schicken.}
\item \all{Ich würde gern dieses Produkt umtauschen.} (ou \all{Ich möchte dieses Produkt umtauschen.})
\item \all{Könnten Sie mir eine schriftliche Bestätigung schicken?}
\end{enumerate}
\end{corrige}

\begin{aretenir}[Bilan de la section]
Le subjonctif II de politesse se forme sur l'imparfait avec Umlaut (\all{hätte}, \all{wäre}, \all{würde}, \all{könnte}, \all{müsste}, \all{dürfte}, \all{sollte}, \all{möchte}) et correspond au conditionnel français. Les tournures à maîtriser pour la lettre formelle sont : \all{ich hätte gern} + nom, \all{ich würde gern} + infinitif, \all{ich möchte} + infinitif, \all{Könnten Sie bitte...?} et \all{Sie sollten...}. Armés de ces formes, nous pouvons maintenant examiner dans la section suivante la structure complète et les formules fixes de la lettre formelle allemande.
\end{aretenir}

\coursec{Grammaire 2 : structure et formules de la lettre formelle}

Dans la section précédente, nous avons appris à formuler des demandes polies grâce au \cle{subjonctif II} : \all{Ich hätte gern...}, \all{Ich würde Sie bitten,...}. Ces formules trouvent leur place naturelle dans un genre d'écrit très codifié en allemand : la \cle{lettre formelle} (\all{der formelle Brief}). Contrairement au français, où la mise en page d'une lettre administrative relève souvent de l'habitude, l'allemand possède des règles précises, presque mathématiques, que les examinateurs du baccalauréat sanctionnent sévèrement. Cette section vous donne toutes les clés pour rédiger une lettre de réclamation ou une lettre de candidature irréprochable.

\courssub{Observation : lettre amicale ou lettre formelle ?}

Lisons d'abord les deux extraits suivants et comparons-les.

\begin{textebox}[Deux lettres, deux registres]
\textbf{Lettre A (extrait)}\\
Liebe Anna,\\
vielen Dank für deinen Brief! Ich freue mich schon sehr auf deinen Besuch in Yaoundé. Wir können zusammen zum Markt gehen und Fußball schauen.\\
Viele Grüße\\
dein Paul

\textbf{Lettre B (extrait)}\\
Sehr geehrte Damen und Herren,\\
bezugnehmend auf meine Bestellung vom 12. März muss ich Ihnen leider mitteilen, dass das Paket bis heute nicht angekommen ist. Ich würde Sie bitten, den Fall zu prüfen und mir eine neue Lieferung zu schicken.\\
Mit freundlichen Grüßen\\
Jean Mbarga
\end{textebox}

\textbf{Questions d'observation :}
\begin{enumerate}
\item Quelle lettre s'adresse à un ami ? Laquelle s'adresse à une entreprise ?
\item Compare les formules d'appel et de conclusion des deux lettres.
\item Observe la première lettre du corps après la virgule de l'appel : majuscule ou minuscule ?
\end{enumerate}

On constate immédiatement que la lettre A tutoie (\all{deinen, dein}) et utilise des formules chaleureuses, tandis que la lettre B vouvoie (\all{Ihnen}), emploie le subjonctif II de politesse et des formules figées. C'est toute la différence entre le \cle{registre familier} et le \cle{registre formel}.

\begin{center}
\begin{tabularx}{\linewidth}{|X|X|X|}
\hline
\rowcolor{popblueL} Élément & Lettre amicale & Lettre formelle \\
\hline
Appel & Liebe Anna, / Lieber Paul, & Sehr geehrte Frau Schmidt, / Sehr geehrte Damen und Herren, \\
\hline
Pronom personnel & du, dich, dein & Sie, Ihnen, Ihr (majuscule !) \\
\hline
Registre & familier, direct & poli, souvent au Konjunktiv II \\
\hline
Conclusion & Viele Grüße / Liebe Grüße / Dein... & Mit freundlichen Grüßen / Hochachtungsvoll \\
\hline
Objet (Betreff) & absent & obligatoire \\
\hline
\end{tabularx}
\end{center}

\courssub{La structure obligatoire de la lettre formelle allemande}

Une lettre formelle allemande (\all{der Geschäftsbrief}) comprend huit zones, toujours dans le même ordre et à la même place sur la page. Les connaître par cœur est indispensable : au baccalauréat, une lettre sans \all{Betreff} ou avec une date mal placée perd des points de présentation.

\begin{definition}[Les huit zones de la lettre formelle]
\begin{enumerate}
\item \all{der Absender} : l'expéditeur (nom, adresse), en haut à gauche.
\item \all{der Empfänger} : le destinataire (nom ou service, adresse), sous l'expéditeur.
\item \all{Ort und Datum} : le lieu et la date, à droite (ex. \all{Yaoundé, den 15. Mai 2024}).
\item \all{der Betreff} : l'objet de la lettre, en gras, sans le mot « Objet : ».
\item \all{die Anrede} : la formule d'appel.
\item \all{der Brieftext} : le corps, divisé en paragraphes (introduction, faits, demande).
\item \all{die Grußformel} : la formule de politesse finale.
\item \all{die Unterschrift} : la signature, suivie éventuellement des \all{Anlagen} (pièces jointes).
\end{enumerate}
\end{definition}

\begin{popfigure}
\begin{tikzpicture}[font=\small]
\draw[thick, popdark] (0,0) rectangle (10,13);
\node[anchor=north west, align=left] at (0.4,12.6) {Jean Mbarga\\ Quartier Nlongkak\\ B.P. 1234 Yaoundé};
\node[anchor=north west, align=left] at (0.4,10.4) {Firma TechnoWelt GmbH\\ Kundendienst\\ Hauptstraße 45\\ 10115 Berlin};
\node[anchor=north east, align=right] at (9.6,8.9) {Yaoundé, den 15. Mai 2024};
\node[anchor=north west, align=left] at (0.4,8.2) {\textbf{Betreff: Reklamation — Bestellung Nr. 4567}};
\node[anchor=north west, align=left] at (0.4,7.4) {Sehr geehrte Damen und Herren,};
\node[anchor=north west, align=left] at (0.4,6.6) {bezugnehmend auf meine Bestellung ...\\ (corps de la lettre en paragraphes)};
\node[anchor=north west, align=left] at (0.4,4.4) {Mit freundlichen Grüßen};
\node[anchor=north west, align=left] at (0.4,3.6) {Jean Mbarga};
\node[anchor=north west, align=left] at (0.4,2.8) {Anlagen: Kopie der Rechnung};
\foreach \n/\x/\y in {1/9.9/12.2, 2/9.9/10.0, 3/9.9/8.9, 4/9.9/8.2, 5/9.9/7.4, 6/9.9/6.0, 7/9.9/4.4, 8/9.9/3.2}
  \node[circle, fill=poporange, text=white, inner sep=2pt] at (\x,\y) {\n};
\end{tikzpicture}
\legende{Schéma d'une page de lettre formelle : 1 Absender, 2 Empfänger, 3 Ort und Datum, 4 Betreff, 5 Anrede, 6 Brieftext, 7 Grußformel, 8 Unterschrift und Anlagen.}
\end{popfigure}

\begin{aretenir}[La date allemande]
La date s'écrit avec l'article défini et l'ordinal suivi d'un point : \all{Yaoundé, den 15. Mai 2024} (« Yaoundé, le 15 mai 2024 »). On écrit \all{den} + jour ordinal (\all{15.} se lit « fünfzehnten »), puis le mois avec une majuscule, puis l'année. Jamais de virgule entre le mois et l'année.
\end{aretenir}

\courssub{Les formules d'appel (die Anrede)}

Le choix de l'appel dépend d'une seule question : connais-tu le nom du destinataire ?

\begin{propriete}[Les trois formules d'appel formelles]
\begin{itemize}
\item Destinataire inconnu : \all{Sehr geehrte Damen und Herren,} (« Mesdames, Messieurs, »).
\item Destinataire masculin connu : \all{Sehr geehrter Herr Schmidt,} (« Monsieur Schmidt, »).
\item Destinataire féminin connu : \all{Sehr geehrte Frau Mbarga,} (« Madame Mbarga, »).
\end{itemize}
Attention à la déclinaison : \all{geehrter} au masculin (nominatif fort), \all{geehrte} au féminin et au pluriel.
\end{propriete}

\begin{propriete}[La règle typographique de la minuscule]
Après la formule d'appel, on met une \cle{virgule}, et la première phrase du corps commence obligatoirement par une \cle{minuscule} :\\
\all{Sehr geehrte Damen und Herren,}\\
\all{bezugnehmend auf Ihre Anzeige bewerbe ich mich um die Stelle.}\\
Cette règle, contraire au réflexe français (« Mesdames, Messieurs, » suivi d'une majuscule), est très souvent testée à l'examen. Seule exception : si le premier mot est un nom commun, il garde naturellement sa majuscule (\all{... die Ware...}).
\end{propriete}

\begin{attention}[« Lieber Herr » : le faux ami du registre]
Ne traduis jamais « Cher Monsieur » par \all{Lieber Herr} dans une lettre formelle ! \all{Lieber/Liebe} est réservé aux proches, aux amis et à la famille (\all{Lieber Paul, Liebe Anna}). Dans un contexte administratif ou commercial, seul \all{Sehr geehrte(r)...} est correct. Écrire \all{Lieber Herr Direktor} à un inconnu serait aussi déplacé que de le tutoyer.
\end{attention}

\courssub{Le corps de la lettre : les formules utiles}

Le corps se construit en trois temps : le motif (pourquoi j'écris), les faits (ce qui s'est passé), la demande (ce que j'attends). Voici les formules canoniques, à mémoriser comme des blocs.

\begin{exemplebox}[Formules du corps, traduites]
\all{Bezugnehmend auf Ihr Schreiben vom 3. April...} — « Me référant à votre courrier du 3 avril... »\\
\all{Bezugnehmend auf Ihre Anzeige bewerbe ich mich um die Stelle.} — « Me référant à votre annonce, je pose ma candidature pour le poste. »\\
\all{Ich hätte gern weitere Informationen über das Angebot.} — « J'aimerais avoir de plus amples informations sur l'offre. »\\
\all{Ich würde Sie bitten, mir die Ware zu ersetzen.} — « Je vous prierais de me remplacer la marchandise. »\\
\all{Leider muss ich Ihnen mitteilen, dass das Gerät defekt ist.} — « Je dois malheureusement vous informer que l'appareil est défectueux. »\\
\all{Ich freue mich auf Ihre Antwort.} — « Dans l'attente de votre réponse. » (litt. « Je me réjouis de votre réponse »)
\end{exemplebox}

Remarque la syntaxe de \all{Ich würde Sie bitten, ... zu + infinitif} : la proposition infinitive se place en fin de phrase, exactement comme nous l'avons vu dans la section sur le subjonctif II de politesse.

\courssub{Les formules de conclusion (die Grußformel)}

\begin{propriete}[Les deux formules de conclusion formelles]
\begin{itemize}
\item \all{Mit freundlichen Grüßen} (abréviation : \all{MfG}) : la formule standard, équivalente de « Cordialement » ou de « Veuillez agréer, Madame, Monsieur, mes salutations distinguées ».
\item \all{Hochachtungsvoll} : très solennel, réservé aux autorités (ministre, ambassadeur, recteur) : « Avec ma haute considération ».
\end{itemize}
La formule de conclusion n'est \textbf{jamais suivie d'une virgule} en allemand. La signature (\all{die Unterschrift}) suit à la ligne.
\end{propriete}

\begin{attention}[L'orthographe de « Mit freundlichen Grüßen »]
Deux fautes reviennent sans cesse dans les copies des francophones :
\begin{itemize}
\item \all{Mit freundliche Grüße} : FAUX. Après la préposition \all{mit} (datif pluriel), l'adjectif prend la désinence \all{-en} : \all{freundlichen Grüßen}. Le mot \all{der Gruß} prend un tréma au pluriel : \all{die Grüße}.
\item La virgule finale : on n'écrit jamais « Mit freundlichen Grüßen, ». Aucune ponctuation après la formule.
\end{itemize}
\end{attention}

\begin{center}
\begin{tabularx}{\linewidth}{|X|X|X|}
\hline
\rowcolor{popblueL} Français & Deutsch & Remarque \\
\hline
Monsieur, Madame (inconnu) & Sehr geehrte Damen und Herren, & toujours au pluriel \\
\hline
Monsieur Dupont, & Sehr geehrter Herr Dupont, & geehrter (masculin) \\
\hline
Madame Dupont, & Sehr geehrte Frau Dupont, & geehrte (féminin) \\
\hline
Objet : réclamation & Betreff: Reklamation & pas de mot « Objet » seul ; le Betreff est souvent en gras \\
\hline
Veuillez agréer... mes salutations distinguées & Mit freundlichen Grüßen & sans virgule finale \\
\hline
Je vous prie de bien vouloir... & Ich würde Sie bitten, ... & Konjunktiv II de politesse \\
\hline
Dans l'attente de votre réponse & Ich freue mich auf Ihre Antwort. & formule idiomatique \\
\hline
Pièces jointes & Anlagen & en bas, après la signature \\
\hline
\end{tabularx}
\end{center}

\begin{savaistu}[Pourquoi tant de rigueur ?]
Dans les pays germanophones, la lettre formelle est normée : la norme DIN 5008 fixe au millimètre près la position de chaque zone sur la page. Cette culture de la \all{Form} (la forme) reflète une valeur sociale forte : respecter la forme, c'est respecter son interlocuteur. Un CV ou une candidature mal présentés peuvent être écartés sans même être lus !
\end{savaistu}

\courssub{Texte d'application : une lettre de réclamation complète}

\begin{textebox}[Eine Reklamation aus Douala]
\textbf{Absender:} Amina Bello, Rue Joffre, Akwa, Douala\\
\textbf{Empfänger:} Firma ElectroMarkt, Kundendienst, Musterstraße 12, 80331 München\\
\textbf{Ort und Datum:} Douala, den 20. April 2024\\
\textbf{Betreff: Reklamation — defekter Laptop, Bestellung Nr. 8892}\\
Sehr geehrte Damen und Herren,\\
am 5. April habe ich in Ihrem Onlineshop einen Laptop bestellt. Das Paket ist am 15. April angekommen, aber leider funktioniert das Gerät nicht: Der Bildschirm bleibt schwarz, und der Akku lädt nicht.\\
Ich habe den Laptop sofort geprüft und alle Anweisungen in der Bedienungsanleitung befolgt, aber ohne Erfolg. Da das Gerät noch unter Garantie steht, würde ich Sie bitten, mir entweder ein neues Gerät zu schicken oder mir den Kaufpreis zurückzuerstatten.\\
Ich hätte auch gern Informationen über das Rücksendeverfahren. Eine Kopie der Rechnung und des Kaufvertrags liegt diesem Brief bei.\\
Ich freue mich auf Ihre schnelle Antwort.\\
Mit freundlichen Grüßen\\
Amina Bello\\
\textbf{Anlagen:} Kopie der Rechnung, Kopie des Kaufvertrags
\end{textebox}

\textbf{Glossaire :} \all{die Reklamation, -en} (réclamation) ; \all{der Kundendienst} (service clientèle) ; \all{defekt} (défectueux) ; \all{der Bildschirm, -e} (écran) ; \all{der Akku, -s} (batterie) ; \all{die Bedienungsanleitung, -en} (notice d'utilisation) ; \all{die Garantie, -n} (garantie) ; \all{zurückerstatten} (rembourser) ; \all{das Rücksendeverfahren, -} (procédure de retour) ; \all{die Rechnung, -en} (facture) ; \all{der Kaufvertrag, -e} (contrat d'achat) ; \all{beiliegen} (être joint).

\textbf{Questions de compréhension :}
\begin{enumerate}
\item Quel est le motif de la lettre d'Amina ?
\item Quelles sont les deux solutions qu'elle propose ?
\item Quelles pièces jointes accompagnent la lettre ?
\item Relevez dans le texte les huit zones de la lettre formelle et deux formules au subjonctif II.
\end{enumerate}

\courssub{Méthode : rédiger une lettre formelle en sept étapes}

\begin{methode}[Rédiger une lettre formelle]
\begin{enumerate}
\item \textbf{L'en-tête} : écris l'expéditeur (en haut à gauche), le destinataire, puis le lieu et la date à droite (\all{Yaoundé, den 15. Mai 2024}).
\item \textbf{Le Betreff} : résume l'objet en deux ou trois mots, en gras (\all{Betreff: Reklamation}).
\item \textbf{L'appel} : choisis la bonne formule selon que tu connais ou non le destinataire, et n'oublie pas la virgule.
\item \textbf{L'introduction} : annonce le motif en commençant par une minuscule (\all{bezugnehmend auf...}, \all{am ... habe ich...}).
\item \textbf{Le développement} : expose les faits dans l'ordre chronologique, au Perfekt ou au Präteritum.
\item \textbf{La demande} : formule-la poliment au subjonctif II (\all{Ich würde Sie bitten,...}, \all{Ich hätte gern...}).
\item \textbf{La formule finale} : \all{Ich freue mich auf Ihre Antwort.} puis \all{Mit freundlichen Grüßen} (sans virgule), ta signature et les \all{Anlagen} éventuelles.
\end{enumerate}
\end{methode}

\courssub{Exercices résolus}

\begin{exoresolu}[Remettre la lettre en ordre]
Les éléments suivants d'une lettre formelle sont mélangés. Remets-les dans le bon ordre :\\
a) Mit freundlichen Grüßen \quad b) Sehr geehrter Herr Ngono, \quad c) Betreff: Bewerbung um ein Praktikum \quad d) Yaoundé, den 3. Juni 2024 \quad e) Anlagen: Lebenslauf, Zeugnisse \quad f) Marie Etoundi, Quartier Bastos, Yaoundé \quad g) ich freue mich auf Ihre Antwort. \quad h) Firma Sahel GmbH, Personalabteilung, Bonn
\tcblower
\textbf{Solution :} l'ordre correct est f → h → d → c → b → g → a → e.\\
1) \all{Marie Etoundi, Quartier Bastos, Yaoundé} (Absender, en haut à gauche) ; 2) \all{Firma Sahel GmbH, Personalabteilung, Bonn} (Empfänger) ; 3) \all{Yaoundé, den 3. Juni 2024} (Ort und Datum, à droite) ; 4) \all{Betreff: Bewerbung um ein Praktikum} (l'objet) ; 5) \all{Sehr geehrter Herr Ngono,} (appel, destinataire masculin connu) ; 6) \all{ich freue mich auf Ihre Antwort.} (fin du corps, avec minuscule après la virgule de l'appel) ; 7) \all{Mit freundlichen Grüßen} (formule finale, sans virgule) ; 8) \all{Anlagen: Lebenslauf, Zeugnisse} (pièces jointes, en dernier).
\end{exoresolu}

\begin{exoresolu}[Version : traduire en français]
Traduis la lettre suivante :\\
\all{Sehr geehrte Damen und Herren,}\\
\all{bezugnehmend auf meine Bestellung vom 10. Februar muss ich Ihnen leider mitteilen, dass die Bücher noch nicht angekommen sind. Ich würde Sie bitten, die Sendung zu prüfen. Ich freue mich auf Ihre Antwort.}\\
\all{Mit freundlichen Grüßen}\\
\all{Paul Essomba}
\tcblower
\textbf{Traduction :} « Mesdames, Messieurs,\\
Me référant à ma commande du 10 février, je dois malheureusement vous informer que les livres ne sont pas encore arrivés. Je vous prierais de vérifier l'envoi. Dans l'attente de votre réponse.\\
Cordialement (ou : Veuillez agréer, Mesdames, Messieurs, mes salutations distinguées),\\
Paul Essomba. »\\
\textbf{Points de méthode :} \all{bezugnehmend auf} se rend par « me référant à » ; \all{die Sendung prüfen} = « vérifier l'envoi » ; on ne traduit jamais \all{Mit freundlichen Grüßen} mot à mot par « avec des salutations amicales », mais par une formule française équivalente.
\end{exoresolu}

\begin{exercice}[À toi de jouer]
1. Corrige les fautes dans cette lettre : \all{Lieber Herr Kamdem, Ich hätte gern Informationen über Ihre Kurse. Mit freundliche Grüße, Sandra}\\
2. Traduis en allemand : « Monsieur Atangana, me référant à votre lettre du 2 mai, je vous prierais de m'envoyer le formulaire. Dans l'attente de votre réponse, veuillez agréer mes salutations distinguées. »
\end{exercice}

\begin{corrige}[Exercice « À toi de jouer »]
1. Trois fautes : \all{Lieber Herr Kamdem} → \all{Sehr geehrter Herr Kamdem,} (registre formel) ; après la virgule de l'appel, minuscule obligatoire : \all{ich hätte gern...} ; \all{Mit freundliche Grüße} → \all{Mit freundlichen Grüßen} (datif pluriel, sans virgule finale).\\
2. \all{Sehr geehrter Herr Atangana, bezugnehmend auf Ihr Schreiben vom 2. Mai würde ich Sie bitten, mir das Formular zu schicken. Ich freue mich auf Ihre Antwort. Mit freundlichen Grüßen}
\end{corrige}

\begin{aretenir}[L'essentiel de la section]
Une lettre formelle allemande suit huit zones fixes : \all{Absender, Empfänger, Ort und Datum, Betreff, Anrede, Brieftext, Grußformel, Unterschrift/Anlagen}. Trois réflexes à garder : minuscule après la virgule de l'appel ; \all{Sehr geehrte(r)} et jamais \all{Lieber} pour un inconnu ; \all{Mit freundlichen Grüßen} sans virgule et avec \all{-en} aux deux mots. Ces règles de forme, combinées au subjonctif II de politesse étudié précédemment, te permettent de rédiger n'importe quelle lettre de réclamation ou de candidature exigée au baccalauréat.
\end{aretenir}

\coursec{Production guidée et méthode du commentaire}

Après avoir étudié la structure et les formules de la lettre formelle, nous passons maintenant à la pratique : comment exploiter un texte de ce type à l'oral et à l'écrit, notamment dans l'épreuve de commentaire du Baccalauréat ? Cette dernière section vous donne une méthode rigoureuse, un commentaire modèle rédigé pas à pas, puis des productions guidées (lettre à trous, mini-\all{Lebenslauf}, jeu de rôle) pour vérifier que vous savez réinvestir le vocabulaire, le \all{Konjunktiv II} de politesse et les formules épistolaires.

\courssub{La méthode du commentaire de texte au Bac}

Au Baccalauréat, le commentaire (\all{die Textanalyse}) porte souvent sur un document authentique : lettre, article, extrait de presse. Il ne s'agit pas de résumer, mais de \emph{présenter}, \emph{analyser} et \emph{juger} le texte de manière organisée.

\begin{methode}[Commenter un texte en quatre étapes]
\begin{enumerate}
\item \textbf{Introduction} : présenter le document (type de texte, auteur, source, date, thème général). Exemple de phrase d'accroche : \all{Der vorliegende Text ist ein Brief aus dem Jahr 2024. Er handelt von einer Beschwerde über ein defektes Gerät.} « Le texte présenté est une lettre de 2024. Elle traite d'une réclamation concernant un appareil défectueux. »
\item \textbf{Annonce du plan} : annoncer clairement les deux ou trois axes d'analyse. \all{Wir analysieren zuerst die Form des Briefes, dann seinen Inhalt und seine Sprache.} « Nous analysons d'abord la forme de la lettre, puis son contenu et sa langue. »
\item \textbf{Analyse (forme + contenu + langue)} : étudier la structure (en-tête, formules, paragraphes), le contenu (problème, arguments, demande) et la langue (registre, subjonctif II de politesse, vocabulaire thématique). Chaque remarque doit être justifiée par une citation allemande traduite.
\item \textbf{Conclusion personnelle} : bilan bref et avis justifié. \all{Der Brief ist höflich und klar strukturiert; die Beschwerde ist überzeugend.} « La lettre est polie et clairement structurée ; la réclamation est convaincante. »
\end{enumerate}
\end{methode}

\begin{popfigure}
\begin{center}
\begin{tikzpicture}[
node distance=0.55cm,
etape/.style={rectangle, rounded corners=4pt, draw=popblue, fill=popblueL, text width=5.6cm, align=center, font=\small, inner sep=6pt},
fleche/.style={-{Stealth[length=2.5mm]}, thick, poporange}
]
\node[etape, align=center] (intro){\textbf{1. Introduction}\\ type de texte, source, thème};
\node[etape, below=of intro, align=center] (plan){\textbf{2. Annonce du plan}\\ les axes d'analyse};
\node[etape, below=of plan, align=center] (analyse){\textbf{3. Analyse}\\ forme + contenu + langue\\ (citations traduites)};
\node[etape, below=of analyse, align=center] (concl){\textbf{4. Conclusion}\\ bilan + avis personnel};
\draw[fleche] (intro) -- (plan);
\draw[fleche] (plan) -- (analyse);
\draw[fleche] (analyse) -- (concl);
\end{tikzpicture}
\end{center}
\legende{Organigramme du commentaire de texte : de l'introduction à la conclusion}
\end{popfigure}

\begin{attention}[La règle d'or : toujours citer le texte]
Une remarque sans citation n'a aucune valeur au Bac. Pour chaque observation, citez le texte \textbf{en allemand} entre guillemets, puis \textbf{traduisez} la citation en français. Exemple de formulation correcte : « Le client exprime son mécontentement : \all{„Ich bin mit dem Gerät nicht zufrieden“} (je ne suis pas satisfait de l'appareil). » Évitez de paraphraser sans preuve textuelle. Attention aussi à la construction : \all{zufrieden sein mit} exige le datif (\all{mit dem Gerät}), jamais l'accusatif.
\end{attention}

\courssub{Commentaire modèle guidé : la lettre de réclamation du texte support}

Reprenons collectivement la lettre de réclamation étudiée dans la section 1 et construisons le commentaire en trois parties rédigées.

\textbf{Introduction (rédigée).}\par
\all{Der vorliegende Text ist ein formeller Brief. Ein Kunde schreibt an den Kundendienst einer Firma, weil sein neuer Computer defekt ist. Der Brief ist eine Beschwerde: Der Kunde ist nicht zufrieden und verlangt eine Lösung.}\par
« Le texte présenté est une lettre formelle. Un client écrit au service après-vente d'une entreprise parce que son nouvel ordinateur est défectueux. La lettre est une réclamation : le client n'est pas satisfait et exige une solution. »

\textbf{Analyse (rédigée).}\par
\begin{itemize}
\item \emph{La forme} : la lettre respecte la structure de la lettre formelle allemande : adresses, date, objet (\all{Betreff}), salutation \all{„Sehr geehrte Damen und Herren“} (Madame, Monsieur) et formule finale \all{„Mit freundlichen Grüßen“} (Cordialement).
\item \emph{Le contenu} : le client décrit le problème (\all{das Gerät ist defekt} — l'appareil est défectueux), donne des preuves (numéro de commande, facture) et formule une demande précise (réparation ou remboursement).
\item \emph{La langue} : le registre est poli et formel ; on remarque le \all{Konjunktiv II} de politesse : \all{„Ich wäre Ihnen dankbar, wenn Sie das Gerät reparieren könnten.“} (Je vous serais reconnaissant si vous pouviez réparer l'appareil.)
\end{itemize}

\textbf{Conclusion (rédigée).}\par
\all{Zusammenfassend kann man sagen: Der Brief ist höflich, präzise und gut strukturiert. Die Beschwerde hat große Chancen auf Erfolg.}\par
« En résumé, on peut dire : la lettre est polie, précise et bien structurée. La réclamation a de grandes chances de succès. »

\begin{exemplebox}[Phrases utiles pour le commentaire]
\all{Der Brief ist eine Beschwerde. Der Kunde ist nicht zufrieden, weil das Gerät defekt ist.} « La lettre est une réclamation. Le client n'est pas satisfait parce que l'appareil est défectueux. »\par
\all{Der Absender benutzt höfliche Formen wie „könnten Sie“ und „ich wäre Ihnen dankbar“.} « L'expéditeur utilise des formes polies comme „pourriez-vous“ et „je vous serais reconnaissant“. »\par
\all{Am Ende steht die Formel „Mit freundlichen Grüßen“.} « À la fin se trouve la formule „Cordialement“. »
\end{exemplebox}

\courssub{Production écrite guidée 1 : lettre de réclamation à trous}

\begin{exoresolu}[Compléter une lettre de réclamation]
Énoncé. Complétez la lettre avec les éléments suivants : \all{Sehr geehrte Damen und Herren} — \all{könnten} — \all{Mit freundlichen Grüßen} — \all{wäre} — \all{Beschwerde}.\par\medskip
\all{\_\_\_\_\_\_\_\_\_\_\_\_,}\par
\all{letzte Woche habe ich bei Ihnen ein Handy gekauft. Leider funktioniert der Akku nicht. Deshalb schreibe ich Ihnen diese \_\_\_\_\_\_\_\_\_\_. Ich \_\_\_\_\_\_ Ihnen dankbar, wenn Sie das Handy reparieren \_\_\_\_\_\_.}\par
\all{\_\_\_\_\_\_\_\_\_\_\_\_}\par
\all{Paul Essono}\par
\tcblower
Solution détaillée.\par
\all{\textbf{Sehr geehrte Damen und Herren},} (salutation obligatoire en tête d'une lettre formelle dont on ne connaît pas le destinataire ; la ligne suivante commence par une minuscule : \all{letzte Woche…})\par
\all{letzte Woche habe ich bei Ihnen ein Handy gekauft. Leider funktioniert der Akku nicht. Deshalb schreibe ich Ihnen diese \textbf{Beschwerde}. Ich \textbf{wäre} Ihnen dankbar, wenn Sie das Handy reparieren \textbf{könnten}.} (subjonctif II de politesse : \all{wäre} = je serais ; \all{könnten} = vous pourriez, verbe conjugué rejeté en fin de subordonnée introduite par \all{wenn})\par
\all{\textbf{Mit freundlichen Grüßen}} (formule de politesse finale, sans virgule après)\par
\all{Paul Essono}\par
Traduction : « Madame, Monsieur, la semaine dernière, j'ai acheté chez vous un téléphone portable. Malheureusement, la batterie ne fonctionne pas. C'est pourquoi je vous écris cette réclamation. Je vous serais reconnaissant si vous pouviez réparer le téléphone. Cordialement, Paul Essono. »
\end{exoresolu}

\courssub{Production écrite guidée 2 : rédiger son mini-Lebenslauf}

À vous maintenant de rédiger votre propre CV allemand en cinq rubriques et huit à dix lignes. Respectez les rubriques canoniques :

\begin{center}
\begin{tabularx}{\linewidth}{|X|X|X|}
\hline
\rowcolor{popblueL} \textbf{Rubrique} & \textbf{Contenu} & \textbf{Exemple allemand} \\
\hline
Persönliche Daten & nom, âge, adresse & \all{Manga Claire, 17 Jahre, Yaoundé} \\
\hline
Schulbildung & scolarité & \all{seit 2022: Lycée de Nkol-Eton, Yaoundé} \\
\hline
Sprachkenntnisse & langues & \all{Französisch (Muttersprache), Deutsch, Englisch} \\
\hline
Praktische Erfahrung & expériences & \all{2024: Praktikum in einem Büro in Douala} \\
\hline
Interessen & loisirs & \all{Fußball, Lesen, Musik} \\
\hline
\end{tabularx}
\end{center}

\begin{attention}[Les pièges du CV allemand]
\begin{itemize}
\item Utilisez l'\textbf{ordre chronologique inverse} : l'expérience la plus récente d'abord. C'est la norme en Allemagne, contrairement à l'habitude de certains CV francophones.
\item Préférez des \textbf{formulations nominales courtes}, sans phrase complète : écrivez \all{seit 2022: Lycée de Nkol-Eton} et non \all{Ich gehe seit 2022 ins Gymnasium}.
\item N'oubliez pas la majuscule à tous les noms : \all{die Schulbildung}, \all{die Sprachkenntnisse}, \all{die Interessen}.
\item Faux ami : \all{die Kenntnisse} signifie « les connaissances, les compétences », pas « les connaissances » au sens de « personnes que l'on connaît » (celles-ci se disent \all{die Bekannten}).
\item Attention à la préposition : \all{seit} (+ datif) marque une action commencée dans le passé et qui continue ; pour une période terminée, utilisez \all{von … bis …} : \all{von 2020 bis 2022: Collège de Biyem-Assi}.
\end{itemize}
\end{attention}

\begin{exercice}[Votre mini-Lebenslauf]
Rédigez votre CV allemand en 8 à 10 lignes avec les cinq rubriques du tableau : \all{Persönliche Daten}, \all{Schulbildung}, \all{Sprachkenntnisse}, \all{Praktische Erfahrung}, \all{Interessen}. Utilisez au moins deux formulations nominales et l'ordre chronologique inverse.
\end{exercice}

\begin{corrige}[Exercice : Votre mini-Lebenslauf]
Modèle de réponse possible (à adapter à votre situation) :\par
\all{\textbf{Lebenslauf}}\par
\all{\textbf{Persönliche Daten:} Ngo Bell Sandrine, 17 Jahre, Yaoundé, Kamerun}\par
\all{\textbf{Schulbildung:} seit 2023: Lycée de Nkol-Eton, Yaoundé (Klasse Tle A); von 2019 bis 2023: Collège d'Ekounou}\par
\all{\textbf{Sprachkenntnisse:} Französisch (Muttersprache), Deutsch (gute Kenntnisse), Englisch (Grundkenntnisse)}\par
\all{\textbf{Praktische Erfahrung:} 2024: Praktikum in einer Bank in Yaoundé}\par
\all{\textbf{Interessen:} Fußball, Lesen, Musik}\par
Traduction : « Curriculum vitae. Données personnelles : Ngo Bell Sandrine, 17 ans, Yaoundé, Cameroun. Scolarité : depuis 2023, lycée de Nkol-Eton, Yaoundé (classe de Tle A) ; de 2019 à 2023, collège d'Ekounou. Connaissances linguistiques : français (langue maternelle), allemand (bonnes connaissances), anglais (notions). Expérience pratique : 2024, stage dans une banque à Yaoundé. Centres d'intérêt : football, lecture, musique. »\par
Points de vigilance : rubriques avec majuscule, ordre chronologique inverse dans chaque rubrique, formulations nominales sans verbe conjugué, \all{seit} + datif pour la scolarité en cours.
\end{corrige}

\courssub{Entraînement au thème : phrases à traduire}

Le thème grammatical est une épreuve décisive au Bac. Entraînez-vous à traduire en allemand des phrases typiques de la correspondance formelle.

\begin{exoresolu}[Thème : la lettre de réclamation]
Énoncé. Traduisez en allemand : « 1. Je vous serais reconnaissant si vous pouviez réparer l'appareil. 2. Pourriez-vous m'envoyer une nouvelle facture ? 3. J'ai acheté ce ventilateur la semaine dernière, mais il ne fonctionne plus. »\par
\tcblower
Solution détaillée.\par
1. \all{Ich wäre Ihnen dankbar, wenn Sie das Gerät reparieren könnten.} — \all{wäre} (Konjunktiv II de \all{sein}) + datif \all{Ihnen} ; dans la subordonnée en \all{wenn}, le verbe conjugué \all{könnten} se place en fin de phrase, après l'infinitif \all{reparieren}.\par
2. \all{Könnten Sie mir eine neue Rechnung schicken?} — question polie au Konjunktiv II ; \all{mir} (datif de la personne) avant \all{eine neue Rechnung} (accusatif de la chose) ; déclinaison adjectivale : \all{eine neue Rechnung}.\par
3. \all{Letzte Woche habe ich diesen Ventilator gekauft, aber er funktioniert nicht mehr.} — le complément de temps en tête provoque l'inversion (\all{habe ich}) ; accusatif masculin : \all{diesen Ventilator} ; \all{nicht mehr} = « ne … plus ».
\end{exoresolu}

\courssub{Production orale : jeu de rôle au Kundendienst}

\begin{experience}[Client contre employé du service après-vente]
Par groupes de deux, jouez la scène suivante devant la classe (3 à 4 minutes).\par
\textbf{Élève A — le client} : vous avez acheté un article défectueux (ordinateur, chaussures, ventilateur). Vous téléphonez au \all{Kundendienst}. Décrivez le problème, donnez la date d'achat, demandez une réparation ou un remboursement avec le \all{Konjunktiv II} : \all{Ich hätte eine Frage…}, \all{Könnten Sie das Gerät umtauschen?}, \all{Ich wäre Ihnen sehr dankbar.}\par
\textbf{Élève B — l'employé(e)} : vous accueillez poliment le client : \all{Kundendienst, guten Tag!}, \all{Wie kann ich Ihnen helfen?}, \all{Haben Sie die Rechnung?}, \all{Wir können das Gerät reparieren oder umtauschen.}\par
Variante : inversez les rôles, puis un binôme joue la version « client difficile » (l'employé doit rester poli malgré tout).
\end{experience}

\begin{savaistu}[« MfG » : l'abréviation à manier avec prudence]
Dans les e-mails professionnels allemands, on rencontre très souvent l'abréviation \all{MfG} pour \all{Mit freundlichen Grüßen}. Elle est pratique et courante entre collègues, mais elle est jugée trop rapide dans une lettre très formelle — et elle est \textbf{à éviter au Bac} : l'examinateur attend la formule complète \all{Mit freundlichen Grüßen}, écrite en toutes lettres, sans virgule après. Réservez \all{MfG} à vos messages informels. À noter aussi : en allemand, contrairement au français, on ne met \emph{jamais} de virgule après la formule finale.
\end{savaistu}

\courssub{Grille d'auto-évaluation}

Après chaque production (lettre, CV, jeu de rôle), évaluez-vous honnêtement avec la grille suivante. Cochez chaque critère et notez-vous sur 2 points par ligne.

\begin{center}
\begin{tabularx}{\linewidth}{|X|c|c|}
\hline
\rowcolor{popblueL} \textbf{Critère} & \textbf{Oui} & \textbf{À revoir} \\
\hline
J'ai utilisé les formules correctes (\all{Sehr geehrte…}, \all{Mit freundlichen Grüßen}) & & \\
\hline
J'ai employé au moins deux \all{Konjunktiv II} de politesse (\all{könnten}, \all{wäre}, \all{hätte}) & & \\
\hline
J'ai placé le verbe conjugué en deuxième position et le participe ou l'infinitif en fin de phrase & & \\
\hline
J'ai utilisé le lexique thématique (\all{die Beschwerde}, \all{der Lebenslauf}, \all{die Kenntnisse}, \all{der Kundendienst}) & & \\
\hline
J'ai respecté la structure (en-tête, objet, paragraphes ; rubriques du CV en ordre chronologique inverse) & & \\
\hline
J'ai mis la majuscule à tous les noms allemands & & \\
\hline
\end{tabularx}
\end{center}

\begin{aretenir}[L'essentiel de la section]
Un bon commentaire de texte suit quatre étapes : introduction, annonce du plan, analyse (forme, contenu, langue) avec citations allemandes traduites, conclusion personnelle. Une bonne lettre de réclamation combine la structure formelle, le \all{Konjunktiv II} de politesse (\all{ich wäre Ihnen dankbar}, \all{könnten Sie…}) et le lexique précis de la réclamation. Un bon \all{Lebenslauf} allemand est bref, nominal et en ordre chronologique inverse. Entraînez-vous avec la grille d'auto-évaluation : c'est exactement ce que le correcteur du Bac regardera.
\end{aretenir}

\newpage

\coursec{Activité d'intégration}

\coursec{Leçon AL29 : Lettres de réclamation et curriculum vitae — forme et contenu}

\courssub{Texte support : Bewerbung und Beschwerde – schriftliche Kommunikation im Alltag}

\begin{textebox}[Texte original — Contexte camerounais et germanophone]
In Yaoundé und Douala, aber auch in München, Wien oder Zürich gehört die formelle schriftliche Kommunikation zu den wichtigsten Kompetenzen für Schule, Beruf und Alltag. Zwei Situationen treten besonders häufig auf: Man möchte sich um einen Ausbildungsplatz, ein Praktikum oder einen Job bewerben (\textbf{sich bewerben um}), oder man muss ein defektes Produkt oder eine schlechte Leistung reklamieren (\textbf{sich beschweren über}).

Für die \textbf{Bewerbung} ist der \textbf{Lebenslauf} (CV) das zentrale Dokument. Im deutschsprachigen Raum ist er tabellarisch (\textbf{tabellarischer Lebenslauf}), lückenlos, wahrheitsgetreu und endet mit Ort, Datum und eigenhändiger Unterschrift. Ein Foto ist üblich, aber in Deutschland durch das Allgemeine Gleichbehandlungsgesetz (AGG) nicht mehr verpflichtend. Wichtige Rubriken sind: \textit{Persönliche Daten}, \textit{Schulbildung}, \textit{Berufserfahrung / Praktika}, \textit{Sprachkenntnisse}, \textit{EDV-Kenntnisse} und \textit{Hobbys / Interessen}.

Die \textbf{Beschwerde} (der \textbf{Beschwerdebrief}) folgt einer strengen Struktur: Absender, Empfänger, Ort und Datum, Betrezzelle, förmliche Anrede, Sachverhaltsschilderung (meist im Perfekt oder Präteritum), konkrete Forderung (Umtausch, Reparatur, Rückerstattung) formuliert im \textbf{Konjunktiv II der Höflichkeit}, Schlussformel und Unterschrift. Höflichkeit und Sachlichkeit erhöhen die Erfolgsaussichten – Aggressivität schadet der eigenen Sache.

\emph{Traduction : À Yaoundé et Douala, mais aussi à Munich, Vienne ou Zurich, la communication écrite formelle est une compétence clé pour l'école, le travail et la vie quotidienne. Deux situations reviennent souvent : postuler pour une formation, un stage ou un emploi (\textbf{sich bewerben um}), ou réclamer pour un produit défectueux ou un mauvais service (\textbf{sich beschweren über}). Pour la \textbf{candidature}, le \textbf{CV} (Lebenslauf) est le document central. Dans l'espace germanophone, il est tabulaire (\textbf{tabellarischer Lebenslauf}), sans trous, véridique et se termine par le lieu, la date et la signature manuscrite. Une photo est courante, mais plus obligatoire en Allemagne (loi AGG). Rubriques principales : \textit{Données personnelles}, \textit{Scolarité}, \textit{Expérience professionnelle / Stages}, \textit{Compétences linguistiques}, \textit{Compétences informatiques}, \textit{Loisirs / Centres d'intérêt}. La \textbf{réclamation} suit une structure stricte : expéditeur, destinataire, lieu/date, objet, formule d'appel formelle, exposé des faits (souvent au Perfekt ou Präteritum), demande concrète (échange, réparation, remboursement) formulée au \textbf{subjonctif II de politesse}, formule de politesse finale et signature. Courtoisie et objectivité augmentent les chances de succès — l'agressivité nuit à sa propre cause.}
\end{textebox}

\courssub{Compréhension du texte}

\begin{exercice}[Questions de compréhension globale et détaillée]
Répondez en allemand par des phrases complètes.
\begin{enumerate}
\item Welche zwei häufigen Situationen für formelle Briefe nennt der Text?
\item Wie ist der Lebenslauf im deutschsprachigen Raum aufgebaut? (Stichworte: Form, Vollständigkeit, Ende)
\item Welche Rubriken gehören in einen tabellarischen Lebenslauf? Nennen Sie mindestens fünf.
\item Wie ist ein Beschwerdebrief strukturiert? Nennen Sie die Elemente in der richtigen Reihenfolge.
\item Warum soll man im Beschwerdebrief den Konjunktiv II benutzen? Welchen Effekt hat das?
\item Was bedeutet „tabellarisch“ im Kontext des Lebenslaufs? Warum ist ein Foto „nicht mehr verpflichtend“ in Deutschland?
\end{enumerate}
\end{exercice}

\begin{corrige}[Exercice 1 — Compréhension]
\begin{enumerate}
\item Der Text nennt die \textbf{Bewerbung} (sich bewerben um) und die \textbf{Beschwerde} (sich beschweren über).
\item Der Lebenslauf ist \textbf{tabellarisch}, \textbf{lückenlos}, \textbf{wahrheitsgetreu} und endet mit \textbf{Ort, Datum und Unterschrift}.
\item Persönliche Daten, Schulbildung, Berufserfahrung / Praktika, Sprachkenntnisse, EDV-Kenntnisse, Hobbys / Interessen.
\item Reihenfolge: Absender – Empfänger – Ort und Datum – Betreff – Anrede – Sachverhaltsschilderung – konkrete Forderung (im Konjunktiv II) – Grußformel – Unterschrift.
\item Der Konjunktiv II drückt \textbf{Höflichkeit} und \textbf{Distanz} aus. Er macht die Forderung zu einer \textbf{Bittenformulierung}, was den Ton sachlich und respektvoll hält und die Erfolgsaussichten erhöht.
\item „Tabellarisch“ bedeutet: Darstellung in einem \textbf{Tableau à deux colonnes} (Rubrik links, Inhalt rechts), nicht als Fließtext. Das Foto ist durch das AGG (Allgemeines Gleichbehandlungsgesetz) nicht mehr Pflicht, um Diskriminierung aufgrund von Aussehen, Herkunft etc. zu vermeiden.
\end{enumerate}
\end{corrige}

\courssub{Lexique thématique : Bewerbung, Lebenslauf, Beschwerde}

\begin{definition}[Vocabulaire de base — Noms, verbes, expressions]
\begin{center}
\begin{tabularx}{\linewidth}{|>{\bfseries}l|X|X|}
\hline
\rowcolor{popblueL} \textbf{Deutsch (Artikel, Plural)} & \textbf{Français} & \textbf{Famille de mots / Exemple} \\
\hline
die Bewerbung, -en & la candidature & sich bewerben um + Akk. (postuler) ; der Bewerber, die Bewerberin (le candidat) ; das Bewerbungsschreiben (la lettre de motivation) \\
\hline
der Lebenslauf, die Lebensläufe & le curriculum vitae (CV) & tabellarisch (tabulaire) ; lückenlos (sans trous) ; der tabellarische Lebenslauf \\
\hline
die Beschwerde, -n & la réclamation & sich beschweren über + Akk. (se plaindre de) ; der Beschwerdebrief (la lettre de réclamation) ; reklamieren (réclamer, technique) \\
\hline
die Kenntnis, -en (meist Pl.) & la connaissance, la compétence & die Sprachkenntnisse (compétences linguistiques) ; die EDV-Kenntnisse (compétences informatiques) ; gute Kenntnisse haben in + Dat. \\
\hline
die Anrede, -n & la formule d'appel & Sehr geehrte Damen und Herren, ; Sehr geehrte Frau Müller, ; Sehr geehrter Herr Meier, \\
\hline
die Grußformel, -n & la formule de politesse finale & Mit freundlichen Grüßen ; Mit freundlichen Grüßen aus Yaoundé \\
\hline
der Betreff, -s / Betreffe & l'objet (du courrier) & Betreff: Beschwerde über ... ; Betreff: Bewerbung um ein Praktikum \\
\hline
der Absender, - & l'expéditeur & Absender: Vorname Name, Adresse \\
\hline
der Empfänger, - & le destinataire & Empfänger: Firmenname, Abteilung, Adresse \\
\hline
die Rückerstattung, -en & le remboursement & die Rückerstattung verlangen / beantragen ; das Geld zurückerstatten \\
\hline
umtauschen (tauscht um, hat umgetauscht) & échanger & der Umtausch ; ich möchte die Ware umtauschen (trennbares Verb) \\
\hline
defekt & défectueux, en panne & ein defektes Gerät ; der rechte Ohrhörer ist defekt \\
\hline
die Bestellnummer, -n & le numéro de commande & Bestellnummer 45821 ; unter Angabe der Bestellnummer \\
\hline
die Teamfähigkeit & l'esprit d'équipe & Teamfähigkeit ist gefragt ; im Team arbeiten können \\
\hline
das Praktikum, -a & le stage & das Sommerpraktikum ; ein Praktikum machen / absolvieren ; der Praktikant, die Praktikantin \\
\hline
\end{tabularx}
\end{center}
\end{definition}

\courssub{Grammaire 1 : Le subjonctif II de politesse (Konjunktiv II der Höflichkeit)}

\begin{propriete}[Formation du Konjunktiv II — Deux voies principales]
Le subjonctif II (Konjunktiv II) exprime l'irréel, le souhait, la politesse. À l'écrit formel, on utilise deux constructions équivalentes :
\begin{enumerate}
\item \textbf{La forme « würde + infinitif »} (valable pour \textbf{tous} les verbes, y compris forts/irréguliers). C'est la forme la plus courante à l'oral et dans les lettres courantes.
\item \textbf{La forme « Präteritum + Umlaut »} (uniquement pour les verbes forts/irréguliers et les modaux, ainsi que \textit{haben} et \textit{sein}). Forme plus soutenue, très fréquente dans les lettres très formelles.
\end{enumerate}
\end{propriete}

\begin{exemplebox}[Tableau de formation — Verbes clés pour la politesse]
\begin{center}
\begin{tabular}{|l|c|c|c|}
\hline
\rowcolor{popblueL} \textbf{Infinitif} & \textbf{Präteritum (Indikatif)} & \textbf{Konjunktiv II (Forme Präteritum + Umlaut)} & \textbf{Konjunktiv II (Forme « würde »)} \\
\hline
\textbf{haben} & ich hatte & \textbf{ich \cle{hätte}} & ich würde haben \\
\hline
\textbf{sein} & ich war & \textbf{ich \cle{wäre}} & ich würde sein \\
\hline
\textbf{können} & ich konnte & \textbf{ich \cle{könnte}} & ich würde können \\
\hline
\textbf{müssen} & ich musste & \textbf{ich \cle{müsste}} & ich würde müssen \\
\hline
\textbf{dürfen} & ich durfte & \textbf{ich \cle{dürfte}} & ich würde dürfen \\
\hline
\textbf{sollen} & ich sollte & \textbf{ich \cle{sollte}} & ich würde sollen \\
\hline
\textbf{wollen} & ich wollte & \textbf{ich \cle{wollte}} & ich würde wollen \\
\hline
\textbf{werden} & ich wurde & \textbf{ich \cle{würde}} & — (auxiliaire) \\
\hline
\textbf{bestellen} (faible) & ich bestellte & (ich bestellte) & \textbf{ich \cle{würde bestellen}} \\
\hline
\textbf{schreiben} (fort) & ich schrieb & \textbf{ich \cle{schriebe}} & ich würde schreiben \\
\hline
\end{tabular}
\end{center}
\emph{Remarque :} Pour les verbes faibles (réguliers), la forme Präteritum et Konjunktiv II sont identiques (\textit{ich bestellte}) → on utilise \textbf{obligatoirement « würde + infinitif »} pour les distinguer.
\end{exemplebox}

\begin{propriete}[Emploi dans la lettre formelle — Formules types]
Le Konjunktiv II remplace l'impératif (trop direct) et l'indicatif (trop affirmatif) pour formuler des \textbf{demandes}, des \textbf{souhaits}, des \textbf{suggestions} ou des \textbf{réclamations} avec courtoisie.
\begin{center}
\begin{tabularx}{\linewidth}{|>{\bfseries}l|X|X|}
\hline
\rowcolor{popblueL} \textbf{Intention} & \textbf{Formule allemande (Konjunktiv II)} & \textbf{Traduction française} \\
\hline
Souhait / Demande polie & \all{Ich \cle{hätte} gern(e) einen neuen Kopfhörer.} / \all{Ich \cle{möchte} Sie bitten, ...} & J'aimerais (avoir) un nouveau casque. / Je voudrais vous prier de... \\
\hline
Question polie (capacité/possibilité) & \all{\cle{Könnten} Sie mir bitte sagen, ...?} / \all{\cle{Wären} Sie so freundlich, ...?} & Pourriez-vous me dire... ? / Auriez-vous l'amabilité de... ? \\
\hline
Demande formelle « Je vous prie de... » & \all{Ich \cle{würde} Sie bitten, mir \cle{zu antworten}.} & Je vous prierais de me répondre. \\
\hline
Condition polie (si + Konjunktiv II) & \all{Ich \cle{wäre} Ihnen sehr dankbar, \cle{wenn} Sie das \cle{prüfen} \cle{könnten}.} & Je vous serais très reconnaissant si vous pouviez vérifier. \\
\hline
Formule d'accompagnement & \all{Ich \cle{danke} Ihnen im Voraus.} (Indikatif, mais figé) / \all{Ich \cle{wäre} Ihnen für eine Rückmeldung dankbar.} & Je vous remercie d'avance. / Je vous saurais gré d'une réponse. \\
\hline
\end{tabularx}
\end{center}
\end{propriete}

\begin{attention}[Pièges francophones — Konjunktiv II]
\begin{itemize}
\item \textbf{Ne pas confondre « würde » (Konjunktiv II de werden) et « wurde » (Präteritum de werden = passif).} \all{Ich würde gern kommen} (j'aimerais venir) ≠ \all{Ich wurde gefragt} (j'ai été interrogé).
\item \textbf{Place du verbe :} En proposition principale, le verbe conjugué (\all{würde}, \all{könnte}, \all{hätte}) est en \textbf{deuxième position},
\end{itemize}
\end{attention}


\newpage

\coursec{Méthodes et exercices résolus}

\courssub{Méthode 1 : Comprendre une lettre de réclamation ou un CV (lecture globale puis détaillée)}

\begin{methode}[Lire un document administratif allemand]
\begin{enumerate}
\item \cle{Identifier le type de document} : lettre de réclamation (\all{der Beschwerdebrief}) ou curriculum vitae (\all{der Lebenslauf}). Indices : formules de politesse, date, adresse, objet (\all{Betreff}), rubriques (\all{Persönliche Daten}, \all{Schulbildung}, \all{Berufserfahrung}).
\item \cle{Repérer l'émetteur et le destinataire} : qui écrit ? à qui ? pourquoi ? (\all{der Absender} / \all{der Empfänger}).
\item \cle{Dégager la situation de communication} : problème exposé, demande formulée, ton (formel, courtois mais ferme).
\item \cle{Faire une deuxième lecture détaillée} : souligner les mots-clés du thème (\all{die Reklamation}, \all{der Mangel, die Mängel}, \all{die Kenntnisse}, \all{die Erfahrung}), repérer les chiffres, dates, lieux.
\item \cle{Reformuler en français} le contenu global en deux ou trois phrases avant de répondre aux questions.
\item \cle{Répondre aux questions} en allemand, par des phrases complètes, en reprenant les mots de la question, sans recopier tout le texte.
\end{enumerate}
\end{methode}

\begin{exoresolu}[Compréhension de l'écrit : une lettre de réclamation]
\textbf{Texte.} \all{Sehr geehrte Damen und Herren, am 12. Mai habe ich in Ihrem Geschäft in Douala einen Laptop gekauft. Leider funktioniert der Akku schon nach zwei Wochen nicht mehr. Ich bitte Sie, das Gerät zu reparieren oder umzutauschen. Ich hoffe auf eine schnelle Lösung. Mit freundlichen Grüßen Marie Ngo Bell}

\textbf{Questions.} 1. Wo hat die Kundin den Laptop gekauft? 2. Was ist das Problem? 3. Was verlangt die Kundin?
\tcblower
\textbf{Raisonnement.} On identifie d'abord le type de document : la formule d'appel \all{Sehr geehrte Damen und Herren} et la salutation finale \all{Mit freundlichen Grüßen} indiquent une lettre formelle ; le verbe \all{gekauft} et la demande de réparation montrent qu'il s'agit d'une réclamation (\all{die Beschwerde}). On repère ensuite les trois informations demandées : le lieu (\all{in Ihrem Geschäft in Douala}), le problème (\all{der Akku funktioniert nicht mehr}), la demande (\all{reparieren oder umtauschen}).

\textbf{Réponses rédigées.}
\begin{enumerate}
\item \all{Die Kundin hat den Laptop in einem Geschäft in Douala gekauft.} (La cliente a acheté l'ordinateur dans un magasin à Douala.)
\item \all{Der Akku funktioniert schon nach zwei Wochen nicht mehr.} (La batterie ne fonctionne plus au bout de deux semaines.)
\item \all{Sie verlangt, dass das Geschäft das Gerät repariert oder umtauscht.} (Elle exige que le magasin répare ou échange l'appareil.)
\end{enumerate}
\textbf{Conclusion.} Chaque réponse est une phrase complète, au Perfekt pour les faits passés et au présent pour l'état actuel, avec le verbe conjugué en deuxième position : c'est ce qu'attend le correcteur.
\end{exoresolu}

\courssub{Méthode 2 : Employer le vocabulaire de la Bewerbung, du Lebenslauf et de la Beschwerde}

\begin{methode}[Réactiver et enrichir le lexique thématique]
\begin{enumerate}
\item \cle{Apprendre chaque nom avec son article et son pluriel} : \all{der Lebenslauf, die Lebensläufe} ; \all{die Bewerbung, -en} ; \all{die Beschwerde, -n} ; \all{die Kenntnisse} (pluriel) ; \all{die Erfahrung, -en} ; \all{das Zeugnis, -se}.
\item \cle{Construire des familles de mots} : \all{sich bewerben} (postuler) $\rightarrow$ \all{die Bewerbung} $\rightarrow$ \all{der Bewerber / die Bewerberin} ; \all{sich beschweren} (se plaindre) $\rightarrow$ \all{die Beschwerde} ; \all{kennen} $\rightarrow$ \all{die Kenntnis}.
\item \cle{Mémoriser les constructions à préposition} : \all{sich bewerben um + Akk.} (\all{Ich bewerbe mich um die Stelle.} — je postule pour le poste) ; \all{sich beschweren über + Akk.} (se plaindre de) ; \all{sich bedanken für + Akk.} (remercier de).
\item \cle{Classer le lexique par rubrique du CV} : données personnelles (\all{Persönliche Daten}), formation (\all{Schulbildung / Ausbildung}), expérience (\all{Berufserfahrung}), compétences (\all{Kenntnisse und Fähigkeiten}), loisirs (\all{Hobbys}).
\item \cle{Réutiliser le lexique dans des phrases personnelles} liées au contexte camerounais (un stage à Yaoundé, un emploi à Douala).
\end{enumerate}
\end{methode}

\begin{definition}[Vocabulaire essentiel du thème]
\begin{center}
\begin{tabularx}{\linewidth}{|X|X|X|}\hline
\rowcolor{popblueL} Deutsch & Français & Remarque \\ \hline
\all{der Lebenslauf, die Lebensläufe} & le curriculum vitae & aussi \all{der CV} \\ \hline
\all{die Bewerbung, -en} & la candidature & \all{sich bewerben um + Akk.} \\ \hline
\all{die Beschwerde, -n} & la réclamation & \all{sich beschweren über + Akk.} \\ \hline
\all{die Reklamation, -en} & la réclamation (commerciale) & ne pas confondre avec \all{die Reklame} \\ \hline
\all{die Stelle, -n} & le poste, l'emploi & \all{eine Stelle als Verkäufer} \\ \hline
\all{die Kenntnisse (Pl.)} & les connaissances & \all{Sprachkenntnisse} \\ \hline
\all{die Berufserfahrung, -en} & l'expérience professionnelle & souvent sans article \\ \hline
\all{das Zeugnis, -se} & le diplôme, le bulletin & \all{das Abiturzeugnis} \\ \hline
\all{der Absender, - / die Empfängerin, -nen} & l'expéditeur / la destinataire & en-tête de lettre \\ \hline
\all{der Betreff, -e} & l'objet (de la lettre) & \all{Betreff: Bewerbung} \\ \hline
\all{der Anhang, die Anhänge} & la pièce jointe & \all{im Anhang} \\ \hline
\all{der Termin, -e} & le rendez-vous & \all{einen Termin bekommen} \\ \hline
\end{tabularx}
\end{center}
\end{definition}

\begin{exoresolu}[Lexique : compléter un CV et une phrase de candidature]
\textbf{Énoncé.} Complétez avec le mot qui convient : \all{Kenntnisse — Bewerbung — Berufserfahrung — Zeugnis — Stelle}.

1. \all{Ich bewerbe mich um die \dots\ als Verkäuferin.} 2. \all{Meine \dots\ schicke ich Ihnen heute per E-Mail.} 3. \all{Ich habe zwei Jahre \dots\ in einem Hotel in Kribi.} 4. \all{Meine Sprach\dots : Deutsch, Englisch, Französisch.} 5. \all{Im Anhang finden Sie mein Abitur\dots .}
\tcblower
\textbf{Raisonnement.} On s'appuie sur le sens et sur la grammaire de chaque phrase : article, cas, contexte.

\textbf{Solutions.}
\begin{enumerate}
\item \all{Stelle} : \all{sich bewerben um eine Stelle} (postuler pour un poste) — accusatif féminin après \all{um}.
\item \all{Bewerbung} : « j'envoie ma candidature aujourd'hui par courriel » — nominatif féminin, sujet de \all{schicke}.
\item \all{Berufserfahrung} : « j'ai deux ans d'expérience professionnelle dans un hôtel à Kribi » — nom féminin, employé ici sans article.
\item \all{Kenntnisse} : \all{Sprachkenntnisse} (connaissances linguistiques) — nom composé, toujours au pluriel dans un CV.
\item \all{Zeugnis} : \all{das Abiturzeugnis} (le diplôme du baccalauréat) — nom composé neutre, d'où \all{mein Abiturzeugnis}.
\end{enumerate}
\textbf{Conclusion.} Les noms composés (\all{Sprachkenntnisse}, \all{Abiturzeugnis}) prennent le genre du dernier élément : vérifier l'article est le réflexe qui fait gagner des points.
\end{exoresolu}

\courssub{Méthode 3 : Le subjonctif II de politesse}

\begin{propriete}[Formation du Konjunktiv II de politesse]
Le Konjunktiv II se forme à partir du \cle{Präteritum}, avec \cle{Umlaut} pour les verbes forts (a, o, u $\rightarrow$ ä, ö, ü) :

\begin{center}
\begin{tabular}{|l|l|l|}\hline
\rowcolor{popblueL} Infinitif & Präteritum & Konjunktiv II \\ \hline
\all{haben} & \all{hatte} & \all{hätte} \\ \hline
\all{sein} & \all{war} & \all{wäre} \\ \hline
\all{können} & \all{konnte} & \all{könnte} \\ \hline
\all{werden} & \all{wurde} & \all{würde} \\ \hline
\all{geben} & \all{gab} & \all{gäbe} \\ \hline
\all{kommen} & \all{kam} & \all{käme} \\ \hline
\end{tabular}
\end{center}

Pour la plupart des verbes, on préfère la forme de remplacement : \all{würde} + \cle{infinitif} (\all{ich würde kaufen}, \all{wir würden uns freuen}). Les formes simples \all{hätte}, \all{wäre}, \all{wüsste} restent très fréquentes.
\end{propriete}

\begin{methode}[Formuler une demande polie avec le Konjunktiv II]
\begin{enumerate}
\item \cle{Retenir les trois formes vedettes} : \all{ich hätte} (j'aurais), \all{ich würde} + infinitif (je voudrais / je ferais), \all{könnten Sie} (pourriez-vous).
\item \cle{Demander poliment} : \all{Ich hätte gern Informationen.} (Je voudrais des renseignements.) ; \all{Ich würde gern einen Termin bekommen.} (J'aimerais obtenir un rendez-vous.)
\item \cle{Poser une question polie} : \all{Könnten Sie mir bitte das Formular schicken?} (Pourriez-vous m'envoyer le formulaire, s'il vous plaît ?)
\item \cle{Conseiller ou suggérer} : \all{Sie sollten das Gerät umtauschen.} (Vous devriez échanger l'appareil.)
\item \cle{Exprimer un souhait ou un contentement} : \all{Ich würde mich über eine schnelle Antwort freuen.} (Je me réjouirais d'une réponse rapide.)
\item \cle{Respecter la place du verbe} : \all{würde} en deuxième position, l'infinitif à la fin de la phrase.
\end{enumerate}
\end{methode}

\begin{exoresolu}[Grammaire : transformer des demandes directes en demandes polies]
\textbf{Énoncé.} Réécrivez au Konjunktiv II de politesse : 1. \all{Geben Sie mir die Rechnung!} 2. \all{Schicken Sie mir ein neues Gerät!} 3. \all{Ich will einen Termin.} 4. \all{Reparieren Sie den Laptop!}
\tcblower
\textbf{Raisonnement.} L'impératif direct est trop brutal dans une lettre formelle allemande. On le remplace par \all{Könnten Sie ... ?} (question polie, infinitif en fin de phrase) ou par \all{Ich hätte gern ...} / \all{Ich würde gern ...}.

\textbf{Solutions.}
\begin{enumerate}
\item \all{Könnten Sie mir bitte die Rechnung geben?} (Pourriez-vous me donner la facture, s'il vous plaît ?) — \all{könnten} en tête, infinitif \all{geben} à la fin.
\item \all{Würden Sie mir bitte ein neues Gerät schicken?} (Pourriez-vous m'envoyer un appareil neuf ?)
\item \all{Ich hätte gern einen Termin.} (Je voudrais un rendez-vous.) — \all{hätte} + \all{gern}, accusatif \all{einen Termin}.
\item \all{Könnten Sie bitte den Laptop reparieren?} (Pourriez-vous réparer l'ordinateur ?)
\end{enumerate}
\textbf{Conclusion.} Au Bac, toute lettre de réclamation doit contenir au moins une formule au Konjunktiv II : c'est un critère explicite de la grille d'évaluation.
\end{exoresolu}

\courssub{Méthode 4 : Rédiger une lettre formelle (structure et formules)}

\begin{methode}[Construire une lettre de réclamation ou de candidature]
\begin{enumerate}
\item \cle{En-tête} : expéditeur (\all{der Absender}) en haut à gauche, destinataire (\all{der Empfänger}) en dessous, lieu et date à droite : \all{Yaoundé, den 15. Juni 2024} (accusatif de date après \all{den}).
\item \cle{Objet} : \all{Betreff: Beschwerde über einen defekten Laptop} ou \all{Betreff: Bewerbung um eine Stelle als Verkäuferin}.
\item \cle{Formule d'appel} : \all{Sehr geehrte Damen und Herren,} (destinataire inconnu) ou \all{Sehr geehrter Herr Mbarga,} (virgule obligatoire, puis ligne suivante qui commence par une minuscule).
\item \cle{Corps en trois paragraphes} : (a) rappel des faits avec date précise ; (b) description du problème ou présentation de sa candidature ; (c) demande claire au Konjunktiv II.
\item \cle{Formule de clôture} : \all{Ich hoffe auf eine schnelle Antwort.} puis \all{Mit freundlichen Grüßen} (sans virgule !) et la signature.
\item \cle{Vérifier} : majuscules aux noms, place du verbe, cas après prépositions (\all{über + Akk.}, \all{mit + Dat.}, \all{um + Akk.}, \all{für + Akk.}).
\end{enumerate}
\end{methode}

\begin{exemplebox}[Modèle de CV allemand (rubriques)]
\all{\textbf{Lebenslauf}}

\all{\textbf{Persönliche Daten:} Name: Amina Fotso; Geburtsdatum: 12.03.2006; Adresse: Yaoundé, Kamerun}

\all{\textbf{Schulbildung:} 2018--2024 Collège/Lycée in Yaoundé; 2024 Abitur (Baccalauréat)}

\all{\textbf{Berufserfahrung:} 2023 Praktikum in einem Reisebüro in Douala}

\all{\textbf{Kenntnisse:} Sprachen: Französisch, Deutsch, Englisch; Informatik: Word, Excel}

\all{\textbf{Hobbys:} Fußball, Lesen, Musik}

« Curriculum vitae — Données personnelles : nom, date de naissance, adresse — Formation scolaire — Expérience professionnelle : stage dans une agence de voyages à Douala — Connaissances : langues, informatique — Loisirs. »
\end{exemplebox}

\begin{exoresolu}[Expression écrite : rédiger une courte lettre de réclamation]
\textbf{Énoncé.} Vous avez acheté un téléphone portable dans un magasin de Douala. Après une semaine, l'écran ne fonctionne plus. Rédigez une lettre de réclamation (environ 60--80 mots) avec les formules adaptées.
\tcblower
\textbf{Raisonnement.} On suit le plan en cinq étapes : en-tête, objet, appel, corps en trois paragraphes (faits — problème — demande au Konjunktiv II), clôture. On utilise le Perfekt pour les faits passés (\all{habe gekauft}) et le présent pour le problème actuel (\all{funktioniert nicht}).

\textbf{Modèle de réponse.}

\all{Yaoundé, den 20. März 2024}

\all{Betreff: Beschwerde über ein defektes Handy}

\all{Sehr geehrte Damen und Herren,}

\all{am 10. März habe ich in Ihrem Geschäft in Douala ein Handy gekauft. Leider funktioniert der Bildschirm schon nach einer Woche nicht mehr. Das Gerät ist also defekt, obwohl es ganz neu ist.}

\all{Ich bitte Sie, das Handy zu reparieren oder umzutauschen. Ich würde mich auch über eine schnelle Antwort freuen.}

\all{Mit freundlichen Grüßen}

\all{Paul Etoundi}

« Yaoundé, le 20 mars 2024 — Objet : réclamation au sujet d'un téléphone défectueux — Mesdames, Messieurs, le 10 mars j'ai acheté un téléphone dans votre magasin à Douala. Malheureusement, l'écran ne fonctionne plus au bout d'une semaine. L'appareil est donc défectueux, alors qu'il est tout neuf. Je vous prie de réparer ou d'échanger le téléphone. Je me réjouirais aussi d'une réponse rapide. Cordialement, Paul Etoundi. »

\textbf{Justifications grammaticales.} \all{habe ... gekauft} : Perfekt avec l'auxiliaire \all{haben} ; \all{über + Akk.} (\all{über ein defektes Handy}) ; \all{Ich würde mich freuen} : Konjunktiv II de politesse avec infinitif en fin de phrase ; \all{Mit freundlichen Grüßen} sans virgule.
\end{exoresolu}

\courssub{Méthode 5 : Traduire (thème et version) les formules de la lettre formelle}

\begin{methode}[Traduire sans calquer le français]
\begin{enumerate}
\item \cle{Repérer les faux amis} : « une réclamation » = \all{die Beschwerde / die Reklamation}, pas \all{die Reklame} (la publicité !) ; « un poste » = \all{die Stelle}, pas \all{der Posten} au sens courant ; « actuellement » = \all{zurzeit / im Moment}, alors que \all{aktuell} signifie « d'actualité ».
\item \cle{Traduire les formules figées apprises par cœur} : « Veuillez agréer... » $\rightarrow$ \all{Mit freundlichen Grüßen} ; « Je vous prie de... » $\rightarrow$ \all{Ich bitte Sie, ... zu + Infinitiv}.
\item \cle{Attention à la place du verbe} : « Je vous serais reconnaissant de bien vouloir échanger l'appareil » $\rightarrow$ \all{Ich wäre Ihnen dankbar, wenn Sie das Gerät umtauschen würden.} (subordonnée avec \all{wenn} : verbe conjugué à la fin).
\item \cle{Datif ou accusatif} : \all{Ich danke Ihnen für Ihre Antwort.} (\all{danken} + datif ; \all{für} + accusatif).
\item \cle{Relire} : majuscule à chaque nom, Umlaute, ponctuation allemande.
\end{enumerate}
\end{methode}

\begin{exoresolu}[Thème : traduire en allemand]
\textbf{Énoncé.} Traduisez : 1. « Je me permets de vous écrire au sujet de ma commande. » 2. « Pourriez-vous m'envoyer un nouveau produit ? » 3. « Je vous remercie de votre compréhension. » 4. « J'ai de bonnes connaissances en allemand et en informatique. »
\tcblower
\textbf{Raisonnement.} Phrase 1 : « au sujet de » se rend simplement par \all{wegen + Gen.} ; phrase 2 : Konjunktiv II de politesse ; phrase 3 : \all{danken + Dat.} avec \all{für + Akk.} ; phrase 4 : \all{Kenntnisse in + Dat.}.

\textbf{Solutions.}
\begin{enumerate}
\item \all{Ich erlaube mir, Ihnen wegen meiner Bestellung zu schreiben.} (\all{wegen} + génitif \all{meiner Bestellung} ; infinitif \all{zu schreiben} à la fin.)
\item \all{Könnten Sie mir ein neues Produkt schicken?} (Konjunktiv II, infinitif final.)
\item \all{Ich danke Ihnen für Ihr Verständnis.} (\all{danken} + datif \all{Ihnen} ; \all{für} + accusatif.)
\item \all{Ich habe gute Kenntnisse in Deutsch und Informatik.} (\all{Kenntnisse} toujours au pluriel ; adjectif \all{gute} au pluriel accusatif.)
\end{enumerate}
\textbf{Conclusion.} Chaque phrase a été vérifiée : cas après préposition, place du verbe, majuscules aux noms (\all{Bestellung}, \all{Produkt}, \all{Verständnis}, \all{Kenntnisse}, \all{Deutsch}, \all{Informatik}).
\end{exoresolu}

\begin{savaistu}[La lettre formelle dans les pays germanophones]
En Allemagne, en Autriche et en Suisse, la lettre formelle reste un genre très codifié : on n'écrit jamais de virgule après \all{Mit freundlichen Grüßen}, et la première ligne du corps de la lettre commence par une minuscule après la virgule de l'appel. Au Cameroun, où l'allemand ouvre des portes (bourses du DAAD, coopération germano-camerounaise, entreprises germanophones à Douala), maîtriser le \all{Lebenslauf} et la \all{Bewerbung} est un atout concret pour les études et l'emploi.
\end{savaistu}

\begin{attention}[Les 5 erreurs qui coûtent des points au Bac]
\begin{enumerate}
\item \cle{Confondre Reklamation et Reklame} : \all{die Reklamation} = la réclamation ; \all{die Reklame} = la publicité. Faux ami classique, très pénalisé.
\item \cle{Oublier le Konjunktiv II de politesse} : écrire \all{Geben Sie mir ...} (impératif brutal) au lieu de \all{Könnten Sie mir bitte ... geben?} dans une lettre formelle.
\item \cle{Mettre une virgule après la formule de clôture} : \all{Mit freundlichen Grüßen} s'écrit SANS virgule en allemand (contrairement au français « Cordialement, »).
\item \cle{Oublier les majuscules et les cas} : \all{mit freundlichen Grüßen} (datif pluriel, sans article), \all{über + Akkusativ} (\all{Beschwerde über einen Laptop}), \all{sich bewerben um + Akkusativ}.
\item \cle{Calquer l'ordre des mots français} : « Je voudrais recevoir une réponse » ne se traduit pas mot à mot ; il faut \all{Ich würde gern eine Antwort bekommen.} — verbe conjugué en deuxième position, infinitif à la fin.
\end{enumerate}
\end{attention}

\newpage

\coursec{Exercices}

\courssub{Je vérifie mes connaissances}

\begin{exercice}[QCM : la lettre formelle]
Entoure la bonne réponse.
\begin{enumerate}
\item Pour commencer une lettre formelle adressée à une entreprise, on écrit :
\begin{enumerate}
\item \all{Liebe Damen und Herren,}
\item \all{Sehr geehrte Damen und Herren,}
\item \all{Hallo zusammen,}
\end{enumerate}
\item La formule de politesse finale correcte est :
\begin{enumerate}
\item \all{Mit freundlichen Grüßen}
\item \all{Viele Grüße}
\item \all{Tschüss und bis bald}
\end{enumerate}
\item \all{Der Betreff} désigne :
\begin{enumerate}
\item l'adresse de l'expéditeur
\item l'objet de la lettre
\item la signature
\end{enumerate}
\item \all{Der Lebenslauf} signifie :
\begin{enumerate}
\item la lettre de motivation
\item le contrat de travail
\item le curriculum vitae
\end{enumerate}
\end{enumerate}
\end{exercice}

\begin{exercice}[Vrai ou faux ? Justifie ta réponse]
Dis si les affirmations suivantes sont vraies (V) ou fausses (F) et justifie en une phrase en français.
\begin{enumerate}
\item \all{Könnten Sie mir bitte helfen?} est plus poli que \all{Helfen Sie mir!}
\item Dans un CV allemand, on ne mentionne jamais la date de naissance.
\item \all{Ich hätte gern einen Termin.} contient un subjonctif II.
\item Une lettre de réclamation se termine toujours par \all{Liebe Grüße}.
\item \all{die Kenntnisse} est un nom qui s'emploie presque toujours au pluriel.
\end{enumerate}
\end{exercice}

\begin{exercice}[Vocabulaire : trouve le mot]
Complète chaque phrase avec le mot thématique qui convient : \all{die Bewerbung, die Beschwerde, der Lebenslauf, die Berufserfahrung, die Sprachkenntnisse}.
\begin{enumerate}
\item Herr Mbarga schreibt einen Brief an die Firma, weil das Paket kaputt ist. Das ist eine \ldots
\item Für die Stelle in Douala schickt Amina ihren \ldots{} mit Foto und Unterschrift.
\item Deutsch, Englisch und Französisch — das sind ihre \ldots
\item Er hat drei Jahre in einer Bank gearbeitet: das ist seine \ldots
\item Sie schickt ihren Lebenslauf und ein Anschreiben: das ist ihre \ldots
\end{enumerate}
\end{exercice}

\begin{exercice}[Conjugaison : le subjonctif II de politesse]
Complète avec la forme correcte du subjonctif II du verbe entre parenthèses.
\begin{enumerate}
\item \all{\ldots{} Sie mir bitte das Formular schicken? (können)}
\item \all{Ich \ldots{} gern einen Termin bei Ihnen. (haben)}
\item \all{\ldots{} Sie so freundlich, mir zu antworten? (sein)}
\item \all{Ich \ldots{} Sie bitten, das Gerät zu reparieren. (mögen)}
\item \all{\ldots{} Sie mir bitte sagen, wann das Paket ankommt? (wollen)}
\end{enumerate}
\end{exercice}

\courssub{Je m'entraîne}

\begin{exercice}[Thème : traduis en allemand]
Traduis les phrases suivantes en allemand correct, en utilisant le subjonctif II de politesse quand c'est nécessaire.
\begin{enumerate}
\item Je voudrais vous demander de m'envoyer le formulaire.
\item Pourriez-vous me donner votre adresse électronique ?
\item Je serais très content(e) d'une réponse rapide.
\item Auriez-vous la gentillesse de réparer l'appareil ?
\item Je vous prie de bien vouloir me rembourser.
\end{enumerate}
\end{exercice}

\begin{exercice}[Transformation : de l'impératif à la politesse]
Transforme ces demandes directes en demandes polies au subjonctif II, selon le modèle.\\[2pt]
\emph{Modèle :} \all{Geben Sie mir das Dokument!} $\rightarrow$ \all{Könnten Sie mir bitte das Dokument geben?}
\begin{enumerate}
\item \all{Schicken Sie mir die Rechnung!}
\item \all{Rufen Sie mich morgen an!}
\item \all{Reparieren Sie das Handy!}
\item \all{Sagen Sie mir den Preis!}
\item \all{Kommen Sie bitte um 10 Uhr!}
\end{enumerate}
\end{exercice}

\begin{exercice}[Texte à trous : un CV]
Complète ce CV avec les mots suivants : \all{Schulbildung, Berufserfahrung, Sprachkenntnisse, Geburtsdatum, Adresse, Unterschrift}.
\begin{textebox}[Lebenslauf]
\textbf{Lebenslauf}\\[4pt]
Name: Amina Fotso\\
\ldots{} (1): 12. März 2007, Yaoundé\\
\ldots{} (2): Quartier Nlongkak, Yaoundé, Kamerun\\[4pt]
\textbf{\ldots{} (3):}\\
2013--2019: Grundschule in Yaoundé\\
2019--2026: Gymnasium, Abitur (Baccalauréat A)\\[4pt]
\textbf{\ldots{} (4):}\\
Deutsch (sehr gut), Englisch (gut), Französisch (Muttersprache)\\[4pt]
\textbf{\ldots{} (5):}\\
2024: Praktikum in einem Reisebüro in Douala\\[4pt]
Yaoundé, den 15. Mai 2026 \hfill \ldots{} (6): \emph{Amina Fotso}
\end{textebox}
\end{exercice}

\begin{exercice}[Appariement : la structure de la lettre formelle]
Remets dans le bon ordre les huit éléments suivants d'une lettre formelle (numérote de 1 à 8).
\begin{enumerate}
\item[\ldots] \all{Mit freundlichen Grüßen}
\item[\ldots] \all{Betreff: Beschwerde über ein defektes Handy}
\item[\ldots] Adresse de l'expéditeur
\item[\ldots] \all{Sehr geehrte Damen und Herren,}
\item[\ldots] Lieu et date : \all{Yaoundé, den 15. Mai 2026}
\item[\ldots] Corps de la lettre (présentation du problème et demande)
\item[\ldots] Adresse du destinataire
\item[\ldots] Signature
\end{enumerate}
\end{exercice}

\courssub{Je me prépare au Bac}

\begin{exercice}[Type Bac — Compréhension écrite]
Lis la lettre suivante, puis réponds aux questions \textbf{en allemand, par des phrases complètes}.
\begin{textebox}[Eine Beschwerde]
Amina Fotso \hfill Reisebüro Kamerun-Reisen\\
Quartier Nlongkak \hfill Avenue Kennedy\\
Yaoundé, Kamerun \hfill Douala, Kamerun\\[4pt]
Yaoundé, den 15. Mai 2026\\[4pt]
\textbf{Betreff: Beschwerde über eine Busreise nach Kribi}\\[4pt]
Sehr geehrte Damen und Herren,\\[4pt]
am 2. Mai habe ich bei Ihnen eine Busreise nach Kribi für meine Familie gebucht. Leider war ich mit dieser Reise sehr unzufrieden. Der Bus kam zwei Stunden zu spät, die Klimaanlage funktionierte nicht, und der Fahrer war unfreundlich. Für den Preis von 25\,000 FCFA pro Person erwarte ich einen besseren Service.\\[4pt]
Ich möchte Sie bitten, mir die Hälfte des Preises zurückzuzahlen. Könnten Sie mir außerdem schriftlich bestätigen, dass solche Probleme in Zukunft nicht mehr vorkommen? Ich wäre Ihnen sehr dankbar für eine schnelle Antwort.\\[4pt]
Mit freundlichen Grüßen\\
\emph{Amina Fotso}
\end{textebox}
\begin{enumerate}
\item Wer schreibt diesen Brief, und an wen? (2 pts)
\item Worum geht es in diesem Brief? (2 pts)
\item Wann hat Amina die Reise gebucht? (1 pt)
\item Nennen Sie die drei Probleme während der Reise. (3 pts)
\item Wie viel hat die Reise pro Person gekostet? (1 pt)
\item Was verlangt Amina von der Firma? (2 pts)
\item Welche Höflichkeitsform benutzt Amina am Ende des Briefes? (1 pt)
\item Trouve dans le texte deux phrases au subjonctif II et recopie-les. (2 pts)
\end{enumerate}
\textbf{Vrai ou faux ?} Justifie ta réponse par une citation du texte.
\begin{enumerate}
\item[a)] Der Bus war pünktlich. (2 pts)
\item[b)] Amina möchte das ganze Geld zurück. (2 pts)
\end{enumerate}
\end{exercice}

\begin{exercice}[Type Bac — Thème et version]
\textbf{A. Thème (10 pts).} Traduis en allemand les cinq phrases suivantes, extraites d'une lettre de réclamation. Attention à la place du verbe et au subjonctif II de politesse.
\begin{enumerate}
\item J'ai commandé un téléphone portable chez vous le 10 avril.
\item Malheureusement, l'appareil ne fonctionne pas.
\item Je voudrais vous demander de me rembourser.
\item Pourriez-vous m'envoyer un nouvel appareil ?
\item Je vous serais reconnaissant(e) d'une réponse rapide.
\end{enumerate}
\textbf{B. Version (10 pts).} Traduis en français l'extrait suivant d'une lettre formelle.
\begin{textebox}[Auszug aus einem offiziellen Brief]
\textbf{Betreff: Ihre Bewerbung um ein Praktikum}\\[4pt]
Sehr geehrter Herr Nkoulou,\\[4pt]
vielen Dank für Ihre Bewerbung und Ihren Lebenslauf. Wir hätten gern mehr Informationen über Ihre Sprachkenntnisse. Könnten Sie uns bitte auch sagen, wann Sie bei uns anfangen könnten? Wir würden uns sehr freuen, Sie bald kennenzulernen.\\[4pt]
Mit freundlichen Grüßen\\
\emph{Dr. Sabine Weber, Personalchefin}
\end{textebox}
\end{exercice}

\begin{exercice}[Type Bac — Expression écrite]
\textbf{Sujet :} Tu as commandé sur Internet un téléphone portable (\all{das Handy}) qui est arrivé cassé (\all{defekt}). Rédige une lettre de réclamation à l'entreprise.\\[4pt]
\emph{Consignes :} rédige en allemand une lettre d'environ \textbf{80 mots}. Ta lettre doit contenir obligatoirement :
\begin{itemize}
\item les adresses, le lieu, la date et le \all{Betreff} ;
\item la formule d'appel \all{Sehr geehrte Damen und Herren,} ;
\item la présentation du problème (au moins deux détails) ;
\item une demande polie au subjonctif II (par ex. \all{Könnten Sie ...?}, \all{Ich möchte Sie bitten, ...}) ;
\item la formule finale \all{Mit freundlichen Grüßen} et ta signature.
\end{itemize}
\emph{Barème indicatif :} respect de la forme de la lettre formelle (4 pts), exactitude grammaticale et place du verbe (6 pts), richesse du vocabulaire thématique (6 pts), cohérence du contenu (4 pts).
\end{exercice}

\newpage

\coursec{Corrigés des exercices}

\begin{corrige}[Exercice 1 — QCM : la lettre formelle]
\begin{enumerate}
\item \textbf{b)} \all{Sehr geehrte Damen und Herren,} — c'est la formule d'appel standard d'une lettre formelle adressée à une entreprise dont on ne connaît pas le destinataire précis. \all{Liebe Damen und Herren} est trop familier, \all{Hallo zusammen} est réservé aux courriels entre amis ou collègues. Attention à la virgule finale et à la majuscule de \all{geehrte} après \all{Sehr}.
\item \textbf{a)} \all{Mit freundlichen Grüßen} — formule de politesse finale obligatoire dans une lettre formelle (souvent abrégée \all{MfG} dans les courriels professionnels, mais on écrit la formule entière dans une lettre). \all{Viele Grüße} et \all{Tschüss} sont familiers. Remarque : pas de virgule après \all{Grüßen}.
\item \textbf{b)} l'objet de la lettre — \all{der Betreff} (« l'objet ») annonce en une ligne le sujet de la lettre, par exemple \all{Betreff: Beschwerde über ein defektes Handy}.
\item \textbf{c)} le curriculum vitae — \all{der Lebenslauf} (littéralement « le cours de la vie ») est le mot allemand pour CV. La lettre de motivation se dit \all{das Anschreiben} ou \all{das Motivationsschreiben}, le contrat de travail \all{der Arbeitsvertrag}.
\end{enumerate}
\end{corrige}

\begin{corrige}[Exercice 2 — Vrai ou faux ?]
\begin{enumerate}
\item \textbf{V.} \all{Könnten Sie mir bitte helfen?} emploie le subjonctif II de \all{können} (\all{könnten}), qui adoucit la demande ; l'impératif \all{Helfen Sie mir!} est un ordre direct, impoli en contexte formel.
\item \textbf{V.} Dans les pays germanophones, le CV mentionne traditionnellement la date et le lieu de naissance (\all{das Geburtsdatum}, \all{der Geburtsort}), même si cette habitude tend à évoluer. \emph{Correction de l'énoncé : l'affirmation était donc vraie.}
\item \textbf{V.} \all{hätte} est le subjonctif II de \all{haben} (Präteritum \all{hatte} + tréma) ; \all{ich hätte gern} signifie « je voudrais ».
\item \textbf{F.} Une lettre de réclamation est une lettre formelle : elle se termine par \all{Mit freundlichen Grüßen}. \all{Liebe Grüße} est réservé aux lettres personnelles (amis, famille).
\item \textbf{V.} \all{die Kenntnisse} (les connaissances) s'emploie presque toujours au pluriel dans les composés courants : on dit \all{Sprachkenntnisse}, \all{Computerkenntnisse}. Le singulier \all{die Kenntnis} existe mais signifie plutôt « la connaissance d'un fait » (\all{Kenntnis von etwas haben}).
\end{enumerate}
\end{corrige}

\begin{corrige}[Exercice 3 — Vocabulaire : trouve le mot]
\begin{enumerate}
\item \textbf{Beschwerde} — une lettre pour se plaindre d'un paquet cassé est une réclamation (\all{die Beschwerde, -n}). Famille : \all{sich beschweren über} + accusatif (« se plaindre de »).
\item \textbf{Lebenslauf} — Amina envoie son CV avec photo et signature (\all{der Lebenslauf, die Lebensläufe}).
\item \textbf{Sprachkenntnisse} — les langues parlées sont les connaissances linguistiques (\all{die Sprachkenntnisse}, pluriel ; composé de \all{die Sprache} + \all{die Kenntnisse}).
\item \textbf{Berufserfahrung} — trois ans de travail dans une banque, c'est l'expérience professionnelle (\all{die Berufserfahrung, -en} ; composé de \all{der Beruf} + \all{die Erfahrung}).
\item \textbf{Bewerbung} — le CV plus la lettre de motivation forment la candidature (\all{die Bewerbung, -en}). Famille : \all{sich bewerben um} + accusatif (« poser sa candidature pour ») : \all{Sie bewirbt sich um ein Praktikum.} « Elle pose sa candidature pour un stage. »
\end{enumerate}
\end{corrige}

\begin{corrige}[Exercice 4 — Conjugaison : le subjonctif II de politesse]
\begin{enumerate}
\item \all{\textbf{Könnten} Sie mir bitte das Formular schicken?} — subjonctif II de \all{können} : Präteritum \all{konnte} + tréma $\rightarrow$ \all{könnten} (2\up{e} pers. plur.).
\item \all{Ich \textbf{hätte} gern einen Termin bei Ihnen.} — subjonctif II de \all{haben} : \all{hatte} + tréma $\rightarrow$ \all{hätte}.
\item \all{\textbf{Wären} Sie so freundlich, mir zu antworten?} — subjonctif II de \all{sein} : \all{war} + tréma $\rightarrow$ \all{wären}.
\item \all{Ich \textbf{möchte} Sie bitten, das Gerät zu reparieren.} — subjonctif II de \all{mögen} : \all{mochte} + tréma $\rightarrow$ \all{möchte} (forme très fréquente : « j'aimerais »).
\item \all{\textbf{Würden} Sie mir bitte sagen, wann das Paket ankommt?} — le subjonctif II de \all{wollen} est identique au Präteritum (\all{wollten}), donc on préfère ici \all{würden} + infinitif, forme de politesse la plus naturelle. (\all{Wollten Sie mir bitte sagen...} est grammaticalement correct mais moins courant.)
\end{enumerate}
\textbf{Rappel de formation :} subjonctif II = radical du Präteritum + tréma sur la voyelle du radical (a, o, u $\rightarrow$ ä, ö, ü) + terminaisons. Pour les verbes faibles (dont le subjonctif II serait identique au Präteritum) et la plupart des verbes lexicaux, on utilise \all{würde} + infinitif. Les formes à connaître par cœur au niveau Terminale : \all{wäre} (sein), \all{hätte} (haben), \all{könnte} (können), \all{müsste} (müssen), \all{dürfte} (dürfen), \all{sollte} (sollen), \all{möchte} (mögen), \all{würde} (werden).
\end{corrige}

\begin{corrige}[Exercice 5 — Thème : traduis en allemand]
\begin{enumerate}
\item \all{Ich möchte Sie bitten, mir das Formular zu schicken.} — \all{jemanden bitten, etwas zu tun} : construction avec infinitif à \all{zu}. Autre solution : \all{Ich würde Sie gern bitten, mir das Formular zu schicken.}
\item \all{Könnten Sie mir bitte Ihre E-Mail-Adresse geben?} — subjonctif II de \all{können} ; l'infinitif \all{geben} se place en fin de phrase.
\item \all{Ich würde mich über eine schnelle Antwort sehr freuen.} — attention : « se réjouir de » se dit \all{sich freuen über} + accusatif. Autre solution : \all{Ich wäre sehr froh über eine schnelle Antwort.}
\item \all{Wären Sie so freundlich, das Gerät zu reparieren?} — ou : \all{Könnten Sie bitte das Gerät reparieren?}
\item \all{Ich möchte Sie bitten, mir das Geld zurückzuzahlen.} — « rembourser » = \all{zurückzahlen} ou \all{erstatten} ; le verbe à particule \all{zurückzahlen} donne l'infinitif \all{zurückzuzahlen} (\all{zu} s'insère entre la particule et le verbe).
\end{enumerate}
\textbf{Points de vigilance :} verbe conjugué en 2\up{e} position, infinitif en fin de phrase, \all{mir} (datif) pour « me/m' », majuscule aux noms (\all{Formular}, \all{Antwort}, \all{Gerät}, \all{Geld}).
\end{corrige}

\begin{corrige}[Exercice 6 — Transformation : de l'impératif à la politesse]
\begin{enumerate}
\item \all{Schicken Sie mir die Rechnung!} $\rightarrow$ \all{Könnten Sie mir bitte die Rechnung schicken?}
\item \all{Rufen Sie mich morgen an!} $\rightarrow$ \all{Könnten Sie mich bitte morgen anrufen?} — attention : \all{anrufen} est un verbe à particule séparable ; à l'infinitif en fin de phrase, la particule se recolle : \all{anrufen}.
\item \all{Reparieren Sie das Handy!} $\rightarrow$ \all{Könnten Sie bitte das Handy reparieren?}
\item \all{Sagen Sie mir den Preis!} $\rightarrow$ \all{Könnten Sie mir bitte den Preis sagen?}
\item \all{Kommen Sie bitte um 10 Uhr!} $\rightarrow$ \all{Könnten Sie bitte um 10 Uhr kommen?}
\end{enumerate}
\textbf{Méthode :} on remplace l'impératif par \all{Könnten Sie ...?} (subjonctif II de \all{können}), le verbe lexical passe à l'infinitif en fin de phrase, et on ajoute \all{bitte}. D'autres formules sont possibles : \all{Würden Sie mir bitte die Rechnung schicken?}, \all{Wären Sie so freundlich, mich morgen anzurufen?} (avec \all{anzurufen} : \all{zu} inséré entre la particule et le verbe).
\end{corrige}

\begin{corrige}[Exercice 7 — Texte à trous : un CV]
\begin{enumerate}
\item \textbf{Geburtsdatum} : 12. März 2007, Yaoundé — la date de naissance (\all{das Geburtsdatum, die Geburtsdaten}).
\item \textbf{Adresse} : Quartier Nlongkak, Yaoundé, Kamerun — l'adresse (\all{die Adresse, -n}).
\item \textbf{Schulbildung} — la rubrique qui liste l'école primaire (\all{die Grundschule}) et le lycée (\all{das Gymnasium}) est la formation scolaire (\all{die Schulbildung}, sans pluriel).
\item \textbf{Sprachkenntnisse} — Deutsch, Englisch, Französisch : les connaissances linguistiques.
\item \textbf{Berufserfahrung} — le stage (\all{das Praktikum, die Praktika}) dans une agence de voyages relève de l'expérience professionnelle.
\item \textbf{Unterschrift} — la signature (\all{die Unterschrift, -en}) termine le CV allemand, précédée du lieu et de la date.
\end{enumerate}
\textbf{À retenir :} un CV allemand (\all{der Lebenslauf}) est structuré en rubriques : données personnelles (\all{Persönliche Daten}), formation scolaire (\all{Schulbildung}), expérience professionnelle (\all{Berufserfahrung}), connaissances linguistiques (\all{Sprachkenntnisse}), et se termine par lieu, date et signature.
\end{corrige}

\begin{corrige}[Exercice 8 — Appariement : la structure de la lettre formelle]
Ordre correct :
\begin{enumerate}
\item[\textbf{7}] \all{Mit freundlichen Grüßen}
\item[\textbf{4}] \all{Betreff: Beschwerde über ein defektes Handy}
\item[\textbf{1}] Adresse de l'expéditeur
\item[\textbf{5}] \all{Sehr geehrte Damen und Herren,}
\item[\textbf{3}] Lieu et date : \all{Yaoundé, den 15. Mai 2026}
\item[\textbf{6}] Corps de la lettre (présentation du problème et demande)
\item[\textbf{2}] Adresse du destinataire
\item[\textbf{8}] Signature
\end{enumerate}
\textbf{Schéma de la lettre formelle :} (1) adresse de l'expéditeur en haut à gauche $\rightarrow$ (2) adresse du destinataire $\rightarrow$ (3) lieu et date à droite $\rightarrow$ (4) \all{Betreff} en gras $\rightarrow$ (5) formule d'appel suivie d'une virgule $\rightarrow$ (6) corps de la lettre (la première ligne commence par une minuscule après la virgule de l'appel) $\rightarrow$ (7) formule de politesse finale $\rightarrow$ (8) signature.
\end{corrige}

\begin{corrige}[Exercice 9 — Type Bac : Compréhension écrite]
\textbf{Barème indicatif : 16 points (questions) + 4 points (vrai/faux).}

\textbf{Réponses modèles (phrases complètes en allemand) :}
\begin{enumerate}
\item \all{Amina Fotso aus Yaoundé schreibt diesen Brief an das Reisebüro Kamerun-Reisen in Douala.} (2 pts : 1 pt pour l'expéditrice, 1 pt pour le destinataire.)
\item \all{In diesem Brief geht es um eine Beschwerde über eine Busreise nach Kribi. Amina war mit der Reise sehr unzufrieden.} (2 pts.)
\item \all{Sie hat die Reise am 2. Mai gebucht.} (1 pt.)
\item \all{Die drei Probleme waren: Der Bus kam zwei Stunden zu spät, die Klimaanlage funktionierte nicht, und der Fahrer war unfreundlich.} (3 pts : 1 pt par problème.)
\item \all{Die Reise hat 25\,000 FCFA pro Person gekostet.} (1 pt.)
\item \all{Sie verlangt, dass die Firma ihr die Hälfte des Preises zurückzahlt. Außerdem möchte sie eine schriftliche Bestätigung, dass solche Probleme nicht mehr vorkommen.} (2 pts.)
\item \all{Am Ende benutzt sie die Formel „Mit freundlichen Grüßen“.} (1 pt.)
\item Deux phrases au subjonctif II : \all{„Könnten Sie mir außerdem schriftlich bestätigen, dass solche Probleme in Zukunft nicht mehr vorkommen?“} et \all{„Ich wäre Ihnen sehr dankbar für eine schnelle Antwort.“} (2 pts : 1 pt par phrase correctement recopiée.)
\end{enumerate}
\textbf{Vrai ou faux :}
\begin{enumerate}
\item[a)] \textbf{Faux.} Le texte dit : \all{„Der Bus kam zwei Stunden zu spät“} — le bus avait deux heures de retard, il n'était donc pas à l'heure. (2 pts : 1 pt pour la réponse, 1 pt pour la citation.)
\item[b)] \textbf{Faux.} Amina écrit : \all{„Ich möchte Sie bitten, mir die Hälfte des Preises zurückzuzahlen“} — elle demande le remboursement de la moitié du prix, pas de la totalité. (2 pts.)
\end{enumerate}
\textbf{Points de vigilance :} répondre par des phrases complètes avec le verbe en 2\up{e} position ; reprendre le vocabulaire du texte sans le recopier intégralement ; au vrai/faux, la justification par une citation exacte est obligatoire pour avoir tous les points.
\end{corrige}

\begin{corrige}[Exercice 10 — Type Bac : Thème et version]
\textbf{A. Thème (10 pts : 2 pts par phrase).}
\begin{enumerate}
\item \all{Ich habe am 10. April ein Handy bei Ihnen bestellt.} — \emph{Points de vigilance :} Perfekt avec \all{haben} + participe \all{bestellt} en fin de phrase ; la date se met avec \all{am} ; \all{bei Ihnen} = « chez vous ».
\item \all{Leider funktioniert das Gerät nicht.} — \emph{Vigilance :} après \all{leider} en tête de phrase, inversion sujet-verbe (\all{funktioniert das Gerät}) ; « l'appareil » = \all{das Gerät} (ne pas confondre avec \all{das Handy}, le téléphone portable).
\item \all{Ich möchte Sie bitten, mir das Geld zurückzuzahlen.} — \emph{Vigilance :} \all{jemanden bitten, etwas zu tun} ; « me rembourser » = \all{mir das Geld zurückzahlen} (datif \all{mir}). On évite \all{mich erstatten} : \all{erstatten} se construit avec une chose (\all{die Kosten erstatten}, « rembourser les frais »), jamais avec une personne à l'accusatif.
\item \all{Könnten Sie mir bitte ein neues Gerät schicken?} — \emph{Vigilance :} subjonctif II \all{könnten} en 1\up{re} position (question sans mot interrogatif), infinitif \all{schicken} en fin de phrase ; adjectif épithète à l'accusatif neutre : \all{ein neues Gerät}.
\item \all{Ich wäre Ihnen für eine schnelle Antwort sehr dankbar.} ou \all{Ich würde mich über eine schnelle Antwort sehr freuen.} — \emph{Vigilance :} « reconnaissant envers qqn » = \all{jemandem (Ihnen) dankbar}, datif obligatoire.
\end{enumerate}
\textbf{B. Version (10 pts).} Traduction proposée :

« \textbf{Objet : votre candidature pour un stage}

Monsieur Nkoulou,

Nous vous remercions de votre candidature et de votre CV. Nous aimerions avoir plus d'informations sur vos connaissances linguistiques. Pourriez-vous également nous dire quand vous pourriez commencer chez nous ? Nous serions très heureux de faire bientôt votre connaissance.

Nous vous prions d'agréer, Monsieur, l'expression de nos salutations distinguées.

Dr Sabine Weber, directrice du personnel »

\emph{Barème indicatif :} compréhension globale du sens (4 pts), exactitude du vocabulaire (3 pts), qualité du français (3 pts). \emph{Points de vigilance :} \all{die Bewerbung um ein Praktikum} = « la candidature pour un stage » ; \all{die Personalchefin} = « la directrice du personnel » (féminin de \all{der Personalchef}) ; \all{kennenlernen} = « faire la connaissance de » ; rendre les subjonctifs II par des conditionnels français (« nous aimerions », « pourriez-vous », « nous serions »).
\end{corrige}

\begin{corrige}[Exercice 11 — Type Bac : Expression écrite]
\textbf{Proposition de rédaction modèle (environ 90 mots) :}
\begin{textebox}[Modèle de lettre de réclamation]
Jean-Paul Etoundi \hfill Elektronik-Shop GmbH\\
Quartier Bastos \hfill Postfach 123\\
Yaoundé, Kamerun \hfill 10115 Berlin, Deutschland\\[4pt]
Yaoundé, den 20. Mai 2026\\[4pt]
\textbf{Betreff: Beschwerde über ein defektes Handy}\\[4pt]
Sehr geehrte Damen und Herren,\\[4pt]
am 5. Mai habe ich bei Ihnen ein Handy bestellt. Leider ist das Gerät gestern defekt angekommen: der Bildschirm ist kaputt, und das Handy lässt sich nicht einschalten. Ich bin sehr unzufrieden, denn ich habe 80\,000 FCFA bezahlt.\\[4pt]
Ich möchte Sie bitten, mir schnell ein neues Gerät zu schicken. Könnten Sie mir bitte auch schriftlich antworten? Ich wäre Ihnen sehr dankbar für eine schnelle Lösung.\\[4pt]
Mit freundlichen Grüßen\\
\emph{Jean-Paul Etoundi}
\end{textebox}
\textbf{Barème indicatif (20 pts) :}
\begin{itemize}
\item Respect de la forme de la lettre formelle (adresses, lieu/date, \all{Betreff}, appel, formule finale, signature) : 4 pts.
\item Exactitude grammaticale et place du verbe (verbe en 2\up{e} position, infinitif final, Perfekt correct) : 6 pts.
\item Richesse du vocabulaire thématique (\all{bestellen}, \all{defekt}, \all{der Bildschirm}, \all{unzufrieden}, \all{die Lösung}) : 6 pts.
\item Cohérence du contenu (problème clair, au moins deux détails, demande précise) : 4 pts.
\end{itemize}
\textbf{Points de vigilance :}
\begin{itemize}
\item Après \all{Sehr geehrte Damen und Herren,} (avec virgule), la première ligne du corps commence par une \textbf{minuscule} : \all{am 5. Mai habe ich...}.
\item Employer au moins un subjonctif II de politesse : \all{Könnten Sie ...?}, \all{Ich wäre Ihnen dankbar...}, \all{Ich möchte Sie bitten, ...}.
\item \all{sich einschalten lassen} : « l'appareil ne s'allume pas » se dit \all{das Handy lässt sich nicht einschalten} ou plus simplement \all{das Handy funktioniert nicht}.
\item Ne pas écrire \all{Liebe Grüße} à la fin : la lettre de réclamation exige \all{Mit freundlichen Grüßen}.
\item Majuscule à tous les noms : \all{das Handy}, \all{der Bildschirm}, \all{die Beschwerde}, \all{das Gerät}.
\end{itemize}
\end{corrige}

\newpage

\coursec{Fiche bilan}

\begin{popfigure}
\begin{tikzpicture}[mindmap, grow cyclic, every node/.style={concept, text=white, font=\small}, concept color=popblue, level 1/.append style={level distance=3.4cm, sibling angle=51}, level 2/.append style={level distance=2.2cm}]
\node[concept, font=\small\bfseries, align=center]{AL29\\Beschwerde\\und Lebenslauf}
  child[concept color=poporange] { node[concept] {Texte} 
    child { node[concept, font=\scriptsize] {Beschwerdebrief} }
    child { node[concept, font=\scriptsize] {Lebenslauf} }
  }
  child[concept color=popgreen] { node[concept] {Lexique}
    child { node[concept, font=\scriptsize] {die Bewerbung} }
    child { node[concept, font=\scriptsize] {die Kenntnisse} }
  }
  child[concept color=poppurple] { node[concept] {Konjunktiv II}
    child { node[concept, font=\scriptsize] {h\"atte, w\"are} }
    child { node[concept, font=\scriptsize] {k\"onnte, m\"ochte} }
  }
  child[concept color=poppink] { node[concept] {Briefaufbau}
    child { node[concept, font=\scriptsize] {Anrede} }
    child { node[concept, font=\scriptsize] {Gru\ss formel} }
  }
  child[concept color=popgold] { node[concept] {Methode}
    child { node[concept, font=\scriptsize] {lire, rep\'erer} }
    child { node[concept, font=\scriptsize] {r\'esumer, produire} }
  }
  child[concept color=popdark] { node[concept] {Expression}
    child { node[concept, font=\scriptsize] {CV + Brief} }
  };
\end{tikzpicture}
\legende{Carte mentale de la le\c{c}on AL29 : lettre de r\'eclamation et curriculum vitae.}
\end{popfigure}

\begin{bilan}
\textbf{1. Lexique cl\'e (article + pluriel \`a conna\^{\i}tre)}

\begin{center}
\begin{tabularx}{\linewidth}{|X|X|}\hline
\rowcolor{popblueL} \textbf{Deutsch} & \textbf{Fran\c{c}ais} \\ \hline
der Lebenslauf, die Lebensl\"aufe & le curriculum vitae (CV) \\ \hline
die Bewerbung, -en & la candidature \\ \hline
sich bewerben (bewarb, hat beworben) um + Akk. & poser sa candidature \`a \\ \hline
die Beschwerde, -n & la r\'eclamation, la plainte \\ \hline
sich beschweren \"uber + Akk. & se plaindre de \\ \hline
die Kenntnisse (Pl.) & les connaissances, comp\'etences \\ \hline
die Berufserfahrung, -en & l'exp\'erience professionnelle \\ \hline
die Ausbildung, -en & la formation \\ \hline
der Abschluss, die Abschl\"usse & le dipl\^ome \\ \hline
die Anzeige, -n / die Stelle, -n & l'annonce / le poste \\ \hline
die Reklamation, -en & la r\'eclamation (commerce) \\ \hline
der Umtausch / die R\"uckerstattung & l'\'echange / le remboursement \\ \hline
der Beleg, -e / die Quittung, -en & le re\c{c}u / la quittance \\ \hline
\end{tabularx}
\end{center}

\textbf{2. Formules de la lettre formelle}

\begin{center}
\begin{tabularx}{\linewidth}{|X|X|}\hline
\rowcolor{popblueL} \textbf{Deutsch} & \textbf{Fran\c{c}ais} \\ \hline
Sehr geehrte Damen und Herren, & Madame, Monsieur, \\ \hline
Sehr geehrter Herr M\"uller, / Sehr geehrte Frau Ngo, & Monsieur M\"uller, / Madame Ngo, \\ \hline
ich schreibe Ihnen, weil \dots & je vous \'ecris parce que \dots \\ \hline
ich m\"ochte mich \"uber \dots\ beschweren & je voudrais me plaindre de \dots \\ \hline
ich bewerbe mich um die Stelle als \dots & je pose ma candidature au poste de \dots \\ \hline
ich bitte Sie um \dots\ / ich w\"are Ihnen dankbar, wenn \dots & je vous prie de \dots / je vous serais reconnaissant(e) si \dots \\ \hline
Mit freundlichen Gr\"u\ss en & Veuillez agr\'eer mes salutations distingu\'ees \\ \hline
Anlage: Lebenslauf, Zeugnisse & Pi\`eces jointes : CV, dipl\^omes \\ \hline
\end{tabularx}
\end{center}

\textbf{3. Grammaire \`a conna\^{\i}tre par c\oe ur}

\begin{itemize}
\item \cle{Konjunktiv II} de politesse : \all{ich h\"atte} (j'aurais), \all{ich w\"are} (je serais), \all{ich k\"onnte} (je pourrais), \all{ich m\"ochte} (je voudrais), \all{ich sollte} (je devrais), \all{d\"urfte ich \dots?} (pourrais-je \dots ?).
\item Formation : radical du Pr\"ateritum + Umlaut (a, o, u $\rightarrow$ \"a, \"o, \"u) + terminaisons : \all{haben $\rightarrow$ hatte $\rightarrow$ h\"atte} ; \all{sein $\rightarrow$ war $\rightarrow$ w\"are} ; \all{k\"onnen $\rightarrow$ konnte $\rightarrow$ k\"onnte}.
\item \all{w\"urde} + infinitif pour les autres verbes : \all{Ich w\"urde Sie bitten, \dots} (je vous prierais de \dots).
\item Structure de la lettre : exp\'editeur en haut \`a gauche, destinataire, lieu et date (\all{Yaound\'e, den 15. M\"arz 2025}), objet (\all{Betreff}), formule d'appel + virgule, premi\`ere ligne en minuscule, corps en paragraphes, formule de politesse, signature.
\item Place du verbe : verbe conjugu\'e en 2\textsuperscript{e} position ; dans la subordonn\'ee (\all{weil, wenn, dass}), verbe conjugu\'e \`a la fin.
\end{itemize}

\textbf{4. M\'ethodes express}

\begin{itemize}
\item Comprendre : lire d'abord l'exp\'editeur, le destinataire et l'objet ; rep\'erer le probl\`eme, les faits (dates, lieux), la demande finale.
\item R\'ediger une r\'eclamation : 1) faits pr\'ecis, 2) probl\`eme, 3) demande concr\`ete (\all{Umtausch, R\"uckerstattung}), 4) ton poli avec Konjunktiv II.
\item R\'ediger un CV : \all{Pers\"onliche Daten} $\rightarrow$ \all{Schulbildung} $\rightarrow$ \all{Berufserfahrung} $\rightarrow$ \all{Sprachen und Kenntnisse} $\rightarrow$ \all{Hobbys} ; ordre antichronologique, sans phrase compl\`ete.
\end{itemize}
\end{bilan}

\begin{aretenir}[Checklist avant le Bac]
\begin{itemize}
\item[$\square$] Je conna\^{\i}s le lexique : \all{der Lebenslauf, die Bewerbung, die Beschwerde, die Kenntnisse, die Stelle} avec article et pluriel.
\item[$\square$] Je sais former le Konjunktiv II : \all{h\"atte, w\"are, k\"onnte, m\"ochte, sollte, w\"urde} + infinitif.
\item[$\square$] J'emploie le Konjunktiv II pour formuler une demande polie : \all{Ich w\"are Ihnen dankbar, wenn \dots}
\item[$\square$] Je ma\^{\i}trise les formules d'appel : \all{Sehr geehrte Damen und Herren,} et de cl\^oture : \all{Mit freundlichen Gr\"u\ss en}.
\item[$\square$] Apr\`es la formule d'appel et la virgule, je commence la phrase par une minuscule.
\item[$\square$] Je place le verbe conjugu\'e en 2\textsuperscript{e} position et \`a la fin dans les subordonn\'ees (\all{weil, wenn, dass}).
\item[$\square$] Je sais conjuguer \all{sich bewerben um + Akkusativ} et \all{sich beschweren \"uber + Akkusativ}.
\item[$\square$] Je mets la majuscule \`a tous les noms et j'utilise les guillemets allemands „\dots“ si je cite.
\item[$\square$] Je sais structurer un CV : donn\'ees personnelles, formation, exp\'erience, langues, loisirs.
\item[$\square$] Dans une r\'eclamation, je cite les faits (date, lieu, produit) et je formule une demande pr\'ecise.
\item[$\square$] Je relis ma copie : Umlaute, \ss{}, accords, orthographe de la date (\all{den 15. M\"arz}).
\end{itemize}
\end{aretenir}

\findecours

\end{document}
